% Options for packages loaded elsewhere
\PassOptionsToPackage{unicode}{hyperref}
\PassOptionsToPackage{hyphens}{url}
\documentclass[
  ukrainian,
]{article}
\usepackage{xcolor}
\usepackage[margin=2.5cm]{geometry}
\usepackage{amsmath,amssymb}
\setcounter{secnumdepth}{-\maxdimen} % remove section numbering
\usepackage{iftex}
\ifPDFTeX
  \usepackage[T1]{fontenc}
  \usepackage[utf8]{inputenc}
  \usepackage{textcomp} % provide euro and other symbols
\else % if luatex or xetex
  \usepackage{unicode-math} % this also loads fontspec
  \defaultfontfeatures{Scale=MatchLowercase}
  \defaultfontfeatures[\rmfamily]{Ligatures=TeX,Scale=1}
\fi
\usepackage{lmodern}
\ifPDFTeX\else
\setmainfont{DejaVu Serif}
\setsansfont{DejaVu Sans}
\setmonofont{DejaVu Sans Mono}
\setmathfont{STIX Two Math}
\fi
\ifPDFTeX\else
  % xetex/luatex font selection
\fi
% Use upquote if available, for straight quotes in verbatim environments
\IfFileExists{upquote.sty}{\usepackage{upquote}}{}
\IfFileExists{microtype.sty}{% use microtype if available
  \usepackage[]{microtype}
  \UseMicrotypeSet[protrusion]{basicmath} % disable protrusion for tt fonts
}{}
\makeatletter
\@ifundefined{KOMAClassName}{% if non-KOMA class
  \IfFileExists{parskip.sty}{%
    \usepackage{parskip}
  }{% else
    \setlength{\parindent}{0pt}
    \setlength{\parskip}{6pt plus 2pt minus 1pt}}
}{% if KOMA class
  \KOMAoptions{parskip=half}}
\makeatother
\usepackage{graphicx}
\graphicspath{{/Users/kalinicenkovitalijmikolajovic/.codex/worktrees/5c15/Script_for_translate_book/Output_Final/}}
\setkeys{Gin}{width=0.82\linewidth,height=0.82\textheight,keepaspectratio}
\makeatletter
\newsavebox\pandoc@box
\newcommand*\pandocbounded[1]{% scales image to fit in text height/width
  \sbox\pandoc@box{#1}%
  \Gscale@div\@tempa{\textheight}{\dimexpr\ht\pandoc@box+\dp\pandoc@box\relax}%
  \Gscale@div\@tempb{\linewidth}{\wd\pandoc@box}%
  \ifdim\@tempb\p@<\@tempa\p@\let\@tempa\@tempb\fi% select the smaller of both
  \ifdim\@tempa\p@<\p@\scalebox{\@tempa}{\usebox\pandoc@box}%
  \else\usebox{\pandoc@box}%
  \fi%
}
% Set default figure placement to htbp
\def\fps@figure{htbp}
\makeatother
\ifLuaTeX
\usepackage[bidi=basic,shorthands=off,provide=*]{babel}
\else
\usepackage[bidi=default,shorthands=off,provide=*]{babel}
\fi
\ifLuaTeX
  \usepackage{selnolig} % disable illegal ligatures
\fi
\setlength{\emergencystretch}{3em} % prevent overfull lines
\providecommand{\tightlist}{%
  \setlength{\itemsep}{0pt}\setlength{\parskip}{0pt}}
\usepackage{bookmark}
\IfFileExists{xurl.sty}{\usepackage{xurl}}{} % add URL line breaks if available
\urlstyle{same}
\hypersetup{
  pdflang={uk},
  hidelinks,
  pdfcreator={LaTeX via pandoc}}

\author{}
\date{}

\begin{document}

\pandocbounded{\includegraphics[keepaspectratio]{images/_page_355_Figure_5.jpeg}}

\begin{verbatim}
\begin{cases} u(0, t) = 0, \\ u(1, t) = 0, \end{cases}
\end{verbatim}

\pandocbounded{\includegraphics[keepaspectratio]{images/_page_355_Figure_7.jpeg}}

\begin{verbatim}
= \ln\left[\frac{z-1}{z+1}\right] + i \arg\left[\frac{z-1}{z+1}\right]
u + iv = (x + iy)^2 = x^2 - y^2 + 2ixy
v = \arg\left[\frac{z-1}{z+1}\right] = \arg\left[\frac{x+iy-1}{x+1+iy}\right] =
\end{verbatim}

\begin{verbatim}
\left\{
\end{verbatim}

\begin{verbatim}
= \arg\left[\frac{x^2+y^2-1+i2y}{(x+1)^2+y^2}\right] =
\end{verbatim}

Чи знаєте ви, як виглядає графік цієї функції з сегментом {[}0, 1{]}?
Подумай про це. Зверніть увагу, що синусоїдальне перетворення функції
\texttt{f(x)} визначається лише для додатних цілих чисел n.~Іншими
словами, скінченний синус і косинус перетворюють функції у чисельні
послідовності.

\pandocbounded{\includegraphics[keepaspectratio]{images/_page_185_Figure_3.jpeg}}

РИС. 25.1. Графики функции \texttt{f(x)} = 1 и ее преобразования.

\section{Свойства
преобразований}\label{ux441ux432ux43eux439ux441ux442ux432ux430-ux43fux440ux435ux43eux431ux440ux430ux437ux43eux432ux430ux43dux438ux439}

Прежде чем приступить к решению задач, мы должны получить некоторые
свойства этих преобразований.

\begin{verbatim}
Якщо you(x, t) — функція  \partial \theta yx  змінних, і ми виконуємо перетворення змінної x, то
\end{verbatim}

\begin{verbatim}
= \arg\left[\frac{2y}{x^2+y^2-1}\right]
\end{verbatim}

(Зверніть увагу, що перетворення було виконано на змінній x, і отримана
послідовність залежить лише від часу t.)

А как быть с производными? Ниже приводятся несколько полезных формул

\begin{verbatim}
\begin{cases} u(x, 0) = 1, \\ u_t(x, 0) = 0, \end{cases}
\varphi(x, y) = \frac{1}{\pi} \arctan\left[\frac{2y}{x^2 + y^2 - 1}\right].
w = \ln\left\{\frac{z-1}{z+1}\right\}
x^{2} + y^{2} = 1,
(x-1)^{2} + y^{2} = 9,
\begin{aligned} & (\text{УЧП}) & \phi_{xx} + \phi_{yy} = 0, & \text{внутри } D, \\ & (\text{ГУ}) & \begin{cases} \varphi(x, y) = 1 & \text{на } x^2 + y^2 = 1, \\ \varphi(x, y) = 2 & \text{на } (x - 1)^2 + y^2 = 9. \end{cases} \end{aligned}
\begin{split} \frac{d^2S_n\left(t\right)}{dt^2} &= -\left(\pi n\right)^2S_n\left(t\right) + 2n\pi\left[u\left(0,\ t\right) + \left(-1\right)^{n+1}u\left(1,\ t\right)\right] + D_n\left(t\right) = \\ &= -\left(n\pi\right)^2S_n\left(t\right) + D_n\left(t\right), \end{split}
(\text{УЧП}) \quad u_{tt} = u_{xx} + \sin(\pi x), \quad 0 < x < 1, \quad 0 < t < \infty,
\end{verbatim}

M A

\pandocbounded{\includegraphics[keepaspectratio]{images/_page_359_Figure_10.jpeg}}

THBLK1203X

\begin{verbatim}
w = 2t \left[ \frac{z-s}{z-t} \right],
u = \gamma [(x-s)(x-t)-\gamma^2],
v = \gamma [y(x-t)],
= \ln\left[\frac{z-1}{z+1}\right] + i \arg\left[\frac{z-1}{z+1}\right]
\gamma = 2t/[(x-t)^2 + y^2].
= \arg\left[\frac{x^2+y^2-1+i2y}{(x+1)^2+y^2}\right] =
\end{verbatim}

\begin{verbatim}
\begin{aligned} \varphi_{uu} + \varphi_{vv} &= 0, & \text{в кольче,} \\ (\Gamma Y) & & & & & & & \\ (\Gamma Y) & & & & & & \\ (\Gamma Y) & & & & & \\ & & & & & \\ (\Gamma Y) & & & & & \\ & & & & \\ & & & & \\ & & & & \\ & & & & \\ & & & \\ & & & \\ & & & \\ & & & \\ & & & \\ & & & \\ & & \\ & & \\ & & \\ & & \\ & & \\ & & \\ & & \\ & & \\ & & \\ & & \\ & & \\ & & \\ & & \\ & & \\ & & \\ & & \\ & & \\ & & \\ & & \\ & & \\ & & \\ & & \\ & & \\ & & \\ & & \\ & & \\ & & \\ & & \\ & & \\ & & \\ & & \\ & & \\ & & \\ & & \\ & & \\ & & \\ & & \\ & & \\ & & \\ & & \\ & & \\ & & \\ & & \\ & & \\ & & \\ & & \\ & & \\ & & \\ & & \\ & & \\ & & \\ & & \\ & & \\ & & \\ & & \\ & & \\ & & \\ & & \\ & & \\ & & \\ & & \\ & & \\ & & \\ & & \\ & & \\ & & \\ & & \\ & & \\ & & \\ & & \\ & & \\ & & \\ & & \\ & & \\ & & \\ & & \\ & & \\ & & \\ & & \\ & & \\ & & \\ & & \\ & & \\ & & \\ & & \\ & & \\ & & \\ & & \\ & & \\ & & \\ & & \\ & & \\ & & \\ & & \\ & & \\ & & \\ & & \\ & & \\ & & \\ & & \\ & & \\ & & \\ & & \\ & & \\ & & \\ & & \\ & & \\ & & \\ & & \\ & & \\ & & \\ & & \\ & & \\ & & \\ & & \\ & & \\ & & \\ & & \\ & & \\ & & \\ & & \\ & & \\ & & \\ & & \\ & & \\ & & \\ & & \\ & & \\ & & \\ & & \\ & & \\ & & \\ & & \\ & & \\ & & \\ & & \\ & & \\ & & \\ & & \\ & & \\ & & \\ & & \\ & & \\ & & \\ & & \\ & & \\ & & \\ & & \\ & & \\ & & \\ & & \\ & & \\ & & \\ & & \\ & & \\ & & \\ & & \\ & & \\ & & \\ & & \\ & & \\ & & \\ & & \\ & & \\ & & \\ & & \\ & & \\ & & \\ & & \\ & & \\ & & \\ & & \\ & & \\ & & \\ & & \\ & & \\ & & \\ & & \\ & & \\ & & \\ & & \\ & & \\ & & \\ & & \\ & & \\ & & \\ & & \\ & & \\ & & \\ & & \\ & & \\ & & \\ & & \\ & & \\ & & \\ & & \\ & & \\ & & \\ & & \\ & & \\ & & \\ & & \\ & & \\ & & \\ & & \\ & & \\ & & \\ & & \\ & & \\ & & \\ & & \\ & & \\ & & \\ & & \\ & & \\ & & \\ & & \\ & & \\ & & \\ & & \\ & & \\ & & \\ & & \\ & & \\ & & \\ & & \\ & & \\ & & \\ & & \\ & & \\ & & \\ & & \\ & & \\ & & \\ & & \\ & & \\ & & \\ & & \\ & & \\ & & \\ & & \\ & & \\ & & \\ & & \\ & & \\ & & \\ & & \\ & & \\ & & \\ & & \\ & & \\ & & \\ & & \\ & & \\ & & \\ & & \\ & & \\ & & \\ & & \\ & & \\ & & \\ & & \\ & & \\ & & \\ & & \\ & & \\ & & \\ & & \\ & & \\ & & \\ & & \\ & & \\ & & \\ & & \\ & & \\ & & \\ & & \\ & & \\ & & \\ & & \\ & & \\ & & \\ & & \\ & & \\ & & \\ & & \\ & & \\ & & \\ & & \\ & & \\ & & \\ & & \\ & & \\ & & \\ & & \\ & & \\ & & \\ & & \\ & & \\ & & \\ & & \\ & & \\ & &
\end{verbatim}

\begin{verbatim}
\begin{cases} u(x, 0) = 1, \\ u_t(x, 0) = 0, \end{cases}
\begin{cases} 4/n\pi, & n = 1, 3, \dots, \\ 0, & n = 2, 4, \dots, \end{cases}
x^{2} + y^{2} = 1,
\varphi(r) = a \ln r + b, \quad r = u^2 + v^2.
\end{verbatim}

Решение задачи выполним за три следующих шага:

ШАГ І (Определение подходящего преобразования.)

Поскольку переменная x меняется от 0 до 1, мы воспользуемся конечным
преобразованием, а именно синус-преобразованием; позже станет ясно
почему. Мы могли бы решить эту задачу с помощью преобразования Лапласа
по переменной t (по трудоемкости этот метод не отличается от конечного
синус-преобразования).

ШАГ 2 (Выполнение преобразования).

Для удобства будем использовать обозначение \texttt{S\_n}(t) = S{[}u{]}
. Применим синус-преобразование к исходному уравнению

\begin{verbatim}
\varphi(u, v) = 0.57 \ln(u^2 + v^2) + 1
\end{verbatim}

и воспользуемся тождеством для синус-преобразования, получим

\begin{verbatim}
\varphi(x, y) = 0.57 \ln(u^2 + v^2) + 1,
\begin{aligned} & (\text{УЧП}) & \phi_{xx} + \phi_{yy} = 0, & \text{внутри } D, \\ & (\text{ГУ}) & \begin{cases} \varphi(x, y) = 1 & \text{на } x^2 + y^2 = 1, \\ \varphi(x, y) = 2 & \text{на } (x - 1)^2 + y^2 = 9. \end{cases} \end{aligned}
w = \ln \left[ \frac{z-1}{z+1} \right].
z = re^{i\theta}
\begin{cases} 1, & n = 1, \\ 0, & n = 2, 3, \dots, \end{cases}
\end{verbatim}

где

\begin{verbatim}
\begin{aligned} (\text{Y} \forall \Pi) & \qquad \phi_{\text{vx}} + \phi_{yy} = 0, \quad 0 < x < \infty, \quad 0 < y < \infty, \\ \left\{ \begin{array}{l} \phi(x, \ 0) = \begin{cases} 1, & 0 < x < 1, \\ 0, & 1 \leqslant x < \infty, \end{cases} \\ \phi(0, \ y) = \begin{cases} 1, & 0 < y < 1, \\ 0, & 1 \leqslant y < \infty. \end{cases} \end{aligned}
\end{verbatim}

\begin{verbatim}
u = \gamma [(x-s)(x-t)-\gamma^2],
\gamma = 2t/[(x-t)^2 + y^2].
\end{verbatim}

Если теперь преобразовать начальные условия краевой задачи, то получим
начальные условия для обыкновенного дифференциального уравнения

\begin{verbatim}
\begin{cases} S_n(0) = \begin{cases} 4/n\pi, & n = 1, 3, \dots, \\ 0, & n = 2, 4, \dots, \end{cases}
\end{verbatim}

\begin{verbatim}
\varphi(r) = a \ln r + b, \quad r = u^2 + v^2.
\varphi(x, y) = 0.57 \ln(u^2 + v^2) + 1,
\end{verbatim}

Итак, решим теперь новое семейство задач Кощи:

(ОДУ)

\begin{verbatim}
w = \ln \left[ \frac{z-1}{z+1} \right].
\begin{cases} \frac{dS_n(0)}{dt} = 0. \end{cases}
\begin{cases} 1, & n = 1, \\ 0, & n = 2, 3, \dots, \end{cases}
\right.
\end{verbatim}

\pandocbounded{\includegraphics[keepaspectratio]{images/_page_362_Figure_1.jpeg}}

,

где

\pandocbounded{\includegraphics[keepaspectratio]{images/_page_362_Figure_2.jpeg}}

\pandocbounded{\includegraphics[keepaspectratio]{images/_page_362_Figure_5.jpeg}}

\pandocbounded{\includegraphics[keepaspectratio]{images/_page_362_Figure_7.jpeg}}

\begin{enumerate}
\def\labelenumi{\arabic{enumi}.}
\setcounter{enumi}{8}
\tightlist
\item
\end{enumerate}

\begin{verbatim}
H(x-a) =
\begin{cases} 0, x < a, \\ 1, x \geqslant a, \end{cases}
\begin{cases} u(0, t) = f(t) \\ u(L, t) = g(t) \end{cases}
\end{verbatim}

В результате получаем

\begin{verbatim}
H(x-a) =
\end{verbatim}

\texttt{f(x)} * \texttt{g(x)}

\begin{verbatim}
\begin{cases} 0, & x < a, \\ 1, & x \geqslant a, \end{cases}
f(x) = \mathcal{F}^{-1}[F] = \frac{1}{\sqrt{2\pi}} \int_{-\infty}^{+\infty} F(\omega) e^{-i\omega x} d\omega \qquad F(\omega) = \mathcal{F}[f] = \frac{1}{\sqrt{2\pi}} \int_{-\infty}^{+\infty} f(x) e^{-i\omega x} dx
\end{verbatim}

\begin{enumerate}
\def\labelenumi{\arabic{enumi}.}
\setcounter{enumi}{1}
\tightlist
\item
  Решите в общем виде задачу
\end{enumerate}

\begin{verbatim}
\begin{cases} 1, |x| < a \\ 0, |x| > a \end{cases}
\end{verbatim}

\section{ЗАМЕЧ
АНИЯ}\label{ux437ux430ux43cux435ux447-ux430ux43dux438ux44f}

\begin{enumerate}
\def\labelenumi{\arabic{enumi}.}
\tightlist
\item
  Для решения задачи методом конечного синус- или косинуспреобразования
  краевые условия при x=0 и x=L должны иметь вид
\end{enumerate}

\begin{verbatim}
\begin{cases} 1, |x| < 1 \\ 0, |x| > 1 \end{cases}
\end{verbatim}

\section{ПРИНЦИП СУПЕРПОЗИЦІЇ Є ОСНОВОЮ ТЕОРІЇ ЛІНІЙНИХ
СИСТЕМ}\label{ux43fux440ux438ux43dux446ux438ux43f-ux441ux443ux43fux435ux440ux43fux43eux437ux438ux446ux456ux457-ux454-ux43eux441ux43dux43eux432ux43eux44e-ux442ux435ux43eux440ux456ux457-ux43bux456ux43dux456ux439ux43dux438ux445-ux441ux438ux441ux442ux435ux43c}

МЕТА ЛЕКЦІЇ: Ознайомити читача з принципом суперпозиції та показати, як
його можна використати для розкладання початкової задачі на підзадачі,
потім розв'язання всіх підзадач одночасно і, підсумувавши отримані
результати, знаходження рішення початкової задачі. Ми також хочемо
показати, що два основні методи розв'язання лінійних рівнянь --- метод
розділення змінних і метод інтегрального перетворення --- використовують
принцип суперпозиції.

Припустимо, що інженеру потрібно знайти відгук лінійної системи на
вхідний ефект f.~Загальний підхід до розв'язання задачі такий:

\begin{verbatim}
- 1) розкласти f на елементарні ефекти  f = \sum f_k  ;
\end{verbatim}

\begin{itemize}
\tightlist
\item
  \begin{enumerate}
  \def\labelenumi{\arabic{enumi})}
  \setcounter{enumi}{1}
  \tightlist
  \item
    знайти реакцію системи \texttt{u\_k} на вплив \texttt{f\_k} ;
  \end{enumerate}
\end{itemize}

\begin{verbatim}
3) сложить отклики  u = \sum_{i=1}^{n} u_{i}  .
\end{verbatim}

Якщо система лінійна, то отримана сума є реакцією системи на вхідний
ефект f.~Це принцип суперпоезії (див. стор. 26.1).

РИС. 26.1. Основна ідея принципу суперпоенції; L --- диференціальний
оператор; f --- вхід; У---виходь.

Ця базова ідея може використовуватися для розв'язання складних змішаних
задач, розбиваючи їх на прості підзавдання, які розв'язуються за
допомогою

кожне з них, а потім підсумовування отриманих результатів (звісно,
диференціальне рівняння і граничні умови мають бути лінійними).

\section{Розкладання змішаної задачі на дві
простіші}\label{ux440ux43eux437ux43aux43bux430ux434ux430ux43dux43dux44f-ux437ux43cux456ux448ux430ux43dux43eux457-ux437ux430ux434ux430ux447ux456-ux43dux430-ux434ux432ux456-ux43fux440ux43eux441ux442ux456ux448ux456}

Предположим, нам необходимо решить линейную задачу (обозначим ее P)

(УЧП)

(Р) (ГУ)

\texttt{f(x)} * \texttt{g(x)}

\begin{verbatim}
(1+x^2)^{-1}
\begin{cases} u(0, t) = 0, \\ u(1, t) = 0, \\ u(x, 0) = \sin(2\pi x), \quad 0 \le x \le 1. \end{cases}
\frac{a}{x^2+a^2}
\frac{2ax}{(x^2+a^2)^2}
\end{verbatim}

\pandocbounded{\includegraphics[keepaspectratio]{images/_page_189_Picture_10.jpeg}}

\begin{enumerate}
\def\labelenumi{\arabic{enumi}.}
\setcounter{enumi}{17}
\tightlist
\item
\end{enumerate}

\begin{verbatim}
\begin{cases} u(0, t) = 0, \\ u(1, t) = 0, \\ u(1, t) = 0, \end{cases}
\begin{cases} 1-|x|, |x| < 1 \\ 0, |x| > 1 \end{cases}
\frac{2}{\pi} \frac{\sin a\omega}{\omega}
\frac{2}{\pi} \frac{\sin a\omega}{(\omega^2+a^2)^2}
\frac{\pi}{2} e^{-a+\omega}
\begin{cases} u_t = u_{xx} + f(x, t), \\ u(0, t) = 0, \\ u(1, t) = 0, \\ u(x, 0) = \varphi(x) \end{cases}
\frac{\pi}{2} e^{-a+\omega}
\frac{2ax}{(x^2+a^2)^2}
\cos(ax)
\frac{\cos(ax), |x| > \pi/2a}{0, |x| > \pi/2a}
\begin{cases} u_t = u_{xx} + f(x, t), \\ u(0, t) = 0, \\ u(1, t) = 0, \\ u(x, 0) = 0 \end{cases}
\cos(ax)
\frac{\pi}{2} [\delta(\omega+a) + \delta(\omega-a)]
\begin{cases} u_t = u_{xx}, \\ u(0, t) = 0, \\ u(1, t) = 0, \\ u(x, 0) = \varphi(x). \end{cases}
\end{verbatim}

f'\,'(0)

\begin{verbatim}
F(\omega) = \frac{?}{\pi} \int_{0}^{\infty} f(x) \sin(\omega x) dx
\end{verbatim}

f'\,'(0)

\begin{verbatim}
\frac{1}{\alpha} F\left(\frac{\omega}{\alpha}\right)
\frac{2\omega}{\pi (\alpha^2 + \omega^2)}
[2/\pi\omega]^{1/2}
\frac{2}{\pi\omega} [1-\cos(\omega a)]
e^{-1}a + \omega
\frac{1}{2}e^{-\alpha}\sin\omega
\begin{split} A_n(t) &= 2\int_0^1 u_t \sin(n\pi x) \, dx = \frac{d}{dt} \bigg[ 2\int_0^1 u(x, t) \sin(n\pi x) \, dx \bigg] = \frac{dU_n(t)}{dt} \,, \\ B_n(t) &= 2\int_0^1 u_{xx} \sin(n\pi x) \, dx = -(n\pi)^2 \, U_n(t) \,+ \\ &\quad + 2n\pi \, \big[ u(0, t) + (-1)^{n+1} \, u(1, t) \big], \end{split}
f(x) = \int_{0}^{\infty} F(\omega) \cos(\omega x) d\omega \qquad \qquad F(x) = \frac{2}{\pi} \int_{0}^{\infty} f(x) \cos(\omega x) dx
\begin{cases} u(0, t) = 0, \\ u(1, t) = 0 \end{cases}
S_n = \frac{2}{\pi} \int_0^{\pi} f(x) \sin(nx) dx.
y = \frac{\pi}{b - a} (x - a).
f(x) = \sum_{n=1}^{\infty} S_n \sin(nx).
f(x) = \sum_{n=1}^{\infty} S_n \sin(nx)
S_n = \frac{2}{\pi} \int_{0}^{\pi} f(x) \sin(nx) dx
0 \le x \le \pi
n = 1, 2, ...
\end{verbatim}

f'\,'(x)

\begin{verbatim}
\frac{\pi}{2} e^{ax}
C_n = \frac{2}{\pi} \int_0^{\pi} f(x) \cos(nx) dx.
0 \le x \le \pi
y = \frac{\pi}{b-a}(x-a).
f(x) = \frac{C_0}{2} + \sum_{n=1}^{\infty} C_n \cos(nx) \qquad C_n = \frac{2}{\pi} \int_{0}^{\pi} f(x) \cos nx \, dx
\end{verbatim}

f'\,'(x)

f'\,'(x)

\begin{verbatim}
\begin{array}{ll} u\left(x, \ 0\right) = \sin{(\pi x)}, \\ u_{t}\left(x, \ 0\right) = 0, \end{array}
\end{verbatim}

\begin{enumerate}
\def\labelenumi{\alph{enumi})}
\tightlist
\item
\end{enumerate}

\begin{verbatim}
5. \cos(mx), m=1, 2, \dots
\end{verbatim}

\texttt{\textbackslash{}mathbf\{\textbackslash{}Pi\}} для яких рівнянь
функція \texttt{u\_1\ +\ u\_2} також буде розв'язком? Які висновки ви
можете зробити зі своїх відповідей?

\begin{enumerate}
\def\labelenumi{\arabic{enumi}.}
\setcounter{enumi}{3}
\item
  Найдите четыре смещанные задачи, сумма решений которых дает решение
  следующей задачи:
\item
  Решите задачу Коши
\end{enumerate}

(OAY)

\begin{verbatim}
\begin{array}{ll} (\mathrm{Y} + \Pi) & u_{tt} = u_{xx} + \sin{(3\pi x)}, & 0 < x < 1, & 0 < t < \infty, \\ (\mathrm{\Gamma} \mathrm{Y}) & \left\{ \begin{array}{ll} u\left(0, \ t\right) = 0, \\ u\left(1, \ t\right) = 0, \\ \end{array}
\end{verbatim}

(HY) \texttt{U\_n}(0) = 0. Можете ли вы проверить полученное решение?

\begin{enumerate}
\def\labelenumi{\arabic{enumi}.}
\setcounter{enumi}{5}
\tightlist
\item
  Предположим, что функции \texttt{u\_1} и \texttt{u\_2} удовлетворяют
  линейным однородным граничным условиям
\end{enumerate}

\begin{verbatim}
\begin{cases}  u_x(1, t) + h_2 u(1, t) = 0.  Будет ли функция   u_i + u_i   удовлетворять этим условиям? , .
\end{verbatim}

Лекция 27

\section{УРАВНЕНИЯ ПЕРВОГО ПОРЯДКА (МЕТОД
ХАРАКТЕРИСТИК)}\label{ux443ux440ux430ux432ux43dux435ux43dux438ux44f-ux43fux435ux440ux432ux43eux433ux43e-ux43fux43eux440ux44fux434ux43aux430-ux43cux435ux442ux43eux434-ux445ux430ux440ux430ux43aux442ux435ux440ux438ux441ux442ux438ux43a}

ЦЕЛЬ ЛЕКЦИИ: Ввести понятие уравнения с частными производными первого
порядка (до сих пор мы рассматривали уравнения второго порядка) и
познакомиться с важным методом решения задач с начальными
условиями---методом характеристик. Задача, которой мы теперь займемся,
имеет вид

(УЧП)

Сразу же отметим, что впервые мы будем решать задачу с переменными
коэффициентами. Оказывается, если от координат (x, t) перейти к новым,
характеристическим координатам (s, \tau) , то наше УЧП превратится в
обыкновенное дифференциальное уравнение. Затем можно решить обыкновенное
дифференциальное уравнение, т. е. найти u(s, \tau) , а на последнем шаге
выразить s и t через x и t и получить \texttt{u(x,\ t)}.

Напомним читателю, что, когда мы решали уравнение дифф

Для каждой из подзадач выберите метод решения по своему BKYCY,

\begin{enumerate}
\def\labelenumi{\arabic{enumi}.}
\setcounter{enumi}{2}
\tightlist
\item
  Нехай \texttt{u\_1} і \texttt{u\_2} є розв'язками наступних рівнянь:
\end{enumerate}

константа \alpha\^{}2 играла роль коэффициента диффузии, а v---была
скоростью среды. Очевидно, что если \alpha=0 (т. е. диффузии нет), то
решение будет смещаться как целое вдоль оси x со скоростью v (у нас
осталась только конвекция). Другими словами, если начальное условие
\texttt{u(x,0)}=\varphi(x) , то соответствующее решение уравнения

;

\begin{verbatim}
b)  u_t = u_{xx} + e^t  ;
\end{verbatim}

\begin{enumerate}
\def\labelenumi{\Alph{enumi})}
\setcounter{enumi}{1}
\tightlist
\item
\end{enumerate}

будет иметь вид \texttt{u(x,\ t)} = \varphi(x-vt) .

Это дает нам основания думать, что решение уравнения первого порядка

;

\begin{verbatim}
F)  u_t = u_{xx} + u^2  .
\end{verbatim}

переносится вдоль потока со скоростью

Конечно, если a и b---константы, то решение представляет собой бегущую с
постоянной скоростью волну. Если же \texttt{a(x,\ t)} и
\texttt{b(x,\ t)}

\begin{verbatim}
\begin{aligned} u_t &= u_{xx} + f(x, t), & 0 < x < 1, & 0 < t < \infty, \\ & \begin{cases} u(0, t) = g_1(t), \\ u(1, t) = g_2(t), \end{cases} & 0 < t < \infty, \\ & u(x, 0) = \varphi(x), & 0 \le x \le 1. \end{aligned}
\end{verbatim}

\begin{verbatim}
РИС. 27.1. Начальное значение в точке x оказывает влияние на решение только вдоль характеристики в плоскости переменных (x,t). Параметр s изменяется вдоль характеристик от нуля до бесконечности, начиная с линии начального условия;   \tau   постоянна вдоль каждой характеристики;   \tau   изменяется вдоль линии начального условия.
\end{verbatim}

зависят от x и t, то скорость потока изменяется как вдоль потока, так и
во времени (читатель сможет увидеть, что начальная кривая сильно
искажается). Здесь так много аналогий с конвекцией!

Вернемся к нашей основной задаче:

(УЧП)

(НУ) \texttt{u(x,\ 0)} = \varphi(x), -\infty \textless{} x \textless{}
\infty.

Решение этого линейного уравнения первого порядка основывается на
следующем физическом факте: начальное условие в некоторой точке x
переносится в tx-плоскости вдоль линии, которая

называется характеристи-кой (см. рис. 27.1)

Этим наше уразнение отличается от других уравнений (таких, как уравнение
теплопроводности \texttt{u\_t\ =\ u\_\{xx\}\ )},для которых начальное
значение в точке x оказывает влияние на решение во всех точках
пространства и во все моменты времени. Если читатель помнит, то
начальное сме- Р И С

x

\texttt{.\ 27.2.\ Начальное\ возмущение\ \ \ \textquotesingle{}u\textbackslash{}left(x,0\textbackslash{}right)\textquotesingle{}\ \ \ в\ точке\ порождает\ две\ волны.\ Начальное\ возмущение\ струны\ согласно\ волновому\ уравнению\ распространяется\ вдоль\ двух\ характеристик.}

щение скрипичной струны в точке x оказывает влияние на решение вдоль
\partial syx линий в tx-плоскости (соответствующих двум бегущим волнам)
(см. рис. 27.2).

Естественно, напрашивается идея ввести две новые координаты s и
\tau   (вместо x и t) так, чтобы

\begin{itemize}
\tightlist
\item
  s изменялась вдоль характеристических кривых,
\item
  т изменялась вдоль начальной кривой (желательно вдоль линии t=0).
\end{itemize}

Рассмотрим сначала новую координату s. Если s выбрать так, что она
удовлетворяет приведенному выше условию, то уравнение

превратится в обыкновенное дифференциальное уравнение К о н

\_ \textbar, -\infty \textless{} x \textless{} \infty  , 0 \textless{} t
\textless{} \infty  ,

\begin{verbatim}
ечно, остается открытым вопрос, как найти эти характеристики. Ответ прост: мы выберем характеристики   \{[x(s), t(s)]: 0 < s < \infty\}   так, чтобы они удовлетворяли системе
\end{verbatim}

\begin{verbatim}
(HV)   u(x, 0) = \varphi(x), \quad -\infty < x < \infty.
\Delta u = u_{xx} + u_{yy},   -\infty < x < \infty  ,   0 < t < \infty
\Delta u = u_{rr} + \frac{1}{r}u_r + \frac{1}{r^2}u_{\theta\theta}
\Delta u = u_{xx} + u_{yy} + u_{zz}
\Delta u = u_{rr} + \frac{1}{r}u_r + \frac{1}{r^2}u_{\theta\theta} + u_{zz}.
-\infty < y < +\infty
\end{verbatim}

M A

\pandocbounded{\includegraphics[keepaspectratio]{images/_page_196_Figure_10.jpeg}}

THBLK1343X

\begin{verbatim}
\Delta u = u_{rr} + \frac{2}{r}u_r + \frac{1}{r^2}u_{\theta\theta} + \frac{\cot\theta}{r^2}
\end{verbatim}

M A

\pandocbounded{\includegraphics[keepaspectratio]{images/_page_197_Figure_5.jpeg}}

THBLK1345X

\begin{verbatim}
\Delta u = u_{rr} + \frac{2}{r}u_r + \frac{1}{r^2}u_{\theta\theta} + \frac{\cot\theta}{r^2}u_0 + \frac{1}{r^2\sin\theta}u_{\phi\phi}
\end{verbatim}

M A

\pandocbounded{\includegraphics[keepaspectratio]{images/_page_370_Figure_1.jpeg}}

THBLK1347X M A

\pandocbounded{\includegraphics[keepaspectratio]{images/_page_370_Picture_2.jpeg}}

THBLK1348X M A

\pandocbounded{\includegraphics[keepaspectratio]{images/_page_370_Figure_3.jpeg}}

THBLK1349X

\begin{verbatim}
\begin{cases} \frac{dx}{ds} = a(x, t), \\ \frac{dt}{ds} = b(x, t). \end{cases}
x(s) = s + c_1, \quad t(s) = s + c_2.
\frac{du}{ds} + 2u = 0, \quad 0 < s < \infty,
\begin{array}{llll} (\forall \Pi \Pi) & u_x + u_t + 2u = 0, & -\infty < x < \infty, & 0 < t < \infty \\ (\forall \Psi) & u(x, 0) = \sin x, & -\infty < x < \infty. \end{array}
-\infty & < x < +\infty & r \geqslant 0, \ 0 \leqslant \theta \leqslant 2\pi, \\
\frac{dx}{ds} = 1, \quad \frac{dt}{ds} = 1, \quad 0 < s < \infty.
-\infty & < y < +\infty & -\alpha \leqslant z \leqslant \alpha, \\
x(s) = s + c_1, \quad t(s) = s + c_2.
x(0) = \tau,
s=t
-\infty & < z < +\infty & x = r \cos \theta, \\
u(x, t) = \sin(x-t)e^{-xt}
x = s + \tau
(\forall \Pi) \quad a(x, t) u_x + b(x, t) u_t + c(x, t) u = 0, \quad -\infty < x < \infty, \\ 0 < t < \infty,
x-t=\tau
u(x, 0) = f(x), \quad -\infty < x < \infty.
y = r \sin \theta, \\
\frac{dx}{ds} = a(x, t), \qquad \frac{dt}{ds} = b(x, t).
\frac{du}{ds} + 2u = 0, \quad 0 < s < \infty,
x = x (s, \tau),
u(0) = \sin \tau.
z = z.
(\text{УЧП})
u\left(s,\,\tau\right)=\sin\tau e^{-2s}
\frac{dx}{ds} = x
\end{verbatim}

\begin{verbatim}
\end{array}
\end{verbatim}

\begin{verbatim}
x = s + \tau
x = \tau e^s,
s=t
t = s.
u(x, t) = \sin(x-t)e^{-xt}
\frac{du}{ds} + su = 0, \quad 0 < s < \infty,
\Delta u = 0
u(0) = F(\tau).
(\forall \Pi) \quad a(x, t) u_x + b(x, t) u_t + c(x, t) u = 0, \quad -\infty < x < \infty, \\ 0 < t < \infty,
u(s, \tau) = F(\tau) e^{-s} \ell^2
\Delta u + \lambda^2 u = 0
s=t
u(x, 0) = f(x), \quad -\infty < x < \infty.
u(x, t) = F(xe^{-t}) e^{-t^2/2}
\Delta u = k
\end{verbatim}

Уравнение одномерной диффузии Уравнение конвективной диффузии

Уравнение теплопроводности с поглощением Уравнение теплопроводности с
источником

\section{Діагональна
матриця}\label{ux434ux456ux430ux433ux43eux43dux430ux43bux44cux43dux430-ux43cux430ux442ux440ux438ux446ux44f}

\begin{verbatim}
Нехай A — квадратна матриця порядку n і нехай всі її власні значення  \lambda_1  ,  \lambda_2  ...  \lambda_n  різні. Побудуємо матрицю P так, щоб у неї були *стовпці* координати власних векторів (k-й стовпець містить власний вектор  X_k  , що відповідає власному значенню  \lambda_k  ), тобто
\end{verbatim}

\begin{verbatim}
X_1 =
\begin{bmatrix} X_1 & X_{\tilde{2}} & \dots & X_n \\ X_{1} & X_{2} & \dots & X_n \end{bmatrix}
\begin{bmatrix} 2 \\ 1 \end{bmatrix}
\begin{bmatrix} x_{11} & x_{1\tilde{2}} & \dots & x_{1n} \\ x_{21} & x_{22} & \dots & x_{2n} \\ \vdots & \vdots & \ddots & \vdots \\ x_{n1} & x_{n\tilde{3}} & \dots & x_{nn} \end{bmatrix}
U(x) = \int_{0}^{\infty} u(x, t) e^{-st} dt
P^{-1}AP =
\end{verbatim}

Тогда матрица Л, определяемая соотношением

где \texttt{P\^{}\{-1\}} --- обратная к матрице P, будет диагональной
матрицей вида

\begin{verbatim}
\begin{bmatrix} 2 & -2 \\ 1 & 1 \end{bmatrix}
\begin{bmatrix} \lambda_1 & & & \\ & \lambda_2 & 0 & \\ & & \ddots & \\ & & & \lambda_n \end{bmatrix}
^{-1}
\end{verbatim}

Например, у матрицы

\begin{verbatim}
\begin{bmatrix} 0 & 8 \\ 2 & 0 \end{bmatrix}
\begin{bmatrix} 2 & -2 \\ 1 & 1 \end{bmatrix}
= \\ =
in  X_2 =
\end{verbatim}

Функція Хевісайду,

\begin{verbatim}
H(a-x) = \begin{cases} 1, & x \leqslant a, \\ 0, & x > a, \end{cases}
\end{verbatim}

дзеркальною функцією Хевісайду.

ТАБЛИЦЯ A. Експоненціальне перетворення Фур'є

.

Значит, произведение

\begin{verbatim}
\begin{bmatrix} 1/4 & 1/2 \\ -1/4 & 1/2 \end{bmatrix}
\begin{bmatrix} 0 & 8 \\ 2 & 0 \end{bmatrix}
\begin{bmatrix} 2 & -2 \\ 1 & 1 \end{bmatrix}
\Lambda =
\end{verbatim}

равно диагональной матрице

\begin{verbatim}
\begin{bmatrix} 4 & 0 \\ 0 & -4 \end{bmatrix}
\end{verbatim}

(Читатель может проверить это самостоятельно.)

Тепер можна розв'язати просту систему з двох диференціальних рівнянь у
похідних у похідних (з відповідними початковими умовами).

\section{\texorpdfstring{Лінійне \texttt{u\_t\ +\ Au\_x\ =\ 0} системне
рішення.}{Лінійне u\_t + Au\_x = 0 системне рішення.}}\label{ux43bux456ux43dux456ux439ux43dux435-u_t-au_x-0-ux441ux438ux441ux442ux435ux43cux43dux435-ux440ux456ux448ux435ux43dux43dux44f.}

Розглянемо задачу Коші для системи, що містить два рівняння та дві
початкові умови:

(УЧП 1)

.

,\\
(29.3) (УЧП 2)
\texttt{\textbackslash{}frac\{\textbackslash{}partial\ u\_2\}\{\textbackslash{}partial\ t\}\ +\ 2\ \textbackslash{}frac\{\textbackslash{}partial\ u\_1\}\{\textbackslash{}partial\ x\}\ =\ 0}
,
\texttt{-\textbackslash{}infty\ \textless{}\ x\ \textless{}\ \textbackslash{}infty}
, \texttt{0\ \textless{}\ t\ \textless{}\ \textbackslash{}infty} ,\\
(HУ 1) \texttt{u\_1(x,\ 0)\ =\ f(x)} ,\\
\texttt{(HУ\ 2)\ \textquotesingle{}u\_2(x,\ 0)\ =\ g(x)\textquotesingle{}\ ,\ \textquotesingle{}-\textbackslash{}infty\ \textless{}\ x\ \textless{}\ \textbackslash{}infty\textquotesingle{}\ .}

Ця задача може відповідати визначенню \texttt{u\_1(x,\ t)} тиску та
густини \texttt{u\_2(x,\ t)} як функції просторової координати x і часу
t за відомими розподілами цих величин у початковий момент.

Запишем систему уравнений в матричной форме

\begin{verbatim}
\frac{\partial u_1}{\partial t} + 8 \frac{\partial u_2}{\partial x} = 0
\end{verbatim}

\begin{itemize}
\tightlist
\item
\end{itemize}

\begin{verbatim}
\begin{bmatrix} \frac{\partial u_1}{\partial x} \\ \frac{\partial u_2}{\partial x} \end{bmatrix}
\begin{bmatrix} \frac{\partial u_1}{\partial t} \\ \frac{\partial u_2}{\partial t} \end{bmatrix}
\end{verbatim}

\begin{itemize}
\tightlist
\item
\end{itemize}

\begin{verbatim}
\begin{bmatrix} 0 & 8 \\ 2 & 0 \end{bmatrix}
\begin{bmatrix} \frac{\partial u_1}{\partial x} \\ \frac{\partial u_2}{\partial x} \end{bmatrix}
=
\begin{bmatrix} 0 & 8 \\ 2 & 0 \end{bmatrix}
\begin{bmatrix} 0 \\ 0 \end{bmatrix}
\begin{bmatrix} \frac{\partial u_1}{\partial t} \\ \frac{\partial u_2}{\partial t} \end{bmatrix}
u_t + Au_s = 0
\begin{bmatrix} \frac{\partial u_1}{\partial x} \\ \frac{\partial u_2}{\partial x} \end{bmatrix}
A =
\end{verbatim}

.

\begin{verbatim}
\begin{bmatrix} 0 & 8 \\ 2 & 0 \end{bmatrix}
, \quad u_1 =
\end{verbatim}

с помощью преобразования

.

де P --- матриця, стовпці якої містять власні вектори матриці A (матриця
P, яку вже знають). Виявляється, що після такого перетворення буде
отримана дуже проста система для визначення v (два рівняння відносно
нового невідомого \texttt{v\_1} і \texttt{v\_2} виявляються
незалежними). Це означає, що \texttt{v\_1} і \texttt{v\_2} легко знайти.
Після \texttt{v\_1} і \texttt{v\_2} знайдено, згідно з формулою

\begin{verbatim}
\begin{bmatrix} \frac{\partial u_1}{\partial t} \\ \frac{\partial u_2}{\partial t} \end{bmatrix}
\end{verbatim}

находятся искомые функции и, и и2.

Спершу ж давайте розглянемо, як виглядає система визначення v.
Диференціюючи дві частини співвідношення u = Pv, отримуємо

(29 5)

\begin{verbatim}
, \quad u_x =
\begin{bmatrix} \frac{\partial u_1}{\partial x} \\ \frac{\partial u_2}{\partial x} \end{bmatrix}
, \quad 0 =
\begin{bmatrix} 0 \\ 0 \end{bmatrix}
\end{verbatim}

.

\begin{verbatim}
v =
\begin{bmatrix} v_1 \\ v_2 \end{bmatrix}
\end{verbatim}

.

\begin{verbatim}
u = Pv
u = Pv
\frac{\partial u}{\partial t} = P \frac{\partial v}{\partial t},
\frac{\partial u}{\partial x} = P \frac{\partial v}{\partial x}.
u_t + Au_x = 0
Pv_t + APv_x = 0.
\begin{bmatrix} v_1 \\ v_2 \end{bmatrix}
v_t + P^{-1}APv_x = 0,
(29.6) v_t + \Lambda v_x = 0.
\Lambda =
\begin{bmatrix} 4 & 0 \\ 0 & -4 \end{bmatrix}
\end{verbatim}

.

\begin{verbatim}
\frac{\frac{\partial v_1}{\partial t} + 4 \frac{\partial v_1}{\partial x} = 0, \\ \frac{\partial v_2}{\partial t} - 4 \frac{\partial v_2}{\partial x} = 0.
\end{verbatim}

.

\begin{verbatim}
v_1(x, t) = \varphi(x-4t),
u = Pv =
\begin{bmatrix} 2 & -2 \\ 1 & 1 \end{bmatrix}
\begin{bmatrix} v_1 \\ v_2 \end{bmatrix}
=
=
\begin{bmatrix} 2 & -2 \\ 1 & 1 \end{bmatrix}
\end{verbatim}

откуда функци і ф и ф легко находятся:

\begin{verbatim}
\begin{bmatrix} \varphi(x - 4t) \\ \psi(x + 4t) \end{bmatrix}
=
=
\end{verbatim}

(29.9)

\begin{verbatim}
\begin{bmatrix} 2\varphi(x - 4t) - 2\psi(x + 4t) \\ \varphi(x - 4t) + \psi(x + 4t) \end{bmatrix}
\end{verbatim}

.

\begin{verbatim}
u_1(x, t) = 2\varphi(x-4t) - 2\psi(x+4t), u_2(x, t) = \varphi(x-4t) + \psi(x+4t).
\phi(\xi) = \sin \xi, \ \psi(\xi) = \xi^2
u_1(x, t) = 2 \sin(x-4t) - 2(x+4t)^2,
u_1(x, 0) = f(x),
2\varphi(x) - 2\psi(x) = f(x),
\end{verbatim}

.

\begin{verbatim}
\varphi(x) + \psi(x) = g(x),
\varphi(x) = \frac{1}{4} \left[ f(x) + 2g(x) \right],
\begin{cases} \frac{\partial u_1}{\partial t} + \frac{\partial u_1}{\partial x} + \frac{\partial u_2}{\partial x} = 0, \\ \frac{\partial u_2}{\partial t} + 4 \frac{\partial u_1}{\partial x} + \frac{\partial u_2}{\partial x} = 0. \end{cases}
\psi(x) = \frac{1}{4} \left[ 2g(x) - f(x) \right].
u_{1}(x, t) = 2\varphi(x-4t) - 2\psi(x+4t) =
= \frac{1}{2} \left[ f(x-4t) + 2g(x-4t) \right] - \frac{1}{2} \left[ 2g(x+4t) - f(x+4t) \right],
u_{2}(x, t) = \varphi(x-4t) + \psi(x+4t) =
= \frac{1}{4} \left[ f(x-4t) + 2g(x-4t) \right] + \frac{1}{4} \left[ 2g(x+4t) - f(x+4t) \right].
u_t + A(x, t) u_x + B(x, t) u = 0
Au_t + Bu_x + Cu = 0,
A =
\begin{bmatrix} 1 & 1 \\ 4 & 1 \end{bmatrix}
\end{verbatim}

.

\begin{verbatim}
\begin{cases} \frac{\partial u_1}{\partial t} + \frac{\partial u_1}{\partial x} + \frac{\partial u_2}{\partial x} = 0, \\ \frac{\partial u_2}{\partial t} + 4 \frac{\partial u_1}{\partial x} + \frac{\partial u_2}{\partial x} = 0. \end{cases}
u_{tt} = c^2 \left( u_{rr} + \frac{1}{r} u_r + \frac{1}{r^2} u_{\theta\theta} \right)
u(r, \theta, t) = P(r) \Theta(\theta) T(t),
T'' + \lambda^2 c^2 T = 0
\begin{array}{llll} (\text{УЧП}) & u_{tt} = c^2 \left( u_{rr} + \frac{1}{r} u_r + \frac{1}{r^2} u_{\theta\theta} \right), & 0 < r < 1, \ 0 < t < \infty, \\ (\text{ГУ}) & u = 0 & \text{при } r = 1, & 0 < t < \infty, \\ (\text{НУ}) & \begin{cases} u = f(r, \ \theta) \\ u_t = g(r, \ \theta) \end{cases} & \text{при } t = 0. \end{array}
u(r, \theta, t) = U(r, \theta) T(t)
\Delta U + \lambda^2 U = 0
T'' + \lambda^2 c^2 T = 0
\Delta U = U_{rr} + \frac{1}{r} U_r + \frac{1}{r^2} U_{\theta\theta}.
\begin{array}{ll} (\mathsf{Y}\mathsf{H}\mathsf{\Pi}) & \Delta U + \lambda^2 U = 0, \\ (\mathsf{\Gamma}\mathsf{Y}) & U(1, \theta) = 0, \end{array}
u(1, \theta, t) = U(1, \theta) T(t) = 0, \quad 0 < t < \infty,
U(1, \theta) = 0.
\Delta U + \lambda^2 U = 0,
U(1, \theta) = 0.
\begin{cases} \Delta U + \mu u = 0 & \text{в круге,} \\ U = 0 & \text{на границе круга} \end{cases}
\begin{array}{ll} (\mathsf{Y}\mathsf{H}\mathsf{\Pi}) & \Delta U + \lambda^2 U = 0, \\ (\mathsf{\Gamma}\mathsf{Y}) & U(1, \theta) = 0, \end{array}
U(r, \theta) = R(r) \Theta(\theta)
r^2R'' + rR' + (\lambda^2r^2 - n^2)R = 0
r^2R'' + rR' + (\lambda^2r^2 - n^2) R = 0
r^2R'' + rR' + (\lambda^2r^2 - n^2)R = 0
R(r) = AJ_n(\lambda r) + BY_n(\lambda r).
R(r) = AJ_n(\lambda r).
\lambda = k_{nm}
U_{nm}(r, \theta) = J_n(k_{nm}r) [A \sin(n\theta) + B \cos(n\theta)].
f_{nm} = k_{nm}c/2\pi.
T_{nm}(t) = A \sin(k_{nm}ct) + B \cos(k_{nm}ct)
T'' + k_{nm}^2 c^2 T = 0.
sU(x) = -V \frac{dU}{dx}, \quad 0 < x < \infty,
U(0) = P/s.
H(\xi) =
\begin{cases} 1, & 0 \le \theta < \pi, \\ 0, & \pi \le \theta < 2\pi. \end{cases}
\begin{cases} 0, & \xi < 0, \\ 1, & \xi \geqslant 0. \end{cases}
\cos(3\varphi) = a_0 P_0(\cos\varphi) + a_1 P_1(\cos\varphi) + \dots
\Delta u = 0, \quad 0 < r < 1, \quad 0 \le \varphi \le \pi, \quad 0 \le \theta < 2\pi.
u(1, \varphi) =
\begin{cases} 1, & 0 \le \varphi \le \pi/2, \\ -1, & \pi/2 < \varphi \le \pi. \end{cases}
\Delta u = 0
u_{rr} + \frac{1}{r} u_r + \frac{1}{r^2} u_{\theta\theta} = f(r, \theta), \quad 0 \le r < 1, \quad 0 \le \theta < 2\pi,
\Delta u = -q
u(1, \theta) = 0, \quad 0 \le \theta \le 2\pi.
= -\int_{0}^{2\pi} u_{r}(r) r d\theta = -2\pi r u_{r}(r).
-2\pi r u_r(r) = q.
u(r) = \frac{-q}{2\pi} \ln r = \frac{q}{2\pi} \ln \frac{1}{r}.
u\left(r\right)=-\frac{q}{2\pi}\ln\frac{1}{r}.
(Y \cup \Pi)
u\left(r,\;\theta\right) = \int_{0}^{2\pi} \int_{0}^{1} G\left(r,\;\theta,\;\rho,\;\varphi\right) f\left(\rho,\;\varphi\right) \rho \; d\rho \; d\varphi.
\frac{1}{2\pi} \ln \left( \frac{1}{r} \right)
\frac{1}{2\pi} \ln \frac{1}{R}
u(r, \theta) = \frac{1}{2\pi} \ln 1/R - \frac{1}{2\pi} \ln 1/\overline{R}
-\frac{1}{2\pi}\ln\rho
G(r, \theta, \rho, \phi) = \frac{1}{2\pi} \ln 1/R - \frac{1}{2\pi} \ln 1/\overline{R} + \frac{1}{2\pi} \ln \rho,
R = V \frac{r^2 - 2\rho r \cos(\theta - \phi) + \rho^2}{R},
\overline{R} = V \frac{r^2 - 2\frac{r}{\rho}\cos(\theta - \phi) + 1/\rho^2}{r^2 - 2\frac{r}{\rho}\cos(\theta - \phi) + 1/\rho^2}
u(r, \theta) = \int_{0}^{2\pi} \int_{0}^{1} G(r, \theta, \rho, \varphi) f(\rho, \varphi) \rho d\rho d\varphi
(Y \cup \Pi)
u(r, \theta) = \frac{1}{2\pi} \int_{0}^{2\pi} \int_{0}^{1} \ln (\rho \overline{R}/R) f(\rho, \varphi) \rho d\rho d\varphi.
\begin{array}{llll} (\text{УЧ\Pi}) & \Delta u = 0, & 0 < r < 1, & 0 \leqslant \theta < 2\pi, \\ (\text{ГУ}) & u(1, \theta) = g(\theta), & 0 \leqslant \theta \leqslant 2\pi, \end{array}
u(r, \theta) = \int_{0}^{2\pi} \frac{\partial G}{\partial r} (r, \theta, 1, \varphi) g(\varphi) d\varphi.
u(r, \theta) = \frac{1}{2\pi} \int_{0}^{2\pi} \left[ \frac{1 - r^2}{1 - 2r\cos(\theta - \varphi) + r^2} \right] g(\varphi) d\varphi.
\begin{array}{llll} (\text{УЧП}) & \Delta u = f(r, \theta), & 0 < r < 1, & 0 \le \theta < 2\pi, \\ (\text{ГУ}) & u(1, \theta) = g(\theta), & 0 < \theta \le 2\pi, \end{array}
\Delta u = 1
u_p = (r, \theta) = Ar^2
\end{verbatim}

Формули, що виражають відстань між двома точками в полярних координатах.
Щоб знайти розв'язок початкової проблеми, потрібно накласти всі
імпульсні характеристики системи. Отже, переходимо до останнього кроку.

КРОК 3. Суперпозиція імпульсних характеристик Цей крок досить простий. З
відношення

\begin{verbatim}
f(x+h) = f(x) + f'(x)h + \frac{f''(x)}{2!}h^2 + \dots
\end{verbatim}

получаем

(36.4)

\begin{verbatim}
f(x+h) \cong f(x) + f'(x) h
\end{verbatim}

\begin{verbatim}
Формула (36.4) виражає розв'язок задачі Діріхле для рівняння Пуассона всередині кола через функцію Гріна цієї задачі. Якщо щільність заряду  f(r, \theta)  відома, то інтеграл можна обчислити, наприклад, чисельно.
\end{verbatim}

\section{ЗАУВАЖЕННЯ}\label{ux437ux430ux443ux432ux430ux436ux435ux43dux43dux44f}

\begin{enumerate}
\def\labelenumi{\arabic{enumi}.}
\tightlist
\item
  Задачу
\end{enumerate}

\begin{verbatim}
(37.1) f'(x) \cong \frac{f(x+h)-f(x)}{h}.
\end{verbatim}

також можна розв'язати за допомогою методу \emph{функції Гріна}. У цьому
випадку рішення фіксується у формі

\begin{verbatim}
f'(x) \cong \frac{f(x) - f(x - h)}{h}
\end{verbatim}

\begin{verbatim}
Останнє відношення можна надати більш обчислювально корисну форму, явно виставляючи вираз у  \partial G/\partial r  . Внаслідок цього ми отримуємо інтегральну формулу Пуассона, знайдену раніше в лекції 33.
\end{verbatim}

(36.5)

\begin{verbatim}
f'(x) \simeq \frac{1}{2h} [f(x+h) - f(x-h)].
\end{verbatim}

\begin{enumerate}
\def\labelenumi{\arabic{enumi}.}
\setcounter{enumi}{1}
\tightlist
\item
  Решение общей задачи Дирихле
\end{enumerate}

\begin{verbatim}
f''(x) \cong \frac{1}{h^2} [f(x+h) - 2f(x) + f(x-h)].
\end{verbatim}

находится как сумма решений (36.4) и (36.5).

\begin{enumerate}
\def\labelenumi{\arabic{enumi}.}
\setcounter{enumi}{2}
\item
  Метод функцій Гріна дозволяє отримати розв'язки багатьох задач у
  областях різних форм. Однак для кожної області і (точніше, для кожного
  оператора зліва від граничної умови) для кожного рівняння потрібно
  знайти власну функцію Гріна, і це не завжди легко зробити.
\item
  Щоб дійсно знайти розв'язок за формулою (36.4), у більшості випадків
  необхідно обчислити інтеграли на комп'ютері.
\end{enumerate}

\section{ЗАВДАННЯ}\label{ux437ux430ux432ux434ux430ux43dux43dux44f}

Найти потенциал точечного источника в трехмерном пространстве.

\begin{verbatim}
2. Знайти функцію Гріна  G(x, y, \xi, \eta)  для рівняння Лапласа у верхній півплощині y > 0. Іншими словами, знайти потенціал у (x, y) у верхній півплощині, якщо заряд розташований у  (\xi, \eta)  , а потенціал на межі y=0 дорівнює нулю (див. рисунок).
\end{verbatim}

\begin{verbatim}
u(x+h, y) = u(x, y) + u_x(x, y)h + u_{xx}(x, y)\frac{h^2}{2!} + \dots,
\end{verbatim}

\begin{verbatim}
ПРИМІТКА. Якщо від'ємний заряд розмістити в точці  \overline{Q} = (\xi, -\eta)  , очевидно, що негативний потенціал на прямій y = 0 буде нульовим. Те саме те, що функція Усмішки залежить від потенціалу цих двох зарядів.
\end{verbatim}

\begin{verbatim}
- 3. Використовуючи результати задачі 2, знайдемо розв'язок рівняння Пуассона  \Delta u = -k  у верхній півплощині з умовою нульової межі.
\end{verbatim}

\begin{itemize}
\tightlist
\item
  \begin{enumerate}
  \def\labelenumi{\arabic{enumi}.}
  \setcounter{enumi}{3}
  \tightlist
  \item
    Як би ви побудували функцію Зеленого для першого квадранта
    x\textgreater0,~y\textgreater0 ?
  \end{enumerate}
\item
  \begin{enumerate}
  \def\labelenumi{\arabic{enumi}.}
  \setcounter{enumi}{4}
  \tightlist
  \item
    Інший підхід до розв'язання рівняння Пуассона такий. Припустимо, ви
    хочете розв'язати проблему
  \end{enumerate}
\end{itemize}

(УЧП)

\begin{verbatim}
u(x-h, y) = u(x, y) - u_x(x, y)h + u_{xx}(x, y)\frac{h^2}{2!} - \dots,
,  0 < r < 1  ,  0 \le \theta < 2\pi  ,  (\Gamma Y)   u(1, \theta) = \sin \theta  ,  0 \le \theta \le 2\pi  .
\end{verbatim}

Сначала попытайтесь найти частное решение уравнения с помощью
подстановки

.

Потім підставляємо цей вираз у рівняння Пуассона і визначаємо константу
\texttt{A\_\{\textbackslash{}cdot\}} числа \texttt{u\_p} , що
знаходиться в частині розв'язку, необхідно пожартувати поза розв'язком з
огляду на \texttt{u=w+u\_p} . Яка межа розуму, яку можна задовольнити
функцією w? Знайте \texttt{w(r,\ \textbackslash{}theta)} . Як вирішити
проблему \texttt{u(r,\ \textbackslash{}theta)} ? Зачекай. Retelno
proanalizite vidpovid'. Зробіть перетин пеніса шкіри.

\section{Частина 5. ЧИСЕЛЬНІ ТА ПРИБЛИЗНІ
МЕТОДИ}\label{ux447ux430ux441ux442ux438ux43dux430-5.-ux447ux438ux441ux435ux43bux44cux43dux456-ux442ux430-ux43fux440ux438ux431ux43bux438ux437ux43dux456-ux43cux435ux442ux43eux434ux438}

Лекция 37

\section{ЧИСЕЛЬНІ РОЗВ'ЯЗКИ (ЕЛІПТИЧНІ
ЗАДАЧІ)}\label{ux447ux438ux441ux435ux43bux44cux43dux456-ux440ux43eux437ux432ux44fux437ux43aux438-ux435ux43bux456ux43fux442ux438ux447ux43dux456-ux437ux430ux434ux430ux447ux456}

МЕТА ЛЕКЦІЇ: Показати, як диференціальні рівняння в похідних похідних
можна звести до систем алгебраїчних рівнянь, якщо похідні з частковими
похідними замінити на наближення з скінченними різницями. Розв'язання
цієї системи алгебраїчних рівнянь ітеративним методом дозволяє отримати
наближене розв'язок

Надається інформація про пакет ELLPACK, розроблений для розв'язання
еліптичних приладів

задач общего вида на ЭВМ.

На сьогодні ми ознайомилися з кількома методами розв'язання лінійних
диференціальних рівнянь у похідних похідних. Однак більшість рівнянь, з
якими ми працювали, були дуже простими. Форми регіонів, у яких
розв'язувалися проблеми, також були простими. Однак багато завдань не
можна спростити настільки, щоб їх зводити до певного набору шаблонних
завдань. Такі задачі слід розв'язувати приблизно чисельними методами.
Прогрес у сфері високопродуктивних комп'ютерів призвів до створення
нових чисельних методів. Сьогодні такі нелінійні проблеми гідродинаміки,
теорії пружності та теорії потенціалу вже розв'язані, про які ніхто
навіть не думав десять років тому.

Під загальною назвою «чисельні методи» поєднуються кілька різних
підходів до розв'язання задач. Детальний опис усіх цих підходів можна
знайти у книзі {[}4{]} з рекомендованого списку джерел. У цій та
наступних двох лекціях ми покажемо, як розв'язувати еліптичні,
гіперболічні та параболічні рівняння методом скінченної різниці.

Спершу ознайомимося з поняттям скінченних різниць, а потім покажемо, як
їх використовувати для розв'язання задачі Діріхле в квадраті.

\section{Апроксимації скінченних
різниць}\label{ux430ux43fux440ux43eux43aux441ux438ux43cux430ux446ux456ux457-ux441ux43aux456ux43dux447ux435ux43dux43dux438ux445-ux440ux456ux437ux43dux438ux446ux44c}

Вспомним ряд Тейлора для функции \texttt{f(x)}:

Якщо ми розірвемо цю суперечку на другому терміні, отримаємо

\begin{verbatim}
\begin{split} u_x(x,\ y) & \cong \frac{u(x+h,\ y)-u(x,\ y)}{h}, \\ u_{xx}(x,\ y) & \cong \frac{1}{h^2} \big[ u(x+h,\ y)-2u(x,\ y)+u(x-h,y) \big], \\ u_u(x,\ y) & \cong \frac{u(x,\ y+k)-u(x,\ y)}{k}, \\ u_{uy}(x,\ y) & \cong \frac{1}{h^2} \big[ u(x,\ y+k)-2u(x,\ y)+u(x,\ y-k) \big]. \end{split}
\end{verbatim}

.

Откуда

Вираз праворуч називається похідною правої різниці. Вона апроксимує
першу похідну f'(x) у точці x.

В разложении Тейлора функции \texttt{f(\textbackslash{}mathbf\{r\})}
можно заменить h на -h и получить левую разностную производную

(37.2)

.

\begin{verbatim}
Buyutan  f(x - h) \cong f(x) - f'(x) h  \nif  f(x + h) \cong f( x) + f'(x) h  ,
\end{verbatim}

получаем центральную разностную производную

(37.3)

\begin{verbatim}
u_{xx} + u_{yy} = 0
\end{verbatim}

Якщо в ряду Тейлора залишиться ще один член, то центральну різницю
похідну для наближення f'\,'(x) можна отримати точно таким самим
способом

(37.4)

Тепер ми можемо розширити поняття наближення скінченних різниць до
\emph{часткових похідних}. Якщо виходити з розкладу Тейлора функції двох
змінних

\begin{verbatim}
\begin{split} u\left(x,\,y\right) &= u_{ij}\,,\\ u\left(x,\,y+k\right) &= u_{i+1,j}\,,\\ u\left(x,\,y-k\right) &= u_{i-1,j}\,,\\ u\left(x-h,\,y\right) &= u_{i,j-1}\,,\\ u\left(x+h,\,y\right) &= u_{i,j+1}\,,\\ u_{x}\left(x,\,y\right) &= \frac{1}{2h}\left(u_{i,j+1}-u_{i,j-1}\right)\,,\\ u_{y}\left(x,\,y\right) &= \frac{1}{2k}\left(u_{i+1,j}-u_{i-1,j}\right)\,,\\ u_{xx}\left(x,\,y\right) &= \frac{1}{h^{2}}\left(u_{i,j+1}-2u_{i,j}+u_{i,j-1}\right)\,,\\ u_{yy}\left(x,\,y\right) &= \frac{1}{h^{2}}\left(u_{i+1,j}-2u_{i,j}+u_{i-1,j}\right)\,. \end{split}
\begin{split} u_x(x,\ y) & \cong \frac{u(x+h,\ y)-u(x,\ y)}{h}, \\ u_{xx}(x,\ y) & \cong \frac{1}{h^2} \big[ u(x+h,\ y)-2u(x,\ y)+u(x-h,y) \big], \\ u_u(x,\ y) & \cong \frac{u(x,\ y+k)-u(x,\ y)}{k}, \\ u_{uy}(x,\ y) & \cong \frac{1}{h^2} \big[ u(x,\ y+k)-2u(x,\ y)+u(x,\ y-k) \big]. \end{split}
u_{xx} + u_{yy} = 0
\Delta u = \frac{1}{h^2} (u_{i,j+1} - 2u_{i,j} + u_{i,j-1}) + \frac{1}{k^2} (u_{i+1,j} - 2u_{i,j} + u_{i-1,j}) = 0.
u_{i,j} = \frac{1}{4} (u_{i+1,j} + u_{i-1,j} + u_{i,j+1} + u_{i,j-1}).
-4u_{22} + 0 + \sin(\pi/3) + u_{23} + u_{32} = 0,
-4u_{23} + u_{22} + \sin(2\pi/3) + 0 + u_{33} = 0,
\begin{split} u\left(x,\,y\right) &= u_{ij}\,,\\ u\left(x,\,y+k\right) &= u_{i+1,j}\,,\\ u\left(x,\,y-k\right) &= u_{i-1,j}\,,\\ u\left(x-h,\,y\right) &= u_{i,j-1}\,,\\ u\left(x+h,\,y\right) &= u_{i,j+1}\,,\\ u_{x}\left(x,\,y\right) &= \frac{1}{2h}\left(u_{i,j+1}-u_{i,j-1}\right)\,,\\ u_{y}\left(x,\,y\right) &= \frac{1}{2k}\left(u_{i+1,j}-u_{i-1,j}\right)\,,\\ u_{xx}\left(x,\,y\right) &= \frac{1}{h^{2}}\left(u_{i,j+1}-2u_{i,j}+u_{i,j-1}\right)\,,\\ u_{yy}\left(x,\,y\right) &= \frac{1}{h^{2}}\left(u_{i+1,j}-2u_{i,j}+u_{i-1,j}\right)\,. \end{split}
-4u_{32} + 0 + u_{22} + u_{33} + 0 = 0,
-4u_{33} + u_{32} + u_{23} + 0 + 0 = 0.
\begin{bmatrix} -4 & 1 & 1 \\ 1 & -4 & 0 & 1 \\ 1 & 0 & -4 & 1 \\ 0 & 1 & 1 & -4 \end{bmatrix}
\begin{bmatrix} u_{22} \\ u_{23} \\ u_{32} \\ u_{33} \end{bmatrix}
=
\begin{bmatrix} -\sin{(\pi/3)} \\ -\sin{(2\pi/3)} \\ 0 \\ 0 \end{bmatrix}
\end{verbatim}

Це потрібно вирішувати у зв'язку з \texttt{u\_\{22\}} ,
\texttt{u\_\{21\}} , \texttt{u\_\{32\}} та \texttt{u\_\{33\}} .
Розв'язок такої системи можна знайти за допомогою \emph{ітеративного
методу}. Метод Лібмана також є одним із ітеративних методів.

\begin{itemize}
\tightlist
\item
  \begin{enumerate}
  \def\labelenumi{\arabic{enumi}.}
  \setcounter{enumi}{1}
  \tightlist
  \item
    Якщо зменшити кроки сітки на h і k (кількість вузлів сітки
    збільшиться), то отримаємо систему рівнянь типу (37,7), але більших
    розмірів.
  \end{enumerate}
\item
  \begin{enumerate}
  \def\labelenumi{\arabic{enumi}.}
  \setcounter{enumi}{2}
  \tightlist
  \item
    Систему рівнянь (37.7) можна переписати у матричній формі
  \end{enumerate}
\end{itemize}

\begin{verbatim}
=
\begin{bmatrix} u_{22} \\ u_{23} \\ u_{32} \\ u_{33} \end{bmatrix}
\begin{bmatrix} -0.86 \\ -0.86 \\ 0 \\ 0 \end{bmatrix}
\end{verbatim}

.

\begin{verbatim}
u_{xx} + u_{yy} = f(x, y)
xu_{xx} + u_{yy} + 2u = \sin(x - y)
\sin xu_{xx} + u_{xy} + 3u = 0
\end{verbatim}

\pandocbounded{\includegraphics[keepaspectratio]{images/_page_285_Figure_8.jpeg}}

Зі збільшенням кількості рівнянь (наприклад, до 1000) матриця
коефіцієнтів стає розрідженою, тобто містить багато нулів. Системи з
розрідженими матрицями розв'язуються спеціальними методами. Зазвичай
використовуються ітеративні методи, такі як метод Якобі, метод
Гаусса-Зайделя або послідовний метод верхньої релаксації.

\begin{enumerate}
\def\labelenumi{\arabic{enumi}.}
\setcounter{enumi}{3}
\item
  При розв'язанні задачі Неймана ті \emph{похідні}, що входять до
  граничної умови, також слід замінити різницею апроксимаціями.
\item
  Таким чином ви зможете вирішувати проблеми
\end{enumerate}

\begin{enumerate}
\def\labelenumi{\alph{enumi})}
\tightlist
\item
\end{enumerate}

\texttt{(УЧ\textbackslash{}Pi)}

\begin{enumerate}
\def\labelenumi{\arabic{enumi})}
\setcounter{enumi}{5}
\tightlist
\item
\end{enumerate}

\begin{verbatim}
u_{xx} + u_{yy} + 2u = 0
\end{verbatim}

\begin{enumerate}
\def\labelenumi{\Alph{enumi})}
\setcounter{enumi}{1}
\tightlist
\item
\end{enumerate}

\begin{verbatim}
u_{xx} + u_{yy} = 0
\end{verbatim}

(гетерогенні рівняння з змінними коефіцієнтами), (рівняння зі змінними
коефіцієнтами).

\begin{enumerate}
\def\labelenumi{\arabic{enumi}.}
\setcounter{enumi}{5}
\tightlist
\item
  Якщо область, в якій розв'язується задача, має неправильну форму, її
  можна покрити сіткою і наближити розв'язок у точках, найближчих до
  межі, інтерполюючи граничні умови. Після цього задачу вирішують
  звичайним способом (див. рисунок 37.2).
\end{enumerate}

□ новые ГУ находятся интерполяцией ГУ на кривой

РИС. 37 2.

\begin{itemize}
\item
  \begin{enumerate}
  \def\labelenumi{\arabic{enumi}.}
  \setcounter{enumi}{6}
  \tightlist
  \item
    Деякі журнали публікують тексти програм для розв'язання
    диференціальних рівнянь у похідних на комп'ютерах. Нижче наведено
    короткий список цих журналів:
  \end{enumerate}

  \begin{itemize}
  \tightlist
  \item
    \begin{enumerate}
    \def\labelenumi{\alph{enumi})}
    \tightlist
    \item
      Транзакції ACM на математичному програмному забезпеченні,
    \end{enumerate}
  \item
    \begin{enumerate}
    \def\labelenumi{\arabic{enumi})}
    \setcounter{enumi}{5}
    \tightlist
    \item
      Комп'ютерний журнал,
    \end{enumerate}
  \item
    \begin{enumerate}
    \def\labelenumi{\Alph{enumi})}
    \setcounter{enumi}{1}
    \tightlist
    \item
      Numerische Mathematik,
    \end{enumerate}
  \item
    \begin{enumerate}
    \def\labelenumi{\alph{enumi})}
    \setcounter{enumi}{17}
    \tightlist
    \item
      БІТ
    \end{enumerate}
  \end{itemize}
\end{itemize}

Крім того, нещодавно було розроблено пакет застосунків ELLPACK для
розв'язання досить широкого класу задач на лінії прикордонних ліній. Цей
пакет дозволить вам розв'язувати різноманітні дво- та тривимірні задачі
у

різні координатні системи для довільних меж і спільних граничних умов.
Користувачу пропонується широкий спектр різних методів розв'язання 1).

\section{ЗАВДАННЯ}\label{ux437ux430ux432ux434ux430ux43dux43dux44f-1}

\begin{itemize}
\tightlist
\item
  1 Отримати наближення (39,4) для другої похідної f'\,'(x)
  \texttt{f\textquotesingle{}\textquotesingle{}(x)\ =\ \textbackslash{}\{f(x+h)\ 2f(x)\ +\ f(x-h)\textbackslash{}\}/h\^{}2}
  .
\item
  \begin{enumerate}
  \def\labelenumi{\arabic{enumi}.}
  \setcounter{enumi}{1}
  \tightlist
  \item
    Виконайте дві ітерації задачі Діріхле (37.5) за ітеративним методом
    Лібмана.
  \end{enumerate}
\item
  \begin{enumerate}
  \def\labelenumi{\arabic{enumi}.}
  \setcounter{enumi}{2}
  \tightlist
  \item
    До якої алгебранічної системи буде зведена задача Діріхле для
    рівнянь Пуассона в квадраті
  \end{enumerate}
\end{itemize}

\begin{verbatim}
\begin{cases} u = 0 & \text{квадрата,} \\ \frac{\partial u}{\partial x}(1, y) = 1, & 0 \leqslant y \leqslant 1, \end{cases}
\end{verbatim}

\texttt{u\_\{xx\}\ +\ u\_\{yy\}\ =\ f(x,y),\ \textbackslash{}quad\ 0\ \textless{}\ x\ \textless{}\ 1,\ \textbackslash{}quad\ 0\ \textless{}\ y\ \textless{}\ 1,}
\texttt{(\textbackslash{}Gamma\ Y)} \texttt{u(x,y)\ =\ g(x,y)} на
границе, если ее решать ме-

Скінченні відмінності? 4. Розв'язати задачу 3, якщо

(УЧП)

, 0 \textless{} x \textless{} 1 , 0 \textless{} y \textless{} 1 , (ГУ)
\texttt{u(x,\ y)} = \texttt{g(x,\ y)} на границе.

\begin{enumerate}
\def\labelenumi{\arabic{enumi}.}
\setcounter{enumi}{4}
\tightlist
\item
  Как бы вы решали задачу Неймана внутри квадрата
\end{enumerate}

(УЧП)

\begin{verbatim}
(\text{УЧП})
\end{verbatim}

, 0 \textless{} x \textless{} 1 , 0 \textless{} y \textless{} 1 , на
верхней, нижней и левой стороне квадрата,

\begin{verbatim}
u_t = u_{xx}
\end{verbatim}

методом конечных разностей?

\begin{enumerate}
\def\labelenumi{\arabic{enumi}.}
\setcounter{enumi}{5}
\tightlist
\item
  Постройте блок-схему алеоритма решения задачи Дирихле в квадрате
\end{enumerate}

\pandocbounded{\includegraphics[keepaspectratio]{images/_page_288_Figure_3.jpeg}}

\begin{verbatim}
u_{xx} + u_{yy} = f(x, y), \quad 0 < x < 1, \quad 0 < y < 1,   (\text{ГУ})   u(x, y) = g(x, y)  на границе,
\end{verbatim}

\begin{itemize}
\item
  з довільною кількістю вузлів сітки. Якщо ви знаєте мову програмування,
  напишіть програму для виконання цих обчислень.
\item
  \begin{enumerate}
  \def\labelenumi{\arabic{enumi})}
  \tightlist
  \item
    Російською мовою публікуються статті з чисельних методів розв'язання
    диференціальних рівнянь у численних журналах і збірках. Наприклад:
    «Журнал обчислювальної математики та математичної фізики»,
    видавництво «Наука», «Журнал інженерної фізики», видавництво
    Академії наук БРСР, збірки серії «Обчислювальні методи і
    програмування», Видавництво Московського державного університету.
    Багато публікацій, переважно статей, присвячені наборам прикладних
    проблем. Як приклад, можна вказати: Горбунов-Посадов М. М., Карпов
    В. Я, Корягін та ін. Пакет SAFRA. Програмне забезпечення для
    обчислювальних експериментів. У: Пакети прикладних програм.
    Обчислювальний експеримент. Москва, Наука, 1983. Примітка. Переклад.
  \end{enumerate}
\end{itemize}

\section{ЯВНІ СХЕМИ
РІЗНИЦІ}\label{ux44fux432ux43dux456-ux441ux445ux435ux43cux438-ux440ux456ux437ux43dux438ux446ux456}

МЕТА ЛЕКЦІЇ: Ввести концепцію явної різницької схеми та показати, як її
можна використовувати для розв'язання гіперболічних і параболічних
задач. Основна ідея полягає в тому, що після заміни рівняння типу

\begin{verbatim}
u_t = u_{xx}, \quad 0 < x < 1, \quad 0 < t < \infty,
\end{verbatim}

Її апроксимація скінченної різниці дає формули, які явно виражають
значення розв'язку за один момент часу через значення розв'язку в
попередній момент часу. Таким чином, змішану задачу для параболічного
або гіперболічного рівняння можна розв'язати, послідовно обчислюючи
розв'язки для всіх наступних точок часу.

Явна схема не позбавлена недоліків. Якщо необхідно підвищити точність
наближення похідних, то при зміні сегмента збільшується не лише обсяг
обчислень, а й похибки округлення.

Попередня лекція була присвячена розв'язанню еліптичних задач на межі
лінії (стаціонарних задач). У таких задачах необхідно знайти розв'язок
диференціального рівняння в похідних у заданій області простору, якщо
розв'язок або його похідна задано на межі області визначення Приблизне
розв'язання еліптичних задач зводиться до розв'язку системи алгебраїчних
рівнянь для значень функції у внутрішніх вузлах сітки. Іншими словами,
значення розв'язку у всіх внутрішніх вузлах визначаються одночасно.

\begin{verbatim}
У цій лекції ми розглянемо диференціальні методи розв'язання задач, залежних від часу. Основна ідея полягає в тому, що якщо розв'язок відомий у початковий момент часу, то можливо знайти розв'язок за допомогою схеми підрахунку  t=\Delta t,\ 2\Delta t,\ 3\Delta t,\ldots  Замінивши часткові похідні часу та просторову змінну на скінченнорізні похідні, можна отримати явні вирази для  u_{tj}  через значення функції u у попередні моменти часу. Такий процес називається явною схемою підрахунку з тикуванням.
\end{verbatim}

Щоб показати цей метод у дії, розглянемо типову задачу теплопровідності.

\section{Явная схема для уравнения
теплопроводности}\label{ux44fux432ux43dux430ux44f-ux441ux445ux435ux43cux430-ux434ux43bux44f-ux443ux440ux430ux432ux43dux435ux43dux438ux44f-ux442ux435ux43fux43bux43eux43fux440ux43eux432ux43eux434ux43dux43eux441ux442ux438}

Розглянемо проблему теплопровідності у стрижні, початкова температура
якого дорівнює нулю. Нехай температура лівого кінця фіксована, а на
правому кінці відбувається теплообмін із навколишнім середовищем, щоб
тепловий фланець був пропорційний P И

С. 38.1. Сетка для уравнения теплопроводности.

\texttt{різниця\ температур\ між\ кінцем\ стрижня\ і\ середовищем.\ Нехай\ температура\ середовища\ задається\ функцією\ \textquotesingle{}g\textbackslash{}left(t\textbackslash{}right)\textquotesingle{}\ .\ Іншими\ словами,\ ми\ розв\textquotesingle{}язуємо\ проблему}

задачу (УЧП)

\begin{verbatim}
\begin{cases} u(0, t) = 1, \\ u_x(1, t) = -[u(1, t) - g(t)], \\ u(x, 0) = 0, \quad 0 \le x \le 1. \end{cases}
\end{verbatim}

(38.1) (ГУ)

Щоб розв'язати цю задачу, ми побудуємо прямокутну сітку методом
скінченної різниці, вершини якої визначаються формулами (див. рис. 38.1)

, \texttt{j\ =\ 0,\ 1,\ 2,\ ...,\ n} , \texttt{y\_i\ =\ ik} ,
\texttt{i\ =\ 0,\ 1,\ 2,\ ...,\ m} .

Зверніть увагу, що значення \texttt{u\_\{ij\}} зліва і внизу сітки на
рис. 38.1 відомі з граничних і початкових умов, і наша вадача полягає у
знаходженні решти значень \texttt{u\_\{ij\}} . Щоб розв'язати цю задачу,
ми замінимо часткові похідні рівняння теплопровідності на їхні
наближення з скінченними різницями

\begin{verbatim}
x_j = jh
u_{t} = \frac{1}{k} [u(x, t+k) - u(x, t)] = \frac{1}{k} [u_{t+1,j} - u_{t,j}],
u_{xx} = \frac{1}{h^{2}} [u(x+h, t) - 2u(x, t) + u(x-h, t)] = \frac{1}{h^{2}} [u_{t,j+1} - 2_{t,j} + u_{t,j-1}].
u_{t+1,j} = u_{t,j} + \frac{k}{h^2} \left[ u_{t,j+1} - 2u_{t,j} + u_{t,j-1} \right].
a_x(1, t) = -[u(1, t) - g(t)].
\frac{1}{h}[u_{t,n} - u_{t,n-1}] = -[u_{t,n} - g_t].
u_{t,n} = \frac{u_{t,n-1} - hg_t}{1 + h}.
u_{2,j} = u_{1,j} + \frac{k}{h^2} [u_{1,j+1} - 2u_{1,j} + u_{1,j-1}], \quad j = 2, 3, \ldots, n-1.
u_{2,n} = \frac{u_{2,n-1} + hg_2}{1+h}.
u_{tt} = u_{xx}, \quad 0 < x < 1, \quad 0 < t < \infty,
\begin{cases} u(0, t) = g_1(t), \\ u(1, t) = g_2(t), \end{cases}
\quad 0 < t < \infty,
\begin{cases} u(x, 0) = \varphi(x), \\ u_t(x, 0) = \psi(x), \end{cases}
\begin{cases} u(0, t) = g_1(t), \\ u(1, t) = g_2(t), \end{cases}
\quad 0 \leqslant x \leqslant 1,
\end{verbatim}

(HY)

\begin{verbatim}
u_{tt} \simeq \frac{1}{h^2} [u(x, t+k) - 2u(x, t) + u(x, t-k)],
\quad 0 \leqslant x \leqslant 1,
u_{xx} \simeq \frac{1}{h^2} [u(x+h, t) - 2u(x, t) + u(x-h, t)].
u_t(x, 0) \approx \frac{1}{k} [u(x, k) - u(x, 0)] = \frac{1}{k} [u(x, k) - \varphi(x)].
u(x, t+k) = 2u(x, t) - u(x, t-k) + + (k/h)^2 [u(x+h, t) - 2u(x, t) + u(x-h, t)].
\end{verbatim}

а начальное условие--- по схеме

\begin{verbatim}
\frac{1}{k}\left[u\left(x,\,k\right)-\varphi\left(x\right)\right]=\psi\left(x\right),
\end{verbatim}

В результате для вычисления величины \texttt{u(x,\ t+k)} получаем
следующую явную схему:

(38.5)

\begin{verbatim}
u_t = u_{xx}
\end{verbatim}

З (38.5) видно, що для обчислення розв'язку на заданому часовому шарі
необхідно знати розв'язок на двох попередніх шарах. Отже, щоб почати
рахувати, потрібно використовувати P И С

. 38.4. Использование массивов при реализации явной схемы.

начальным условнем на скорость

\texttt{з\ якої\ ми\ виводимо\ \textquotesingle{}u(x,\ k)\ =\ \textbackslash{}varphi(x)\ +\ k\textbackslash{}psi(x)\textquotesingle{}\ ,\ тобто\ цінність\ розв\textquotesingle{}язку\ на\ \textquotesingle{}t\ =\ \textbackslash{}Delta\ t\textquotesingle{}\ .\ Розв\textquotesingle{}язок\ у\ наступні\ часи\ можна\ знайти\ за\ явною\ формулою\ (38.5).}

\section{ЗАВДАННЯ}\label{ux437ux430ux432ux434ux430ux43dux43dux44f-2}

\begin{enumerate}
\def\labelenumi{\arabic{enumi}.}
\tightlist
\item
  Постройте явную конечно-разностную схему для задачи
\end{enumerate}

(УЧП)

\begin{verbatim}
\begin{cases} u(0, t) = 0, \\ u(1, t) = 0, \\ u(x, 0) = \sin(\pi x), \quad 0 \le x \le 1. \end{cases}
\end{verbatim}

\texttt{,\ \textquotesingle{}0\ \textless{}\ x\ \textless{}\ 1\textquotesingle{}\ ,\ \textquotesingle{}0\ \textless{}\ t\ \textless{}\ \textbackslash{}infty\textquotesingle{}\ ,}
(ГУ)

\begin{verbatim}
(HУ)  u(x, 0) = \sin(\pi x), \quad 0 \le x \le 1.  Полагая  h = \Delta x = 0.1  , найдите решение при  t_1 = 0.005  ,  t_2 = 0.010  ,  t_3 = 0.015  . Постройте график полученного решения на сетке x = 0; 0.1; 0.2; 0.3; . .; 0.9: 1 при t = 0.015.
\end{verbatim}

\pandocbounded{\includegraphics[keepaspectratio]{images/_page_293_Figure_6.jpeg}}

РИС. 38.5. Залежність повної похибки від кроку сітки.

\begin{itemize}
\tightlist
\item
  \begin{enumerate}
  \def\labelenumi{\arabic{enumi}.}
  \setcounter{enumi}{1}
  \tightlist
  \item
    Розв'яжіть задачу 1 аналітично (ділячи змінні). Знайдіть значення
    розв'язку у вузлах сітки x=0; 0,1; 0,2; 0,3; \ldots; 0,9; 1 при
    t=0,015. Порівняйте ці результати з результатами чисельного
    розв'язку задачі 1. (Ви можете скористатися калькулятором або
    написати невелику комп'ютерну програму для табуляції аналітичного
    розв'язку.)
  \end{enumerate}
\item
  \begin{enumerate}
  \def\labelenumi{\arabic{enumi}.}
  \setcounter{enumi}{2}
  \tightlist
  \item
    Відповідно до примітки 3, побудуйте блок-схему для розв'язання
    гіперболічної задачі.
  \end{enumerate}
\item
  \begin{enumerate}
  \def\labelenumi{\arabic{enumi}.}
  \setcounter{enumi}{3}
  \tightlist
  \item
    Розв'яжіть задачу 1 з новою граничною умовою при x=1
    \texttt{u\_x}(1, t) = -{[}\texttt{u(1,\ t)} 1{]}. \# СХЕМИ НЕЯВНИХ
    РІЗНИЦЬ (СХЕМА КРАНКА--НІКОЛСОНА)
  \end{enumerate}
\end{itemize}

МЕТА ЛЕКЦІЇ: Показати, як можна розв'язувати задачі, що містять функції,
залежні від часу, використовуючи неявну різницю схеми. У цій схемі, як і
в попередній схемі, часткові похідні замінюються скінченно-різницею, але
\texttt{u\_\{l+1,\textbackslash{}\ J\}} більше не виражаються явно через
значення на попередніх шарах. Тепер, щоб визначити значення
\texttt{u\_\{l+1,\textbackslash{}\ J\}} , необхідно розв'язати систему
рівнянь. Іншими словами, щоб визначити розв'язок на кожному часовому
nage, необхідно розв'язати систему рівнянь.

Перевага неявних схем над явними полягає в тому, що в неявних схемах
танг сітки можна зробити достатньо великим без страху, що помилки
округлення «знищать» розв'язок.

Для розв'язання параболічної задачі ми використаємо відому схему
\emph{Кранк}-\emph{Ніколсона}.

Як уже згадувалося в попередній лекції, щоб явна схема працювала
надійно, часовий крок сітки має бути невеликим. Зокрема, якщо просте
завдання для рівняння теплопровідності

(39.1)

\texttt{розв\textquotesingle{}язується\ за\ явною\ схемою,\ тоді\ кроки\ сітки\ \textquotesingle{}\textbackslash{}Delta\ t\textquotesingle{}\ і\ \textquotesingle{}\textbackslash{}Delta\ x\textquotesingle{}\ повинні\ задовольняти\ умову}

\begin{verbatim}
\begin{array}{ccccccccccccccccccccccccccccccccccc
\end{verbatim}

\texttt{Інакше\ метод\ буде\ чисельно\ нестабільним\ (а\ помилки\ округлення\ зростуть\ нескінченно).\ (Питання\ чисельної\ стабільності\ розглядаються\ у\ {[}53{]}.)\ Іншими\ словами,\ якщо\ відстань\ сітки\ на\ осі\ x\ вибрана\ \textquotesingle{}\textbackslash{}Delta\ x=0,1\textquotesingle{}\ ,\ то\ часовий\ крок\ \textquotesingle{}\textbackslash{}Delta\ t\textquotesingle{}\ не\ може\ бути\ більшим\ за\ \textquotesingle{}0.5\textbackslash{}Delta\ x\^{}2=0.005\textquotesingle{}\ .\ Це\ означає,\ що\ ви\ можете\ пройти\ від\ t=0\ до\ t=1,\ зробивши\ 200\ кроків.}

Однак існують такі схеми (неявні схеми), які дозволяють виконувати
обчислення з набагато більшим кроком, тоді як

Правда, обсяг розрахунків, які виконуються на одному кроці, зростає.
Такі схеми можна використовувати для виконання обчислень із досить
великим кроком, але на кожному кроці потрібно розв'язати систему
алгебраїмних рівнянь. Щоб проілюструвати цей метод, розв'яжемо наступну
проблему теплопровідності.

\section{Неявная схема для уравнения
теплопроводности}\label{ux43dux435ux44fux432ux43dux430ux44f-ux441ux445ux435ux43cux430-ux434ux43bux44f-ux443ux440ux430ux432ux43dux435ux43dux438ux44f-ux442ux435ux43fux43bux43eux43fux440ux43eux432ux43eux434ux43dux43eux441ux442ux438}

Рассмотрим задачу

(39 2)

\begin{verbatim}
\frac{\Delta t}{(\Delta x)^2} \leqslant 0.5.
\end{verbatim}

Використаємо наступні апроксимації скінченних різниць для часткових
похідних \texttt{u\_t} та \texttt{u\_\{xx\}} :

\begin{verbatim}
\begin{array}{ccccccccccccccccccccccccccccccccccc
\end{verbatim}

де \texttt{\textbackslash{}lambda} взято з сегмента {[}0,1{]}. Зверніть
увагу, що \texttt{u\_\{xx\}} наближається зваженим середнім центральних
різницевих похідних у момент часу t і t+k.
\texttt{\textbackslash{}lambda=0.5} отримується звичайне середнє для цих
двох центральних похідних, і при \texttt{\textbackslash{}lambda=0.75}
одна з різницевих похідних бере вагу 0,75, а інша --- 0,25. У
\texttt{\textbackslash{}lambda=0} отримується звичайна явна схема, яка
була обговорена в попередній лекції.

Після заміни часткових похідних \texttt{u\_t} і \texttt{u\_\{xx\}} у
задачі (39.2) отримуємо різницю (див. рисунок 39.1).

Разностное уравнение

\begin{verbatim}
\begin{split} u_t(x,\ t) &= \frac{1}{k} \big[ u(x,\ t+k) - u(x,\ t) \big], \\ u_{xx}(x,\ t) &= \frac{\lambda}{h^2} \big[ u(x+h,\ t+k) - 2u(x,\ t+k) + u(x-h,\ t+k) \big] + \\ &\quad + \frac{(1-\lambda)}{h^2} \big[ u(x+h,\ t) - 2u(x,\ t) + u(x-h,\ t) \big], \end{split}
\end{verbatim}

(39.3)

\begin{verbatim}
\frac{1}{k} (u_{i+1, j} - u_{i, j}) = \frac{\lambda}{h^2} (u_{i+1, j+1} - 2u_{i+1, j} + u_{i+i, j-1}) + \frac{(1-\lambda)}{h^2} (u_{i, j+1} - 2u_{i, j} + u_{i, j-1}),
\end{verbatim}

298

Перенесіть усі невідомі значення u з верхнього часового шару (i+1) на
ліву частину рівняння (39.3) і отримайте (39.4)

(39.3)

де вводиться позначення \texttt{r\ =\ k/h\^{}2} . Зверніть увагу, що
якщо i фіксоване, а j змінюється від 2 до n-1, співвідношення (39.4)
визначають систему рівнянь n-2 з n-2 невідомим
\texttt{u\_\{t+1,2\},\ u\_\{t+1,3\},}

\begin{verbatim}
(\Gamma Y) \begin{cases} u_{i, 1} = 0, & i = 1, 2, ..., m, \\ u_{i, n} = 0, & i = 1, 2, ..., m, \\ u_{i, j} = 1, & j = 2, ..., n-1. \end{cases}
\end{verbatim}

РИС. 39.1. Сітка для схеми неявних різниць.

\begin{verbatim}
-\lambda r u_{i+1, j+1} + (1+2r\lambda) u_{i+1, j} - \lambda r u_{i+1, j-1} = = r(1-\lambda) u_{i, j+1} + [1-2r(1-\lambda)] u_{i, j} + r(1-\lambda) u_{i, j-1},
\end{verbatim}

РИС. 39.2. Шаблон неявної схеми
\texttt{u\_\{i+1,\ 4\},\ \textbackslash{}ldots,\ u\_\{i+1,\ n-1\}} , які
є розв'язком проблеми у внутрішній

\texttt{ренних\ узлах\ сетки\ на\ временном\ слое\ \textquotesingle{}t=(i+1)\textbackslash{}Delta\ t\textquotesingle{}\ .}

Фіг. 39.2 дає візуальне представлення структури кожного рівняння системи
(39.4). Тепер перейдемо до розв'язку системи (39.4).

\section{Алгоритм розв'язання задачі
(39.2)}\label{ux430ux43bux433ux43eux440ux438ux442ux43c-ux440ux43eux437ux432ux44fux437ux430ux43dux43dux44f-ux437ux430ux434ux430ux447ux456-39.2}

КРОК 1. Виберіть якусь цінність \texttt{\textbackslash{}lambda} (
\texttt{0\ \textbackslash{}le\ \textbackslash{}lambda\ \textbackslash{}le\ 1}
). Якщо \texttt{\textbackslash{}lambda\ =\ 0} , то рівняння (39.4)
перетворюються на явні формули з лекції 38.

\texttt{КРОК\ 2.\ Припустимо,\ наприклад,\ \textquotesingle{}h=\textbackslash{}Delta\ x=0,2\textquotesingle{}\ і\ \textquotesingle{}k=\textbackslash{}Delta\ t=0,08\textquotesingle{}\ (у\ цьому\ випадку\ \textquotesingle{}r=k/h\^{}2=2\textquotesingle{}\ ).\ У\ цьому\ випадку\ сітка\ містить\ 6\ вузлів\ уздовж\ осі\ (4\ внутрішні\ вузли)\ (див.\ рисунок\ 39.1).\ Візьмемо\ параметр\ ваги\ \textquotesingle{}\textbackslash{}lambda=0,5\textquotesingle{}\ (отримана\ схема\ називається\ схемою\ Кренка-Ніколсона).\ Відповідно\ до\ обчислювальної\ схеми\ (її\ зазвичай\ називають\ шаблоном),\ показаної\ на\ рис.\ На\ рисунку\ 39.2,\ рухаючись\ зліва\ направо\ (j=2,3,4,5)\ уздовж\ перших\ двох\ шарів\ (i=1),\ отримуємо\ наступні\ чотири\ рівняння:}

\pandocbounded{\includegraphics[keepaspectratio]{images/_page_296_Figure_5.jpeg}}

Перепишем их в матричной форме

(39.5)

\pandocbounded{\includegraphics[keepaspectratio]{images/_page_296_Figure_7.jpeg}}

\begin{verbatim}
\begin{array}{l} -u_{21}+3u_{22}-u_{23}=u_{11}-u_{12}+u_{13}=1,\\ -u_{22}+3u_{23}-u_{24}=u_{12}-u_{13}+u_{14}=1,\\ -u_{23}+3u_{24}-u_{25}=u_{13}-u_{14}+u_{15}=1,\\ -u_{24}+3u_{25}-u_{26}=u_{14}-u_{15}+u_{16}=1. \end{array}
\begin{bmatrix} u_{2\bar{z}} \\ u_{\bar{z}a} \\ u_{24} \\ u_{\bar{z}x} \end{bmatrix}
\begin{bmatrix} 3 & -1 & 0 & 0 \\ -1 & 3 & -1 & 0 \\ 0 & -1 & 3 & -1 \\ 0 & 0 & -1 & 3 \end{bmatrix}
\end{verbatim}

Матриця цієї системи називається тридіагональною. Щоб розв'язати систему
з трьома діагоналями

\begin{verbatim}
\begin{bmatrix} u_{2\bar{z}} \\ u_{\bar{z}a} \\ u_{24} \\ u_{\bar{z}x} \end{bmatrix}
\begin{bmatrix} x_1 \\ x_{\bar{z}} \\ x_{\bar{s}} \\ \vdots \\ x_n \end{bmatrix}
=
\begin{bmatrix} 1 \\ 1 \\ 1 \\ 1 \end{bmatrix}
\begin{bmatrix} d_t \\ d_{\bar{z}} \\ d_{\bar{s}} \\ \vdots \\ \vdots \\ d_n \end{bmatrix}
\end{verbatim}

.

преобразуем ее в эквивалентную систему вида

\begin{verbatim}
\begin{bmatrix} b_1 & c_1 & 0 & 0 & \dots & & & 0 \\ a_1 & b_2 & c_2 & 0 & \dots & & & 0 \\ 0 & a_{\bar{z}} & b_3 & c_3 & \dots & & & 0 \\ & & & & & & & & \\ \vdots & & & & & & & & \\ 0 & 0 & 0 & \dots & & & & a_{n-1} & b_n \end{bmatrix}
\begin{bmatrix} x_1 \\ x_2 \\ x_3 \\ \vdots \\ x_n \end{bmatrix}
\begin{bmatrix} x_1 \\ x_{\bar{z}} \\ x_{\bar{s}} \\ \vdots \\ x_n \end{bmatrix}
=
\begin{bmatrix} d_1^* \\ d_2^* \\ d_3^* \\ \vdots \\ d_n^* \end{bmatrix}
\begin{bmatrix} d_t \\ d_{\bar{z}} \\ d_{\bar{s}} \\ \vdots \\ \vdots \\ d_n \end{bmatrix}
\end{verbatim}

где

,

\texttt{,\ \textquotesingle{}c\_\{j+1\}\^{}*\ =\ \textbackslash{}frac\{c\_\{j+1\}\}\{b\_\{j+1\}\ -\ a\_j\ c\_j\^{}*\}\textquotesingle{}\ ,\ \textquotesingle{}j\ =\ 1,\ 2,\ \textbackslash{}ldots,\ n-2\textquotesingle{}\ ,\ и\ \textquotesingle{}d\_1\^{}*\ =\ d\_1/b\_1\textquotesingle{}\ ,\ \textquotesingle{}d\_\{j+1\}\^{}*\ =\ \textbackslash{}frac\{d\_\{j+1\}\ -\ a\_j\ d\_j\^{}*\}\{b\_\{j+1\}\ -\ a\_j\ c\_j\^{}*\}\textquotesingle{}\ ,\ \textquotesingle{}j\ =\ 1,\ 2,\ \textbackslash{}ldots,\ n-1\textquotesingle{}\ .}

У цьому перетворенні немає нічого дивного, оскільки друга система
повністю еквівалентна першої. Матриця нової системи організована так, що
цю систему дуже легко розв'язати. Розв'язуючи рівняння послідовно знизу
вгору, отримуємо

\begin{verbatim}
Застосовуючи цей метод до системи чотирьох рівнянь (39,5), отримуємо розв'язок  u_{22}=0,60  ;  u_{23}=0,80  ;  u_{24}=0,80  ;  u_{23}=0,60  . Ми отримали наближені значення розв'язку у внутрішніх точках сітки на  t=\Delta t  . Тепер ми можемо зробити наступний крок у часі, але для цього нам доведеться розв'язати нову систему рівнянь.
\end{verbatim}

У неявній схемі кількість обчислень на кожному кроці більша, ніж на
явному, але хорошу точність можна отримати навіть при значно більшому
кроці.

\section{ЗАВДАННЯ}\label{ux437ux430ux432ux434ux430ux43dux43dux44f-3}

\begin{itemize}
\tightlist
\item
  \begin{enumerate}
  \def\labelenumi{\arabic{enumi}.}
  \tightlist
  \item
    Из дифференциального уравнения (39.3) получите уравнение (39.4).
  \end{enumerate}
\item
  \begin{enumerate}
  \def\labelenumi{\arabic{enumi}.}
  \setcounter{enumi}{1}
  \tightlist
  \item
    Постройте неявную конечно-разностную схему для задачи
  \end{enumerate}
\end{itemize}

\begin{verbatim}
\begin{bmatrix} 1 & c_1^* & 0 & 0 & \dots & 0 \\ 0 & 1 & c_2^* & 0 & \dots & 0 \\ 0 & 0 & 1 & c_2^* & \dots & 0 \\ \vdots & \vdots & & & \vdots & \vdots \\ 0 & 0 & \dots & & 1 \end{bmatrix}
\begin{bmatrix} x_1 \\ x_2 \\ x_3 \\ \vdots \\ x_n \end{bmatrix}
=
\end{verbatim}

\begin{verbatim}
- 4. Как выглядит шаблон для уравнения (39.4), когда  \lambda = 1  ?
\end{verbatim}

\begin{itemize}
\tightlist
\item
  \begin{enumerate}
  \def\labelenumi{\arabic{enumi}.}
  \setcounter{enumi}{4}
  \tightlist
  \item
    Постройте блок-схему розв'язки рівняння теплопровідності (39.2).
    Якщо есть возможность, то напишите программу для ЭВМ Було бы очень
    полезно решить эту задачу численно при различных значениях λ с
    начальным условием и (x, 0) = sin (πx). Сравните полученные при
    различных λ численные результаты с аналитическим решением, яке в
    данном случае имеет вид
  \end{enumerate}
\end{itemize}

\begin{verbatim}
u_{\xi\eta} = \frac{\eta u_{\xi} - \xi u_{\eta}}{2(\xi^2 - \eta^2)}.
\end{verbatim}

6 Воспользовавшись формулами, приведенными в тексте, решите систему
алгебраических уравнений (39.5).

Лекція 40

\section{СРАВНЕНИЕ АНАЛИТИЧЕСКИХ РЕШЕНИЙ С ЧИСЛЕННЫМИ МЕТА ЛЕКЦІЇ:
Сравнить достоинства и недостатки аналитических и численных методов
розв'язки уравнений з частинними похідними. Приведена важная
математическая задача об идентификации физических величин
(параметрическая идентификация). В связи с цієї задачей рассматривается
один важный пример из
биологии.}\label{ux441ux440ux430ux432ux43dux435ux43dux438ux435-ux430ux43dux430ux43bux438ux442ux438ux447ux435ux441ux43aux438ux445-ux440ux435ux448ux435ux43dux438ux439-ux441-ux447ux438ux441ux43bux435ux43dux43dux44bux43cux438-ux43cux435ux442ux430-ux43bux435ux43aux446ux456ux457-ux441ux440ux430ux432ux43dux438ux442ux44c-ux434ux43eux441ux442ux43eux438ux43dux441ux442ux432ux430-ux438-ux43dux435ux434ux43eux441ux442ux430ux442ux43aux438-ux430ux43dux430ux43bux438ux442ux438ux447ux435ux441ux43aux438ux445-ux438-ux447ux438ux441ux43bux435ux43dux43dux44bux445-ux43cux435ux442ux43eux434ux43eux432-ux440ux43eux437ux432ux44fux437ux43aux438-ux443ux440ux430ux432ux43dux435ux43dux438ux439-ux437-ux447ux430ux441ux442ux438ux43dux43dux438ux43cux438-ux43fux43eux445ux456ux434ux43dux438ux43cux438.-ux43fux440ux438ux432ux435ux434ux435ux43dux430-ux432ux430ux436ux43dux430ux44f-ux43cux430ux442ux435ux43cux430ux442ux438ux447ux435ux441ux43aux430ux44f-ux437ux430ux434ux430ux447ux430-ux43eux431-ux438ux434ux435ux43dux442ux438ux444ux438ux43aux430ux446ux438ux438-ux444ux438ux437ux438ux447ux435ux441ux43aux438ux445-ux432ux435ux43bux438ux447ux438ux43d-ux43fux430ux440ux430ux43cux435ux442ux440ux438ux447ux435ux441ux43aux430ux44f-ux438ux434ux435ux43dux442ux438ux444ux438ux43aux430ux446ux438ux44f.-ux432-ux441ux432ux44fux437ux438-ux441-ux446ux456ux454ux457-ux437ux430ux434ux430ux447ux435ux439-ux440ux430ux441ux441ux43cux430ux442ux440ux438ux432ux430ux435ux442ux441ux44f-ux43eux434ux438ux43d-ux432ux430ux436ux43dux44bux439-ux43fux440ux438ux43cux435ux440-ux438ux437-ux431ux438ux43eux43bux43eux433ux438ux438.}

Вероятно, пришло время обсудить достоинства и недостатки аналитических и
численных решений уравнений з частинними похідними. Сначала выясним, що
мы подразумеваем под этими двумя типами решений.

\section{Аналитические
розв'язки}\label{ux430ux43dux430ux43bux438ux442ux438ux447ux435ux441ux43aux438ux435-ux440ux43eux437ux432ux44fux437ux43aux438}

Под аналитическими решениями мы подразумеваем такие розв'язки, в которых
неизвестная функція и выражена через независимые змінні и параметры
системи в виде формул, бесконечных рядов и интегралов.

\section{Численные
розв'язки}\label{ux447ux438ux441ux43bux435ux43dux43dux44bux435-ux440ux43eux437ux432ux44fux437ux43aux438}

Під чисельним ми зрозуміємо чисельні розв'язки, отримані чисельно після
приблизної заміни початкового рівняння іншим, простішим рівнянням.
Наприклад, у методі скінченної різниці похідні виражаються через
скінченні різниці, а розв'язок диференціального рівняння в похідних у
похідних апроксимується розв'язком різницького рівняння. Результатом
такої процедури зазвичай є таблиця значень розв'язків, а для деяких
значень пояснювальних змінних.

Таблина 40.1

\textbar{} × \textbar{} \textbar{} \textbar{} \textbar{} \textbar{}
\textbar{} \textbar{} \textbar{} \textbar{} \textbar{} \textbar{}

\textbar---------------------------\textbar-------------\textbar------------------\textbar-----------\textbar----------\textbar----------\textbar-----\textbar-----\textbar-----\textbar-----------\textbar------------------\textbar------------------\textbar{}

\textbar{} 0 \textbar{} 0,1 \textbar{} 0,2 \textbar{} 0,3 \textbar{} 0,4
\textbar{} 0,5 \textbar{} 0,6 \textbar{} 0,7 \textbar{} 0,8 \textbar{}
0,9 \textbar{} 1 \textbar{}

00,010,020,03 \textbar{} 100 \textbar{} 10,20,15 \textbar{} 10,34
\textbar{} l: \textbar{} : \textbar{} : \textbar{} ! \textbar{} :
\textbar{} 10,34 \textbar{} 10,20,15 \textbar{} 1000 \textbar{}

; \textbar{} \textbar{} : \textbar{} : \textbar{} : \textbar{} :
\textbar{} : \textbar{} : \textbar{} : \textbar{} : \textbar{} :
\textbar{} : \textbar{}

Теперь, когда мы знаем, що подразумевается под решением каждого из цих
типов, давайте сравним их между собой.

\section{\texorpdfstring{\textbf{Сравнительный анализ численного} и
аналитического
решений}{Сравнительный анализ численного и аналитического решений}}\label{ux441ux440ux430ux432ux43dux438ux442ux435ux43bux44cux43dux44bux439-ux430ux43dux430ux43bux438ux437-ux447ux438ux441ux43bux435ux43dux43dux43eux433ux43e-ux438-ux430ux43dux430ux43bux438ux442ux438ux447ux435ux441ux43aux43eux433ux43e-ux440ux435ux448ux435ux43dux438ux439}

Рассмотрим простейшую смещанную параболическую задачу

(40.1)

Задача решена. Теперь мы знаем, что нужно сделать, чтобы найти новые
координаты и привести уравнение к канонической форме.

\section{ЗАМЕЧАНИЯ}\label{ux437ux430ux43cux435ux447ux430ux43dux438ux44f}

\begin{enumerate}
\def\labelenumi{\arabic{enumi}.}
\tightlist
\item
  На самом деле для гиперболических уравнений существует две
  канонические формы. Вторую можно получить из первов заменой переменных
  вида
\end{enumerate}

Розв'язок якого показано на рис. 40.1. Що краще --- мати аналітичне
рішення цієї проблеми?

(40.2)

\begin{verbatim}
\alpha = \alpha (\xi, \eta) = \xi + \eta,
\end{verbatim}

\begin{verbatim}
или численное, представленное при  \alpha = 1  в табл. 40.1?
\end{verbatim}

Вопрос хороший, а ответ на него зависит от того, що мы хотим делать с
этим решением дальше. Однако каждому из типов решений присущи очевидные
преимущества.

\section{Преимущества аналитического
розв'язки}\label{ux43fux440ux435ux438ux43cux443ux449ux435ux441ux442ux432ux430-ux430ux43dux430ux43bux438ux442ux438ux447ux435ux441ux43aux43eux433ux43e-ux440ux43eux437ux432ux44fux437ux43aux438}

\begin{enumerate}
\def\labelenumi{\arabic{enumi}.}
\tightlist
\item
  Очевидно, що розв'язок (40.2) є більш інформативним, ніж таблиця
  чисел. Якщо ми хочемо обчислити розв'язок у певній точці (x, t), ми
  можемо зробити це так точно, як хочемо, просто збільшуючи кількість
  членів ряду, які розглядаються. Водночас легко отримати оцінку зверху
  щодо значення допущеної помилки.
\end{enumerate}

\pandocbounded{\includegraphics[keepaspectratio]{images/_page_301_Figure_4.jpeg}}

Рис. 40.1. Розв'язок рівняння теплопровідності (40.1) в различные
моменты времени.

\begin{itemize}
\tightlist
\item
  \begin{enumerate}
  \def\labelenumi{\arabic{enumi}.}
  \setcounter{enumi}{1}
  \tightlist
  \item
    Аналітичне розв'язання завжди дозволяє обчислити значення розв'язку
    в одній точці (x, t) без необхідності обчислювати значення
    розв'язків в інших точках, як це відбувається при розв'язанні задач
    із явними або неявними різницею схем.
  \end{enumerate}
\item
  \begin{enumerate}
  \def\labelenumi{\arabic{enumi}.}
  \setcounter{enumi}{2}
  \tightlist
  \item
    Аналітичне рішення дозволяє визначити розв'язок у будь-якій точці, а
    не лише на вузлах сітки.
  \end{enumerate}
\item
  \begin{enumerate}
  \def\labelenumi{\arabic{enumi}.}
  \setcounter{enumi}{3}
  \tightlist
  \item
    Найважливішою перевагою аналітичного розв'язку для нас є можливість
    простежити вплив фізичних параметрів, початкових і граничних умов на
    природу розв'язку.
  \end{enumerate}
\end{itemize}

Чисельні методи не виявляють цих регулярностей, оскільки дозволяють
знаходити розв'язок лише за заданих параметрів --- початкових і
граничних умовах. Іноді важливо знати взаємозв'язок між параметрами
моделі та розв'язком, особливо коли йдеться про оцінку фізичних
параметрів за типом розв'язку. Наприклад, припустимо, що розв'язок і
визначається експериментально, і ми знаємо аналітичний розв'язок

u = функція от параметров.

Тогда можна поставить вопрос об определении параметров нак функції
входных данных по схеме

Параметры = функція от u = функція исходных данных.

Такая задача називається параметрической идентификацией. Для
параметрической идентификации необходимо уметь решать уравнение з
частинними похідними. Несколько позже в цієї лекции мы приведем пример
параметрической идентификации в биологии, но сначала посмотрим, чем же
хороши численные розв'язки.

\section{Преимущества численных
решений}\label{ux43fux440ux435ux438ux43cux443ux449ux435ux441ux442ux432ux430-ux447ux438ux441ux43bux435ux43dux43dux44bux445-ux440ux435ux448ux435ux43dux438ux439}

Головна перевага чисельних розв'язків полягає в тому, що їх можна
отримати навіть якщо аналітичне рішення неможливо. Майже всі нелінійні
диференціальні рівняння в похідних похідних потрібно розв'язувати
чисельними методами, і більшість реальних моделей фізики, хімії,
біології тощо є нелінійними. Як правило, лінійні моделі наближаються до
нелінійних лише якщо в них відкинути нелінійні члени. Наприклад, ось
кілька нелінійних рівнянь:

\begin{enumerate}
\def\labelenumi{\arabic{enumi})}
\tightlist
\item
  нелинейное хвильове рівняння \texttt{u\_\{tt\}\ =\ u\_\{xx\}\ +}f(u)
\end{enumerate}

\begin{verbatim}
;

2) уравнение реакции с нелинейной диффузией  u_t = u_{xx} +f(u)
\end{verbatim}

;

\begin{enumerate}
\def\labelenumi{\arabic{enumi})}
\setcounter{enumi}{2}
\tightlist
\item
  система Ходжкина---Хаксли
\end{enumerate}

\texttt{Аналитические\ розв\textquotesingle{}язки\ для\ цих\ уравнений\ не\ известны\ ни\ при\ каких\ нелинейных\ функциях\ f\ и}
\texttt{g\^{}\{(1)\}} . Поэтому общий подход к решению нелинейных (а в
ряде случаев и линейных) задач обычно базируется на численных решениях.

Тепер давайте розглянемо приклад того, як можна використати аналітичне
рішення для визначення фізичних параметрів. Детальніший опис цих завдань
можна знайти у рекомендованій літературі.

\section{Параметрическая идентификация (в
биологии)}\label{ux43fux430ux440ux430ux43cux435ux442ux440ux438ux447ux435ux441ux43aux430ux44f-ux438ux434ux435ux43dux442ux438ux444ux438ux43aux430ux446ux438ux44f-ux432-ux431ux438ux43eux43bux43eux433ux438ux438}

Припустимо, що біолог намагається визначити швидкість, з якою
\texttt{(K\^{}+)} іони калію дифундують у екзоплазмовому розчині. Якщо
ви знаєте цей коефіцієнт дифузії, можна багато сказати про те, як
нервові імпульси передаються вздовж аксонів. Проблема в тому, що це
співвідношення майже неможливо знайти прямим вимірюванням. Однак можливо
знайти математичний

1) Це не так. В качестве примера можна привести уравнение
\texttt{u\_\{ii\}\ =\ u\_\{xx\}\ -\ \textbackslash{}sin\ u\_x\ -\ \textbackslash{}Pi\ pum}
, ped.

\texttt{связь\ между\ концентрацией\ поташа\ \textquotesingle{}u\textbackslash{}left(x,\textbackslash{},t\textbackslash{}right)\textquotesingle{}\ и\ коэффициентом\ диффузии\ \textquotesingle{}D,\textbackslash{},\textbackslash{},a\textbackslash{},\textbackslash{},\textquotesingle{}\ затем\ по\ измерениям\ величины\ \textquotesingle{}u\textbackslash{}left(x,\textbackslash{},t\textbackslash{}right)\textquotesingle{}\ определить\ величину\ D.}

Біологи Ходжкін і Кейз виявили, що після розміщення гігантського аксону
у спеціальному соляному розчині концентрація радіохктивного калію (**2K)
вздовж аксону приблизно описується кривою в початковий момент

.

Отже, параметри A і A слід обирати так, щоб крива відповідала
експериментальним даним (наприклад, у сенсі методу найменших квадратів).
Також було встановлено, що P И С

\pandocbounded{\includegraphics[keepaspectratio]{images/_page_303_Figure_6.jpeg}}

\begin{verbatim}
. 40.2. Начальная концентрация ионов (  ^{42}K  ): a—теоретическая кривая;  \delta  —экспериментальные точки.
\end{verbatim}

с течением времени ноны поташа под действием диффузии и конвекции
растекаются вдоль аксона. Ходжкин и Кейс предположили, що концентрация и
является решением следующей задачи с начальными данными для уравнения
диффузии:

(40.3)

Цю проблему вирішили у лекції 15, переключившись на рухому систему
координат. Якщо читач пам'ятає, загальна схема методу виглядає так.
Спочатку відкидаємо конвективний член і отримуємо розв'язок суто
дифузійної задачі (при V = 0):

TO

(23 9)

\begin{verbatim}
= \frac{A\sqrt{a}}{\sqrt{a+4Dt}} e^{-[x^2/(a+4Dt)]},
u_{\alpha\alpha} - u_{\beta\beta} = \frac{-\beta u_{\alpha} - \alpha u_{\beta}}{2\alpha\beta}.
u(x, t) = \frac{A\sqrt{a}}{\sqrt{a+4Dt}}e^{-[(x-Vt)^{2}/(a+4Dt)]}.
\end{verbatim}

\begin{verbatim}
При желании можно выразить   \alpha   и   \beta   через исходные координаты x и y. В результате получаем
\end{verbatim}

\begin{verbatim}
SS = \sum_{i=1}^{n} [y_i - (a + bx_i)]^2.
\alpha = \xi + \eta = (y^2 - x^2) + (y^2 + x^2) = 2y^3,
u_t = \frac{1}{r} \frac{\partial}{\partial r} \left( Dr u_r - s \omega^2 r^2 u \right), \quad 0 < r < 1,
\beta = \xi - \eta = (y^2 - x^2) - (y^2 + x^2) = -2x^2.
Au_{xx} + Bu_{xy} + Cu_{yy} + Du_x + Eu_y + Fu = G
u_{\xi\xi} - u_{\eta\eta} = \Psi(\xi, \, \eta, \, u_{\xi}, \, u, u_{\eta}).
\begin{cases} 0, & t < x/V, \\ P, & t \ge x/V. \end{cases}
3u_{xx} + 7u_{xy} + 2u_{yy} = 0
\end{verbatim}

Конечно, все це очень просто! Не сложнее, чем сбрасывать що-нибудь на
ленту транспортера, а затем наблюдать за его движением. Однако все
становится очень интересным, якщо примесь диффундирует в среде. Чтобы
разобраться в том, що про-исходит с движущейся волной при наличии
диффузии, решим

следующую задачу

(15.2)

\begin{verbatim}
u_{\xi\eta} = \Phi(\xi, \eta, u, u_1, u_2).
\end{verbatim}

З формальною заміною змінних
\texttt{\textbackslash{}xi\ =\ \textbackslash{}xi(x,\ y)} та
\texttt{\textbackslash{}eta\ =\ \textbackslash{}eta(x,\ y)} (41.1)
зводиться до тієї ж форми

(41.2)

\begin{verbatim}
u_{\alpha\alpha} - u_{\beta\beta} = \Psi(\alpha, \beta, u, u_{\alpha}, u_{\beta}).
\end{verbatim}

где, как обычно, \texttt{H(x)} --- функція Хевисайда Начальное
распределение концентрации изображено на рис 15 4.

\pandocbounded{\includegraphics[keepaspectratio]{images/_page_114_Figure_4.jpeg}}

Рис. 15.3 Конвективная волна a---передиції фронт примеси,
\texttt{\textbackslash{}delta} ---на одну единицу вперед волнового
фронта

Обратим внимание на то, що в задаче (15.2) мы отодвинули границу в ---
\texttt{\textbackslash{}infty} (теперь решается задача Коши), так що она
не оказывает теперь влияние на оценку розв'язок. Чтобы решить задачу
(15.2), можна воспользоваться преобразованием Лапласа Р И С

\pandocbounded{\includegraphics[keepaspectratio]{images/_page_114_Figure_7.jpeg}}

. 15 4 початкові умови в задаче конвективной диффузии.

ло переменной t или преобразованием Фурье по переменной x. Однако в
цьому случае гораздо интереснее поступить совсем по-другому Введем новую
систему координат, яка движется вдоль старой со скоростью V. Другими
словами, вместо системи координат, привязанной к берегу реки, мы
рассмотрим систему координат, яка движется со скоростью фронта примеси
(конечно, при налични диффузии фронт буде «смазан»). С точки зрения
математики це означает, що мы заменяем пространственную координату x на
новую \texttt{\textbackslash{}xi\ =\ x\ -\ Vt} .

Теперь ясно, що

когда \texttt{\textbackslash{}xi\ =\ 0} , мы находимся па фронте
распространяющейся примеси.

когда \texttt{\textbackslash{}xi\ =\ 1} , мы находимся на одну единицу
длины впереди фронта.

когда \texttt{\textbackslash{}xi\ =\ -1} , мы находимся на одну единицу
длины позади фронта.

Наша задача --- преобразовать исходную задачу с НУ

\begin{verbatim}
u_{xx} + 4u_{xy} = 0.
\overline{D} = A\xi_{xx} + B\xi_{xy} + C\xi_{yy} + D\xi_{x} + E\xi_{y},
u_{tt} = e^{2} (u_{xx} + u_{yy} + u_{zz}),
[\xi_x/\xi_y] = -B/2A
\xi(x, y) = c
\frac{dy}{dx} = -\left[\xi_x/\xi_y\right] = B_t 2A.
\end{verbatim}

(НУ)

\begin{verbatim}
\overline{B} = 2A\xi_x\eta_x + B(\xi_x\eta_y + \xi_y\eta_x) + 2C\xi_y\eta_y,
\end{verbatim}

1) Завдяки рівності

І з B\^{}2 - 4AC = 0 році , то

\texttt{Оскільки\ \textquotesingle{}\textbackslash{}xi\_x/\textbackslash{}xi\_y\ =\ -B\ \textbackslash{}\ 2A\ =\ -2\ \textbackslash{}sqrt\{A\textbackslash{}tilde\{C\}\}/2A\ =\ -\textbackslash{}sqrt\{C/A\}\textquotesingle{}\ останній\ вираз\ для\ B\ буде\ записаний\ як}

в трехмерном пространстве и показать, что полученное решение
удовлетворяет \emph{принципу Гюйгенса}. Для решения соответствующей
двумерной задачи

(УЧП)

І оскільки \texttt{\textbackslash{}overline\{A\}} дорівнює нулю, то
\texttt{\textbackslash{}overline\{B\}} , звісно, також дорівнює нулю.

Оскільки при нашому виборі \texttt{\textbackslash{}xi} обидва
коефіцієнти \texttt{\textbackslash{}overline\{A\}} і
\texttt{\textbackslash{}overline\{B\}} повертаються до нуля, ми можемо
взяти будь-яку функцію як \texttt{\textbackslash{}eta} (якщо вона не
пропорційна координаті \texttt{\textbackslash{}xi} ). У прикладі можна
вибрати \texttt{\textbackslash{}eta\ =\ y} .

Залишається знайти канонічну форму рівняння в нових координатах. Для
цього потрібно просто підставити формули \texttt{\textbackslash{}xi} і
\texttt{\textbackslash{}eta} (41.3) і знайти всі шанси
\texttt{\textbackslash{}overline\{A\}} ,
\texttt{\textbackslash{}overline\{B\}} ,
\texttt{\textbackslash{}overline\{C\}} ,
\texttt{\textbackslash{}overline\{D\}} ,
\texttt{\textbackslash{}overline\{E\}} ,
\texttt{\textbackslash{}overline\{F\}} і
\texttt{\textbackslash{}overline\{G\}} . На цьому вивчення параболічного
випадку завершується, але перш ніж перейти до іншої теми, розглянемо
простий приклад.

\section{\texorpdfstring{Перетворення параболічного рівняння
\texttt{u\_\{xx\}\ +\ 2u\_\{xy\}\ +\ u\_\{yy\}\ =\ 0} у канонічну
форму}{Перетворення параболічного рівняння u\_\{xx\} + 2u\_\{xy\} + u\_\{yy\} = 0 у канонічну форму}}\label{ux43fux435ux440ux435ux442ux432ux43eux440ux435ux43dux43dux44f-ux43fux430ux440ux430ux431ux43eux43bux456ux447ux43dux43eux433ux43e-ux440ux456ux432ux43dux44fux43dux43dux44f-u_xx-2u_xy-u_yy-0-ux443-ux43aux430ux43dux43eux43dux456ux447ux43dux443-ux444ux43eux440ux43cux443}

Розглянемо рівняння

\begin{verbatim}
u_{tt} = c^2 (u_{xx} + u_{yy}), \quad
\end{verbatim}

Його коефіцієнти мають такі значення: A=1, B=2, C=1, D=E=F=G=0. Отже,
\texttt{B\^{}2-4AC=0} для всіх значень x і y. Щоб знайти нові координати
\texttt{\textbackslash{}xi} і \texttt{\textbackslash{}eta} та канонічну
форму рівняння, ми рухаємося наступним чином.

КРОК 1. Запишемо \emph{характеристичне рівняння} (тепер лише одне, а не
два)

\begin{verbatim}
\begin{cases} -\infty < x < \infty, \\ -\infty < y < \infty, \end{cases}
\end{verbatim}

Розв'язуючи це рівняння відносно y, отримуємо

\begin{verbatim}
\begin{cases} -\infty < x < \infty, \\ -\infty < y < \infty, \\ -\infty < z < \infty, \end{cases}
\end{verbatim}

,

Де --- характеристика \texttt{\textbackslash{}xi\ =\ y\ -\ x} визначена.
Використання цієї характеристики забезпечує досягнення рівності
\texttt{\textbackslash{}overline\{A\}\ =\ 0} . Координату
\texttt{\textbackslash{}eta} можна вибрати довільно (якщо вона не
залежить від \texttt{\textbackslash{}xi} ). Ми оберемо його наступним
чином:

\begin{verbatim}
\begin{cases} u(x, y, z, 0) = \varphi(x, y, z), \\ u_{t}(x, y, z, 0) = \psi(x, y, z). \end{cases}
\end{verbatim}

.

Нові координати

\begin{verbatim}
P_{2}
\end{verbatim}

изображены на рис. 41 1.

\pandocbounded{\includegraphics[keepaspectratio]{images/_page_309_Figure_3.jpeg}}

Рис. 41.1. Новая система координат
\texttt{\textbackslash{}xi\ =\ y\ -\ x} ,
\texttt{\textbackslash{}eta\ =\ y} .

ШАГ 2. Осталось найти каноническую форму уравнения. Подставляя
\texttt{\textbackslash{}xi} и \texttt{\textbackslash{}eta} в формулы для
коефіцієнтів \texttt{\textbackslash{}overline\{A\}} ,
\texttt{\textbackslash{}overline\{B\}} ,
\texttt{\textbackslash{}overline\{C\}} ,
\texttt{\textbackslash{}overline\{D\}} ,
\texttt{\textbackslash{}overline\{E\}} ,
\texttt{\textbackslash{}overline\{F\}} и
\texttt{\textbackslash{}overline\{G\}} , получаем
\texttt{\textbackslash{}overline\{A\}\ =\ 0} (так и должно быть; мы
находили \texttt{\textbackslash{}xi} из условия равенства нулю цього
коэффициента), \texttt{\textbackslash{}overline\{B\}=0} (мы уже раньше
показали, що этог коефіцієнт також обращается в нуль),

(41.5)

\begin{verbatim}
\begin{aligned} u_{tt} &= c^2 \Delta u, \\ (HY) & \begin{cases} u(x, y, z, 0) = 0, \\ u_t(x, y, z, 0) = \psi(x, y, z), \end{cases} \end{aligned}
\overline{D} = A \xi_{xx} + B \xi_{xy} + C \eta_{yy} + D \xi_x + E \xi_y = 0,
\end{verbatim}

\begin{verbatim}
Цю задачу можна розв'язати методом  \Phi  Юр'єра  ^{1}  ), і виявляється, що розв'язок можна представити у вигляді
\end{verbatim}

Отже, нове рівняння

\begin{verbatim}
P_{1}
\xi = x - Vt
\end{verbatim}

\begin{verbatim}
где  \overline{\psi}  — *среднее* значение начального распределения  \psi  *по сфере* радиуса ct с центром в точке (x, y, z), т. е.
\end{verbatim}

\begin{verbatim}
u(x, t) =
\end{verbatim}

\pandocbounded{\includegraphics[keepaspectratio]{images/_page_177_Picture_3.jpeg}}

є загальним розв'язком рівняння П е

\pandocbounded{\includegraphics[keepaspectratio]{images/_page_177_Picture_7.jpeg}}

ред нами розв'язок задачи конвективной диффузии (15.2). Его очень легко
интерпретировать, якщо представить, що мы движемся относительно графика,
изображенного на рис. 15.6.

\pandocbounded{\includegraphics[keepaspectratio]{images/_page_117_Figure_3.jpeg}}

Рис. 15.6. Диффузня из области высокой концентрации в область низкой
концентрации. Чем больше коефіцієнт дифузії, тем быстрее установится
стационарное значение.

Другими словами, в зависимости от относительной величины D (коэффициента
диффузии) и V (скорости потока) розв'язок движется слева направо со
скоростью V и в то же самое время передний фронт расплывается со
скоростью, определяемой величиной D.

\pandocbounded{\includegraphics[keepaspectratio]{images/_page_117_Figure_6.jpeg}}

Рис. 15.7. Розв'язок задачи конвективнои диффузии. Вещество одновременно
движется и диффундирует.

(На рис. 15.7 показано, как расплывается передний фронт концентрации.)

\section{ЗАМЕЧАНИЯ}\label{ux437ux430ux43cux435ux447ux430ux43dux438ux44f-1}

Преобразование координат является важным методом розв'язки уравнений з
частинними похідними. Выбрав подходящую систему координат, можна
существенно упростить уравнение.

\section{ЗАДАЧИ}\label{ux437ux430ux434ux430ux447ux438}

\begin{enumerate}
\def\labelenumi{\arabic{enumi}.}
\tightlist
\item
  Решите задачу Коши:
\end{enumerate}

\begin{verbatim}
\begin{aligned} (\text{YAII}) & u_{tt} = c^2 \Delta u, \\ (\text{HY}) & \begin{cases} u = 0, \\ u_t = \varphi, \end{cases} \end{aligned}
\end{verbatim}

Але тепер у \texttt{B\^{}2-4AC\textless{}0} . Переходячи до нових
пояснювальних змінних, ми хочемо перетворити її у нову форму

\begin{verbatim}
u_h(r, \theta) = \sum_{n=0}^{\infty} r^n [a_n \cos(n\theta) + b_n \sin(n\theta)].
\end{verbatim}

\begin{enumerate}
\def\labelenumi{\arabic{enumi}.}
\setcounter{enumi}{1}
\tightlist
\item
  Найдите розв'язок задачи Коши для уравнения конвективной диффузии:
\end{enumerate}

\begin{verbatim}
\Delta u = -\frac{r^2}{2} - \frac{r^3}{2} \cos{(2\theta)}
\end{verbatim}

Це можливо, але лише шляхом введення складних координат. Щоб знайти ці
комплексні координати \texttt{\textbackslash{}xi} і
\texttt{\textbackslash{}eta} , розв'яжемо характеристичні рівняння, так
само як у гіперболічному випадку:

\begin{verbatim}
u = \frac{\partial}{\partial t} [t\overline{\varphi}].
\end{verbatim}

Воспользуйтесь преобразованием из лекции 8.

\begin{enumerate}
\def\labelenumi{\arabic{enumi}.}
\setcounter{enumi}{2}
\tightlist
\item
  Найдите розв'язок задачи переноса:
\end{enumerate}

\begin{verbatim}
u_p(r, \theta) = Ar^4 + Br^4 \cos(2\theta).
\frac{dy}{dx} = \sqrt{-4x^2} = 2ix
u(x, t) = \frac{1}{2c} \int_{x-ct}^{x+ct} \varphi(s) ds.
\\ -\infty < x < \infty, \quad 0 < t < \infty,
u_p(r, \theta) = -\frac{r^4}{32} - \frac{r^4}{24} \cos(2\theta).
u_t = Du_{xx} - Vu_x, \quad -\infty < x < \infty, \quad 0 < t < \infty,
u_t(x, t) = \frac{1}{2} [\varphi(x + ct) + \varphi(x - ct)].
u(x, 0) = e^{-x^2}, \quad -\infty < x < \infty.
\sum_{n=0}^{\infty} [a_n \cos(n\theta) + b_n \sin(n\theta)] - \frac{1}{32} - \frac{1}{24} \cos(2\theta) = 0.
y^2 u_{xx} + x^2 u_{yy} = 0,
u_1(r, \theta) = \frac{1}{32} + \frac{1}{24}r^2\cos(2\theta) - \frac{r^4}{24}\cos(2\theta) - \frac{r^4}{32} = \frac{(r^4 - 1)}{32} - \frac{(r^4 - r^2)}{24}\cos(2\theta).
\frac{dy}{dx} = \frac{B - \sqrt{B^2 - 4AC}}{2A} = -\frac{\sqrt{-4x^2y^2}}{2y^2} = -i\frac{x}{y},
u = u_0 + u_1 = r \cos \theta - \frac{(r^4 - 1)}{32} - \frac{(r^4 - r^2)}{24} \cos (2\theta).
u(\xi, \tau) = \frac{1}{2\sqrt{D\pi\tau}} \int_{-\infty}^{+\infty} e^{-\beta^2} e^{-(\xi-\beta)^2/4D\tau} d\beta.
u_{tt} = \alpha^2 u_{xx}
u_{\xi\eta} = \psi(\xi, \eta, u, u_{\xi}, u_{\eta}).
P_{2}
\xi(x, y) = y^2, \quad \eta(x, y) = x^2.
\Delta u = 0, \quad 0 < r < 1 + \frac{1}{4} \sin \theta,
T \sin \theta_2 - T \sin \theta_1 =
\end{verbatim}

записывается в виде

\begin{verbatim}
\Delta x \rho u_{tt} = T \left[ u_x \left( x + \Delta x, \ t \right) - u_x \left( x, \ t \right) \right] + \Delta x F \left( x, \ t \right) - \frac{1}{\Delta x \rho u_t} \left( x, \ t \right) - \Delta x \gamma \dot{u} \left( x, \ t \right)
u \left( 1 + \frac{1}{4} \sin \theta, \ \theta \right) = \cos \theta, \quad 0 \le \theta \le 2\pi.
u_{tt} = \alpha^{3} u_{xx} - \beta u_{t} - \gamma u + F(x, t).
\end{verbatim}

\begin{verbatim}
где  \overline{\phi}  и  \overline{\psi}  — средние значения функций  \phi  и  \psi  на сфере раднусом ct с центром в точке (x, y, z).
\end{verbatim}

\begin{verbatim}
Цей розв'язок зазвичай називають формулою Пуассона для хвильового рівняння у тривимірному просторі. Це природне узагальнення формули д'Аламбера на тривимірний випадок. Найважливішим у формулі Пуассона є те, що обидва інтеграли в  \bar{\phi}  і  \bar{\psi}  взяті з поверхні сфери. Цей факт дозволяє нам надати таку важливу інтерпретацію рішення.
\end{verbatim}

\begin{verbatim}
У момент  t=t_1  розв'язок u у (x,y,z) залежить лише від початкових розподілів сфери з радіусом  ct_1  центрованим на (x,y,z) (рисунок 24.3). Тепер припустимо, що початкові розподіли  \phi  і  \psi  дорівнюють нулю всюди, крім малої сфери (див. рисунок 24.3). З часом радіус P И С
\end{verbatim}

\pandocbounded{\includegraphics[keepaspectratio]{images/_page_179_Picture_3.jpeg}}

. 24.3. Схема поясияет, как начальное возмущение влияет на точку (к, ч,
г).

\begin{verbatim}
Сфера, що оточує точку (x, y, z), зростає зі швидкістю c, так що через  t_{\hat{z}}  секунди ця сфера починає перетинати область початкового збурення, і тому розв'язок u(x, y, z, t) стає ненульовим. У  t_{\hat{z}} < t < t_{\hat{a}}  розв'язок U у (x, y, z) залишається ненульовим, оскільки сфера продовжує перетинати область початкового збурення, але на  t=t_{\hat{a}}  розв'язок знову стає нульовим. Іншими словами, розв'язок, що поширюється з початкової області збурення, має чітко визначену задню кромку. Цей загальний принцип відомий як принцип Гюйгенса для тривимірного простору. Виявляється, що передній край хвилі завжди чітко окреслений, але задній край чітко окреслений лише в просторах розмірності 3, 5, 7.... Ми вже знаємо з формули д'Аламбера, що початкові умови
\end{verbatim}

\pandocbounded{\includegraphics[keepaspectratio]{images/_page_349_Figure_5.jpeg}}

\begin{verbatim}
(16.2) u_{tt} = \alpha^2 u_{xx},
u\left(1+\frac{1}{4}\sin\theta,\ \theta\right)=\cos\theta
u_{tt} = ku_{xx}
\end{verbatim}

\begin{verbatim}
приводят, вообще говоря, к нерезкому заднему фронту (поскольку формула Даламбера подразумевает интегрирование функции  \psi(x)  в пределах от x-ct до x+ct).
\end{verbatim}

Покажем теперь, что принцип Гюйгенса неприменим к цилиндрическим волнам.
Такого рода решения описывают волны на воде от точечного источника,
когда задний фронт волны не резкий, а постепенно затухает до нуля.

\section{Двумерное волновое
уравнение}\label{ux434ux432ux443ux43cux435ux440ux43dux43eux435-ux432ux43eux43bux43dux43eux432ux43eux435-ux443ux440ux430ux432ux43dux435ux43dux438ux435}

Чтобы решить двумерную волновую задачу

(УЧП)

\begin{verbatim}
u_{tt} = \frac{\partial}{\partial x} [\alpha^2(x) u_x],
u(1+\epsilon\sin\theta, \theta) = \cos\theta, \quad 0 \le \epsilon \le 1/4.
u(x, 0) = f(x)
f(x+h) = f(x) + f'(x)h + f''(x)\frac{h^2}{2!} + \dots,
i_x + Cv_t + Gv = 0, \\ v_x + Li_t + Ri = 0,
i_{xx} + Gv_x - CLi_{tt} - CRi_t = 0.
v_{x} = -Li_{t} - Ri,
\end{verbatim}

\pandocbounded{\includegraphics[keepaspectratio]{images/_page_350_Figure_2.jpeg}}

\begin{verbatim}
i_{t,s} = CLi_{tt} + (CR + GL)i_t + GRi.
u(1+\varepsilon\sin\theta,\ \theta)=u(1,\ \theta)+u_r(1,\ \theta)(\varepsilon\sin\theta)+u_{rr}(1,\ \theta)\frac{(\varepsilon\sin\theta)^2}{21}+\ldots
v_{xx} = CLv_{tt} + (CR + GL)v_t + GRv.
\Delta u = 0, \quad 0 < r < 1 + \varepsilon \sin \theta,
v_{tt} = \alpha^2 v_{xx}, \\ i_{tt} = \alpha^2 i_{xx}, \qquad \alpha^2 = 1/CL,
u(1, \theta) + u_r(1, \theta) (\varepsilon \sin \theta) + u_{rr}(1, \theta) \frac{(\varepsilon \sin \theta)^2}{2!} + \dots = \cos \theta.
u_{tt} = u_{xx} - u_t
(46.8) u = u_0 + \varepsilon u_1 + \varepsilon^2 u_2 + \dots.
\end{verbatim}

Що ви можете сказати про точність?

\begin{enumerate}
\def\labelenumi{\arabic{enumi}.}
\setcounter{enumi}{3}
\tightlist
\item
  Як отримати послідовність випадкових точок, що лежать всередині
  трикутника T?
\end{enumerate}

\pandocbounded{
\includegraphics[keepaspectratio]{images/_page_319_Figure_13.jpeg}}

\pandocbounded{\includegraphics[keepaspectratio]{images/_page_319_Figure_14.jpeg}}

\begin{enumerate}
\def\labelenumi{\arabic{enumi}.}
\setcounter{enumi}{4}
\tightlist
\item
  Як отримати послідовність випадкових чисел, розподіл яких показаний на
  рисунку?
\end{enumerate}

\pandocbounded{\includegraphics[keepaspectratio]{images/_page_319_Figure_16.jpeg}}

Іншими словами, як отримати послідовність чисел \{0, 1, 2\}, якщо
ймовірність виникнення 0 і 2 дорівнює 0,25, а ймовірність появи 1 ---
0,5?

Лекція 43

\section{РОЗВ'ЯЗАННЯ ДИФЕРЕНЦІАЛЬНИХ РІВНЯНЬ У ПОХІДНИХ ПОХІДНИХ
МОНТЕ-КАРЛО}\label{ux440ux43eux437ux432ux44fux437ux430ux43dux43dux44f-ux434ux438ux444ux435ux440ux435ux43dux446ux456ux430ux43bux44cux43dux438ux445-ux440ux456ux432ux43dux44fux43dux44c-ux443-ux43fux43eux445ux456ux434ux43dux438ux445-ux43fux43eux445ux456ux434ux43dux438ux445-ux43cux43eux43dux442ux435-ux43aux430ux440ux43bux43e}

МЕТА ЛЕКЦІЇ: Показати, як побудувати таку гру, результатом якої є
приблизне розв'язання диференціального рівняння. Одна з цих ігор (модель
випадкової ходьби) веде до скінченної різниці апроксимації задачі
Діріхле в квадраті. Цю модель можна узагальнити для отримання розв'язків
інших задач.

У попередній лекції було сказано, що можливо винайти такі ігри на удачу,
результатом яких є розв'язки (можливо, приблизні) задач для рівнянь з
частковими коефіцієнтами проникнення. У цій лекції ми покажемо, як
побудувати таку гру для наближеного розв'язання задачі Діріхле (PPP)
\texttt{u\_\{xx\}\ +\ u\_\{yy\}\ =\ 0} , 0 \textless{} x \textless{} 1,
0 \textless{} y \textless{} 1,

(GI)

\begin{verbatim}
P_{0}
R(A) = g_1 P_A(p_1) + g_2 P_A(p_2) + \dots + g_{12} P_A(p_{12}).
\begin{cases} \Delta u_{0} = 0, & 0 < r < 1 \text{ (внутри круга),} \\ u_{0}(1, \theta) = \cos \theta, & u_{0}(r, \theta) = r \cos \theta, \end{cases}
R(A) = 1 \cdot (0.04) + 1 \cdot (0.15) + 1 \cdot (0.03) + 0 \cdot (0.06) + \dot + 0 \cdot (0.04).
\begin{cases} u(x, \ 0) = \varphi(x), \\ u_t(x, \ 0) = 0 \end{cases}
R(A) = \frac{1}{4} [R(B) + R(C) + R(D) + R(E)]
P_{1}
u_{i,j} = \frac{1}{4} \left[ u_{i-1,j} + u_{i+1,j} + u_{i,j-1} + u_{i,j+1} \right],
\begin{cases} \Delta u_{1} = 0, & 0 < r < 1 \text{ (внутри круга),} \\ u_{1}(1, \theta) = -\sin \theta \frac{\partial u_{0}(1, \theta)}{\partial r} = -\sin \theta \cos \theta \end{cases}
\end{verbatim}

(PPP)

\begin{verbatim}
u_{xx}=u_{l,\ j+1}-2u_{l,\ j}+u_{l,\ j-1},
u_{i, j} = \frac{u_{i, j+1} + u_{i, j-1} + \sin x_j (u_{i+1, j} + u_{i-1, j})}{2 (1 + \sin x_j)}.
u = u_0 + \frac{1}{4}u_1 + \left(\frac{1}{4}\right)^2 u_1 + \cdots
\end{verbatim}

(i, j+1)

\begin{verbatim}
u_u + \frac{1}{4} u_t
u(A) = g_1 P_A(p_1) + g_2 P_A(p_2) + \dots + g_N P_A(p_N),
\end{verbatim}

\begin{enumerate}
\def\labelenumi{\arabic{enumi}.}
\setcounter{enumi}{1}
\tightlist
\item
  Результаты задачи 1 используйте для решения следующей задачи:
\end{enumerate}

(УЧП)

\begin{verbatim}
u_{xx} + u_{yy} = 0
u_t = (1+x) u_{xx}, \\ u(x, 0) = q(x), \quad -\infty < x < \infty,
u_{xx} + x^2 u_{yy} = 0
\end{verbatim}

(НУ)

\begin{verbatim}
(\text{UBP}) \qquad
u_1 = (1 + \varepsilon_1) u_{xx}
\end{verbatim}

\begin{enumerate}
\def\labelenumi{\arabic{enumi}.}
\setcounter{enumi}{4}
\tightlist
\item
  Разработайте схему метода Монте-Карло для розв'язки смешанной
  параболической задачи
\end{enumerate}

\begin{verbatim}
u = u_1 + \varepsilon u_1 + \varepsilon^2 u_2 + \cdots
\end{verbatim}

\begin{verbatim}
\right. \end{array}
\end{verbatim}

\begin{verbatim}
J[y] = \int_{b}^{a} F(x, y, y') dx,
J[y] = \int_{0}^{1} [y^{2}(x) + y'^{2}(x)] dx.
P_{0}
\frac{df(x)}{dx} = 0
\end{verbatim}

\begin{verbatim}
,  (x,  y, z) \in R^3  ,  u = 0   u_t =
\end{verbatim}

\begin{verbatim}
T = \int_{0}^{T} dt = \int_{a}^{L} \frac{dt}{ds} ds = \int_{0}^{L} \frac{ds}{v} = \frac{1}{\sqrt{2mg}} \int_{0}^{L} \frac{ds}{\sqrt{y}} = \frac{1}{\sqrt{2mg}} \int_{a}^{b} \sqrt{\frac{1+y'^{2}}{y}} dy
g(x) = e^{-x^2}.
J[y] = \int_{a}^{b} F(x, y, y') dx.
\end{verbatim}

\begin{enumerate}
\def\labelenumi{\arabic{enumi}.}
\setcounter{enumi}{3}
\tightlist
\item
  Решите аналогичную (см. задачу 3) задачу для двумерного уравнения
\end{enumerate}

(УЧП)

\begin{verbatim}
\frac{df(x)}{dx} = 0
\end{verbatim}

Зверніть увагу, що в перших двох задачах коефіцієнти є сталими. Читач
повинен розуміти, що параметр b має бути достатньо малим, а необережний
ряд може відрізнятися

\section{ЗАДАЧИ}\label{ux437ux430ux434ux430ux447ux438-1}

\begin{verbatim}
- 1. Подставьте разложение (46.8) в задачу (46.7) и получите последовательность задач P_0,\ P_1,\ P_2,\ \dots . 2. Покажите, что нелинейную задачу
\end{verbatim}

\begin{verbatim}
J[y] = \int_{a}^{y} F(x, y, y') dx
\end{verbatim}

(НУ)

\begin{verbatim}
J[y] = \int_{a}^{b} F(x, y, y') dx
u(r, \theta) = u_0(r, \theta) + \frac{1}{4}u_1(r, \theta)
y(a) = A
\Delta u = 0,
J[\overline{y}] \leqslant J[\overline{y} + \varepsilon \eta]
q_{xx} + \varphi_{uu} = 0
\varphi(\varepsilon) = J[\bar{y} + \varepsilon \eta]
q_{uu} + q_{vv} = 0
\frac{d\varphi\left(\varepsilon\right)}{d\varepsilon} = \frac{d}{d\varepsilon} J\left[\overline{y} + \varepsilon\eta\right]|_{\varepsilon=0} =
\end{verbatim}

\pandocbounded{\includegraphics[keepaspectratio]{images/_page_353_Figure_9.jpeg}}

(Читач має робити це самостійно.) Інтегруючи на частини, ми отримуємо

\begin{verbatim}
w = f(z)
\end{verbatim}

Оскільки цей інтеграл є нульовим у будь-якій функції
\texttt{\textbackslash{}eta(x)} , що задовольняє граничні умови
\texttt{\textbackslash{}eta(a)\ =\ \textbackslash{}eta(b)\ =\ 0} , ми Р
и с

\pandocbounded{\includegraphics[keepaspectratio]{images/_page_330_Figure_6.jpeg}}

. 44 3. Граф функції
\texttt{J{[}y+{]}\ \textbackslash{}in\ \textbackslash{}eta} {]} в
околиці \texttt{\textbackslash{}varepsilon\ =\ 0} .

Ми приходимо до висновку, що друга частина підоб'єктивного виразу має
дорівнювати нулю, тобто

(44 2)

\pandocbounded{\includegraphics[keepaspectratio]{images/_page_355_Figure_5.jpeg}}

(Рівняння Ейлера--Лагранжа).

Рівняння (44 2) називається рівнянням Ейлера-Лагранжа, і хоча загалом
воно здається комплексним, коли в нього підставляється конкретна функція
F(x, y, y'), воно стає звичайним диференціальним рівнянням другого
порядку відносно невідомої функції
\texttt{\textbackslash{}overline\{y\}(x)} . Отже, щоб визначити
мінімізуючу функцію \texttt{\textbackslash{}overline\{y\}} , необхідно
розв'язати рівняння Ейлера-Лагранжа.

Отже, ми показали, що

якщо функція \texttt{y(x)} мінімізує функціональний
\texttt{J{[}y{]}\ =\ \textbackslash{}int\_a\^{}b\ F(x,\ y,\ y\textquotesingle{})\ dx}
(у класі гладких функцій з граничними умовами \texttt{y(a)} = A;
\texttt{y(b)} = B), то вона повинна задовольняти рівняння

\pandocbounded{\includegraphics[keepaspectratio]{images/_page_355_Figure_7.jpeg}}

(з граничною умовою).

(Щоб спростити позначення, ми прибрали лінію над
\texttt{\textbackslash{}bar\{y\}} .) Щоб проілюструвати теорію, наведемо
приклад.

Пошук мінімальної функціональності
\texttt{J{[}y{]}\ =\ \textbackslash{}int\_\{0\}\^{}\{1\}\ {[}y\^{}2\ +\ y\textquotesingle{}\^{}2{]}\ dx}
.

Спробуємо знайти функцію \texttt{y(x)}, яка проходить через точки (0,0)
і (1,1) і мінімізує функціональність

\begin{verbatim}
u + iv = (x + iy)^2 = x^2 - y^2 + 2ixy
\end{verbatim}

З контексту зрозуміло, що шукана функція має бути диференційовною,
оскільки інтегративний вираз залежить від y'. Щоб визначити
\texttt{\textbackslash{}bar\{y\}(x)} , запишемо рівняння
Ейлера--Лагранжа (разом із граничною умовою \texttt{y(0)}=0 та
\texttt{y(1)}=1).

\begin{verbatim}
\left\{
\end{verbatim}

\texttt{y(0)\ =\ 0,} \texttt{y(1)\ =\ 1.} Але
\texttt{F(x,\ y,\ y\textquotesingle{})\ =\ y\^{}2\ +\ y\textquotesingle{}\^{}2}
і, отже,

Чи знаєте ви, як виглядає графік цієї функції з сегментом {[}0, 1{]}?
Подумай про це. Зверніть увагу, що синусоїдальне перетворення функції
\texttt{f(x)} визначається лише для додатних цілих чисел n.~Іншими
словами, скінченний синус і косинус перетворюють функції у чисельні
послідовності.

\pandocbounded{\includegraphics[keepaspectratio]{images/_page_185_Figure_3.jpeg}}

РИС. 25.1. Графики функции \texttt{f(x)} = 1 и ее преобразования.

\section{Свойства
преобразований}\label{ux441ux432ux43eux439ux441ux442ux432ux430-ux43fux440ux435ux43eux431ux440ux430ux437ux43eux432ux430ux43dux438ux439-1}

Прежде чем приступить к решению задач, мы должны получить некоторые
свойства этих преобразований.

\begin{verbatim}
Якщо you(x, t) — функція  \partial \theta yx  змінних, і ми виконуємо перетворення змінної x, то
\end{verbatim}

(Продифференцировали \texttt{F} по \texttt{y} и
\texttt{y\textquotesingle{}.} )

Уравнение Эйлера --- Лагранжа принимает вид

(Зверніть увагу, що перетворення було виконано на змінній x, і отримана
послідовність залежить лише від часу t.)

А как быть с производными? Ниже приводятся несколько полезных формул

откуда

Решая це простое дифференциальное уравнение с граничными условиями
\texttt{y(0)}=0 и \texttt{y(1)}=1, получаем

.

График цієї функції изображен на рис. 44.4. Ясно, що любая другая
гладкая функція, удовлетворяющая тем же граничным условиям, даст большее
значение функционала \texttt{J(y)}.

\pandocbounded{\includegraphics[keepaspectratio]{images/_page_332_Figure_5.jpeg}}

Рис. 44 4. Все допустимые гладкие кривые удовлетворяют граничным
условиям \texttt{y(0)} = 0 и \texttt{y(1)} = 1: a---допустимая кривая;
\texttt{\textbackslash{}delta} ---минимизирующая функція
\texttt{y(x)\ =\ 0.42e\^{}x\ -\ 0.42e\^{}\{-x\}} . \# ЗАУВАЖЕННЯ 1.
Уравнение Эйлера---Лагранжа аналогично равенству нулю производной в
дифференциальном исчислении. Якщо читатель помнит, то на самом деле
равенство нулю производной не является достаточным условием экстремума.
Наприклад, производная функції \texttt{f(x)\ =\ x\^{}3} при x = 0 равна
нулю, но ця точка не является нп точкой минимума, ни точкой максимума.
Так же и с уравнением Эйлера---Лагранжа: оно является лише необходимым
условием экстремума, но не достаточным. Значит, экстремум может
достигаться на других функциях. Очень часто, однако, розв'язок уравнения
Эйлера---Лагранжа по существу задачи дает локальный или даже глобальный
минимум.

\begin{enumerate}
\def\labelenumi{\arabic{enumi}.}
\setcounter{enumi}{1}
\tightlist
\item
  Якщо у (х) минимизирует функционал
\end{enumerate}

то площадь поверхности, образованной вращением кривой \texttt{y(x)}
вокруг оси x, буде минимальной (рис. 44.5). Решением уравнения Эйлера
--- Лагранжа в цьому случае буде

цепная линия

где константы \texttt{\textbackslash{}alpha} и
\texttt{\textbackslash{}beta} определяются из условий \texttt{y(a)} = A
и \texttt{y(b)} = B (т. е. цепная линия должна проходить через заданные
точки). Решить уравнение Эйлера---Лагранжа в данном случае достаточно
трудно, поэтому мы рекомендуем читателю лише проверить, що цепная линия
ему удовлетворяет.

\begin{enumerate}
\def\labelenumi{\arabic{enumi}.}
\setcounter{enumi}{2}
\tightlist
\item
  В цієї лекции мы нашли, що функція
\end{enumerate}

доставляет минимум функционалу

\begin{verbatim}
(\text{УЧП}) \quad u_{tt} = u_{xx} + \sin(\pi x), \quad 0 < x < 1, \quad 0 < t < \infty,
\end{verbatim}

\begin{verbatim}
При цьому оказывается, що  J[\overline{y}] = 0.46  . Якщо подставить в функционал любую другую гладкую функцию, график которой проходит через точки (0, 0) и (1, 1), то значение функционала J[y] станет больше.
\end{verbatim}

\pandocbounded{\includegraphics[keepaspectratio]{images/_page_333_Picture_8.jpeg}}

Рис. 44 5. Минимальная поверхность вращения.

\begin{enumerate}
\def\labelenumi{\arabic{enumi}.}
\setcounter{enumi}{3}
\tightlist
\item
  Основные законы физики чаще всего формулируются на языке вариационных
  принципов, а не дифференциальных уравнений. В качестве примеров можна
  привести принцип Ферма (свет при распространении из одной точки в
  другую выбирает путь, которому соответствует наименьшее время
  распространения) или принцип Гамильтона (в консервативном поле частица
  движется так, що интеграл действия
\end{enumerate}

\begin{verbatim}
(\text{ГУ}) \quad
\end{verbatim}

(кинетическая энергия--- потенциальная энергия) dt буде минимальным).
Таким образом, физические явления развиваются лише так, що эти
функционалы принимают минимальные значения.

Основные идеи вариационного исчисления можна распространить на
функционалы, зависящие от функції нескольких независимых змінних, вида

Соответствующее уравнение Эйлера --- Лагранжа для такого функционала
буде уже уравнением з частинними похідними

Такого рода функционалы мы рассмотрим в следующей лекции. Однако общая
философия метода буде иной. Некоторым новым способом мы найдем функцию
\texttt{u(x,\ y)}, яка минимизирует функционал, а поскольку она является
решением уравнения Эйлера --- Лагранжа, то тем самым мы найдем розв'язок
уравнения з частинними похідними. Другими словами, мы будем искать
розв'язок уравнения з частинними похідними методом минимизации
функционала. В цієї лекции все було наоборот: задачу минимизации
функционала, зависящего от функції одной переменной, мы сводили к
решению уравнения Эйлера --- Лагранжа. Методы розв'язки дифференциальных
уравнений путем минимизации соответствующих функционалов принято
называть прямыми методами вариационного исчисления. Методы минимизации
функционалов путем розв'язки соответствующих уравнений Эйлера ---
Лагранжа називаються непрямыми методами вариационного исчисления.
Следующая лекция посвящена хорошо известному прямому методу Ритца.

\section{ЗАДАЧІ}\label{ux437ux430ux434ux430ux447ux456}

\begin{enumerate}
\def\labelenumi{\arabic{enumi}.}
\tightlist
\item
  Среди всех кривых, удовлетворяющих условиям \texttt{y(0)} = 0, и
  \texttt{y(1)} = 1, найдите ту, яка минимизирует функционал
\end{enumerate}

Как вы интерпретируете получелный результат? Чему равна величина
\texttt{J{[}\textbackslash{}overline\{y\}{]}} ? Каков смысл величины
\texttt{J{[}\textbackslash{}overline\{y\}{]}} ?

\texttt{2.\ Кинетическая\ энергия\ колеблющейся\ материальной\ точки\ определяется\ выражением\ \textquotesingle{}KE\ =\ \textbackslash{}frac\{1\}\{2\}\ m\textbackslash{}dot\{y\}\^{}2\textquotesingle{}\ ,\ где\ y\ =\ dy/dt,\ а\ потенциальная\ равна\ \textquotesingle{}PE\ =\ \textbackslash{}frac\{1\}\{2\}\ ky\^{}2\textquotesingle{}\ .}

\pandocbounded{\includegraphics[keepaspectratio]{images/_page_335_Picture_2.jpeg}}

Принцип Гамильтона утверждает, що материальная точка буде двигаться так,
що интеграл

буде минимален. Якщо цей закоп справедлив, то получите дифференциальное
уравнение движения материальной точки.

\begin{verbatim}
3. Покажите, що в классе гладких функцій, удовлетворяющих граничным условиям y(0)=0 и  y(\pi/2)=1  , минимум функционала
\end{verbatim}

достигается на функції
\texttt{\textbackslash{}overline\{y\}(x)\ =\ \textbackslash{}sin\ x} .
Вычислите \texttt{J(\textbackslash{}sin\ x)} . 4. Исходя из функционала

получите уравнение Эйлера --- Лагранжа

\# ВАРИАЦИОННЫЕ МЕТОДЫ РОЗВ'ЯЗКИ УРАВНЕНИЙ З ЧАСТИННИМИ ПОХІДНИМИ МЕТА
ЛЕКЦІЇ: Показать, как можна решать дифференциальное уравнение,
рассматривая его как уравнение Эйлера---Лагранжа для некоторого
функционала и находя некоторым новым методом функцию, минимизирующую цей
функционал. Тогда минимизирующая функція буде решением уравнения з
частинними похідними. Задача, конечно, состоит в том, чтобы найти такой
функционал, для которого исходное уравнение является уравнением
Эйлера---Лагранжа. Хорошо известен результат (теорема о минимуме
энергии), який гласит, що нахождение розв'язки и некоторой эллиптической
краевой задачи

, в области D , u = 0 , на границе D ,

эквивалентно нахождению такой функції u (тоже обращающейся в нуль на
границе области D), яка минимизирует функционал потенинальной энергии

Ми можемо розв'язати задачу наступними трьома кроками:

ШАГ І (Определение подходящего преобразования.)

Оскільки змінна x змінюється від 0 до 1, ми використаємо фінальне
перетворення, а саме синусоїдальне перетворення; Пізніше стане
зрозуміло, чому. Ми могли б розв'язати цю задачу, використовуючи
перетворення Лапласа на змінній t (цей метод не відрізняється від
скінченного синусоїдального перетворення за інтенсивністю праці).

КРОК 2 (Виконати конвертацію).

Для удобства будем использовать обозначение
\texttt{S\_n(t)\ =\ S{[}u{]}} . Применим синус-преобразование к
исходному уравнению

\begin{verbatim}
То есть,  \Delta u = f  —уравнение Эйлера—Лагранжа для J[u]. Функція, приближенно минимизирующая функционал J[u], ищется методом Ритца; тем самым мы получаем розв'язок (приближениое) уравнения з частинними похідними. В цієї лекции мы познакомимся с методом Ритца и посмотрим, как с помощью цього метода можна и и и и и и и уравнения.
\end{verbatim}

Очень хорошо решать граничные задачи (как, наприклад, задачу о
закрепленной мембране) путем отыскания гладкой поверхности, яка
минимизирует потенциальную энергию мембраны. То есть, якщо мы
рассматриваем исходное уравнение как уравнение Эйлера --- Лагранжа для
искоторого функционала J{[}u{]}, то розв'язок дифференциального
уравнения можна получить путем минимизации функционала (поскольку
минимизирующая функція функционала является решением соответствующего
уравнения

Эйлера --- Лагранжа). В лекции 44 мы рассматривали такие функционалы, у
которых уравнение Эйлера --- Лагранжа було обыкновенным дифференциальным
уравнением. В цієї лекции рассматриваются функционалы, для которых
уравнение Эйлера --- Лагранжа относится к классу уравнений з частинними
похідними. Наприклад для функционала

и воспользуемся тождеством для синус-преобразования, получим

уравнение Эйлера---Лагранжа (подобно случаю обыкновенного
дифференциального уравнения в последней лекции) имеет вид

и, отже, для розв'язки задачи Дирихле в единичном квадрате

(УЧП)

, 0 \textless{} x \textless{} 1 , 0 \textless{} y \textless{} 1 , (ГУ) u
= g , на границе квадрата,

можна, наоборот, найти такую функцию \texttt{u(x,\ y)}, яка минимизирует
J{[}u{]} и равна g на границе. Вероятно, не стоит удивляться тому, що
функционал

представляет собой потенциальную эпергию мембраны и фактически мы
находим поверхность с минимальной потенциальной энерешей. Вопрос,
конечно, состоит в следующем: якщо дано дифференциальное уравнение, как
мы находим функционал, выражающий потенциальную энергию розв'язки? Ответ
на цей вопрос дает хорошо известная теорема (теорема о минимуме
энергии), яка утверждает:

Розв'язок и задачи Дирихле

где

\texttt{\^{}\{\textbackslash{}Lambda\}u\ =\ f} в области
\texttt{\^{}\{D\}} \texttt{(ГУ)} \texttt{\^{}\{U\}=0} на границе
\texttt{\^{}\{D\}} является той же функциен, яка минимизирует (среди
функцій, удовлетворяющих граничному условию u=0) энергегический
функционал

\begin{verbatim}
\begin{cases} u_{tx}, \quad 0 = \begin{cases} 1, \quad -1 < x < 1, \\ 0 \quad \text{в остальных точках,} \end{cases} \end{cases}
\end{verbatim}

и, отже, для розв'язки задачи (45.1) достаточно найти функцию
\texttt{\textbackslash{}overline\{u\}} (среди функцій, обращающихся в
нуль на граннце), минимизирующую
\texttt{J\textbackslash{}\{u\textbackslash{}\}} . Этим замечанием
завершается первая часть лекции. В оставшейся части мы покажем, как
находить минимизирующую функцию \texttt{\textbackslash{}overline\{u\}}
для J{[}u{]}, используя метод Ритца.

\section{Метод Ритца минимизации
функционалов}\label{ux43cux435ux442ux43eux434-ux440ux438ux442ux446ux430-ux43cux438ux43dux438ux43cux438ux437ux430ux446ux438ux438-ux444ux443ux43dux43aux446ux438ux43eux43dux430ux43bux43eux432}

Цей метод --- один из многих вариационных методов, описанных в
литературе по вариационному исчислению. Идея, предложенная математиком
В. Ритцем, очень проста. Метод состоит из следующих шагов.

ШАГ 1. Выбираем n и функцию u, минимизирующую функционал

Якщо тепер перетворити початкові умови задачі на граничне значення,
отримаємо початкові умови для звичайного диференціального рівняння

ищем в виде

\begin{verbatim}
где функції  \phi_1, \phi_2, \phi_3, \ldots, \phi_n  принадлежат к классу достаточно хороших функцій и все обращаются в нуль на границе, так що из них можна построить разумное приближение для розв'язки задачи. Эти функції принято называть пробными функциями. Типичный набор пробных функцій для задачи Дирихле в единичном квадрате приведен ниже:
\end{verbatim}

\texttt{—\ обращаются\ в\ нуль\ на\ границе,\ \textquotesingle{}\textbackslash{}varphi\_2(x,\ y)\ =\ x\textbackslash{}varphi\_1(x,\ y),\textquotesingle{}\ \textquotesingle{}\textbackslash{}varphi\_3(x,\ y)\ =\ y\textbackslash{}varphi\_1(x,\ y),\textquotesingle{}\ \textquotesingle{}\textbackslash{}varphi\_4(x,\ y)\ =\ x\^{}2\textbackslash{}varphi\_1(x,\ y),\textquotesingle{}\ \textquotesingle{}\textbackslash{}varphi\_5(x,\ y)\ =\ xy\textbackslash{}varphi\_1(x,\ y),\textquotesingle{}\ \textquotesingle{}\textbackslash{}varphi\_6(x,\ y)\ =\ y\^{}2\textbackslash{}varphi\_1(x,\ y),\textquotesingle{}\ \textquotesingle{}\textbackslash{}vdots\textquotesingle{}\ \textquotesingle{}\textbackslash{}vdots\textquotesingle{}\ Иными\ словами,\ первые\ четыре\ приближения\ имеют\ вид}

\begin{verbatim}
\begin{aligned} & u_{ti} = c^{2}u_{xx}, & -\infty < x < +\infty, & 0 < t < \infty, \\ & (\text{НУ}) & \begin{cases} u_{t}(x, 0) = 0, \\ u_{t}(x, 0) = \begin{cases} 1, & -1 < x < 1, \\ 0, & \text{в остальных случаях.} \end{cases} \end{aligned} \end{cases}
\end{verbatim}

в координатах (x, t).

\pandocbounded{\includegraphics[keepaspectratio]{images/_page_134_Figure_7.jpeg}}

Рис. 18.4. Розв'язок задачи (18.4) в xt-плоскости.

Задача (18.4) описывает поведение струны, которой сообщается рачальная
единичная скорость на отрезке -1 \textless{} x \textless{} 1. Смещение .
.

описывается формулой Даламбера

Итак, решим теперь новое семейство задач Кощи:

(ОДУ)

\texttt{является\ теперь\ функцией\ коефіцієнтів\ \textquotesingle{}a\_1a\_2\textbackslash{}dots\ a\_n\textquotesingle{}\ .\ Отже,\ для\ того\ чтобы\ найти\ минимум\ функционала\ J,\ приравняем\ нулю\ частные\ производные}

\begin{verbatim}
Aa = b
b_i = -\int_0^1 \int_0^1 f(x, y) \varphi_i(x, y) dx dy,
\end{verbatim}

Це розв'язок в различные моменты времени изображено на рис. 18.5.

\pandocbounded{\includegraphics[keepaspectratio]{images/_page_135_Figure_5.jpeg}}

Рис. 18.5. Розв'язок задачи (18.4) в различные моменты времени.

\begin{verbatim}
Этим завершается наша интерпретация формулы Даламбера в плоскости змінних  \boldsymbol{x}  и t. В оставшейся части лекции мы займемся решением смешанной задачи для полубесконечной струны:
\end{verbatim}

(18.6)

Единственное отличие в том, що аппроксимирующие функції
\texttt{\textbackslash{}phi\_1,\ \textbackslash{}phi\_2,\ \textbackslash{}phi\_3,\ \textbackslash{}ldots,\ \textbackslash{}phi\_n}
повинні удовлетворять граничным условиям:

\begin{verbatim}
\right.
\end{verbatim}

является иравнение Лапласа, но це доказательство проводится аналогично
тому, как це делалось в лекции 44 для функционала, зависящего от функції
одной переменной.

\begin{enumerate}
\def\labelenumi{\arabic{enumi}.}
\setcounter{enumi}{2}
\item
  В книгах по вариационному исчислению показано, как строить
  энергетический функционал не лише для уравнения Пуассона с граничным
  условием Дирихле u=0, но и для многих других дифференциальных
  уравнений и многих других типов граничных условий. Отже, большое число
  краевых задач можна решить, минимизируя соответствующие энергетические
  функционалы.
\item
  Чем большее значение п мы выбираем, тем меньшее значение
  \texttt{J{[}\textbackslash{}bar\{u\}\_n{]}} получается. Значит, с
  увеличением n растет точность л приближенного розв'язки
  дифференциального уравнения. Один из методов определения и, основан на
  последовательном вычислении функционала \texttt{J{[}u\_n{]}} для
  возрастающих значений n.~Обчислення прекращаются, якщо с ростом п
  функционал перестает практически уменьшаться.
\item
  Якщо и велико, то проведение вычислений по методу Ритца требует
  привлечения ЭВМ. На рис. 45.2 (с. 345) приведена блок-
\end{enumerate}

схема розв'язки краевой задачи методом Ритца.

Пользователь повинен позаботиться о подпрограмме, яка должна вычислять
значения пробных функцій. Такая подпрограмма может иметь вид:

SUBROUTINE BC (X.Y.PHI)

\begin{itemize}
\tightlist
\item
  SUBROUTINE PROVIDED BY THE USER TO C - EVALUATE THE - FUNCTIONS
  PHI(1), PHI(2), \ldots, PHI(N) DIMENSION PHI(20)
\end{itemize}

PHI(1) = X*Y*(1-X)*(1-Y) \texttt{PHI(2)\ =\ X\^{}*PHI(1)} (используемые
функції)

\begin{verbatim}
PHI(3) = Y\*PHI(1)
\end{verbatim}

PHI(N) = (що получится)RETURN END

\section{\texorpdfstring{\textbf{ЗА}ДАЧИ}{ЗАДАЧИ}}\label{ux437ux430ux434ux430ux447ux438-2}

Найдите энергетический функционал для задачи

(УЧП)

\pandocbounded{\includegraphics[keepaspectratio]{images/_page_362_Figure_1.jpeg}}

, \texttt{0\ \textless{}\ x\ \textless{}\ 1} ,
\texttt{0\ \textless{}\ y\ \textless{}\ 1} , (ГУ) \texttt{u\ =\ 0} на
границе.

\begin{enumerate}
\def\labelenumi{\arabic{enumi}.}
\setcounter{enumi}{1}
\tightlist
\item
  Как минимизировать функционал
\end{enumerate}

методом Ритца?

УКАЗАНИЕ. Якщо ввести новую функцию \texttt{z(x)}:

то можна заметить, що она удовлетворяет граничным условиям \texttt{z(0)}
= 0 и \texttt{z(1)} = 0.1

\begin{itemize}
\tightlist
\item
  \begin{enumerate}
  \def\labelenumi{\arabic{enumi}.}
  \setcounter{enumi}{2}
  \tightlist
  \item
    Напишите программу для проведения вычислений по блоксхеме,
    изображенной на рис. 45 2.
  \end{enumerate}
\item
  \begin{enumerate}
  \def\labelenumi{\arabic{enumi}.}
  \setcounter{enumi}{3}
  \tightlist
  \item
    Покажите, що уравнение я в
  \end{enumerate}
\end{itemize}

\pandocbounded{\includegraphics[keepaspectratio]{images/_page_362_Figure_2.jpeg}}

ляется уравнением Эйлера --- Лагранжа для функционала

\pandocbounded{\includegraphics[keepaspectratio]{images/_page_362_Figure_5.jpeg}}

\begin{enumerate}
\def\labelenumi{\arabic{enumi}.}
\setcounter{enumi}{4}
\tightlist
\item
  Задачу Дирихле г д
\end{enumerate}

\pandocbounded{\includegraphics[keepaspectratio]{images/_page_362_Figure_7.jpeg}}

\texttt{е\ A=0.06\ и\ B=0.04.\ Как\ найти\ потенциальную\ энергию\ цього\ розв\textquotesingle{}язки?\ Мы\ рекомендуем\ читателю\ вспомнить\ синуспреобразование\ и\ самостоятельно\ получить\ розв\textquotesingle{}язок\ \textquotesingle{}u\textbackslash{}left(x,y\textbackslash{}right)\textquotesingle{}\ .}

*) При такой замене возникает особенность в т. x=1. Лучше ввести функцию
\texttt{z(x)} по формуле \texttt{z(x)}=\texttt{y(x)}=x. Прим. ред.

\pandocbounded{\includegraphics[keepaspectratio]{images/_page_343_Figure_2.jpeg}}

Далсе. Приближенное розв'язок
\texttt{u\_n(x,y)\ =\ a\_1\ \textbackslash{}varphi\_1\ +\ a\_2\ \textbackslash{}varphi\_2\ +\ \textbackslash{}dots\ +\ a\_n\ \textbackslash{}varphi\_n}
. В це время можна численно найти эначение функционала
\texttt{J{[}u\_n{]}} . Якщо необходимо найти значение розв'язки
\texttt{u\_n(x,y)} в некоторой точке, можна написать небольшую программу
эля быполнения отой операции. При необходимости можна повторить
обчислення по цієї программе с другими значениями п Рис. 452. Блок-схема
метода Ритца.

\section{РОЗВ'ЯЗОК УРАВНЕНИЙ З ЧАСТИННИМИ ПОХІДНИМИ МЕТОДАМИ ТЕОРИИ
ВОЗМУЩЕНИЙ МЕТА ЛЕКЦІЇ: Показать, как различные сложные задачи
(нелинейные уравнения с змінними коэффициентами, области неправильной
формы и т. д.) можна решить методом возмущения более простых задач.
Иными словами, мы хотим показать, как следует изменить розв'язок простой
задачи, чтобы оно давало приближенное розв'язок трудной
задачи.}\label{ux440ux43eux437ux432ux44fux437ux43eux43a-ux443ux440ux430ux432ux43dux435ux43dux438ux439-ux437-ux447ux430ux441ux442ux438ux43dux43dux438ux43cux438-ux43fux43eux445ux456ux434ux43dux438ux43cux438-ux43cux435ux442ux43eux434ux430ux43cux438-ux442ux435ux43eux440ux438ux438-ux432ux43eux437ux43cux443ux449ux435ux43dux438ux439-ux43cux435ux442ux430-ux43bux435ux43aux446ux456ux457-ux43fux43eux43aux430ux437ux430ux442ux44c-ux43aux430ux43a-ux440ux430ux437ux43bux438ux447ux43dux44bux435-ux441ux43bux43eux436ux43dux44bux435-ux437ux430ux434ux430ux447ux438-ux43dux435ux43bux438ux43dux435ux439ux43dux44bux435-ux443ux440ux430ux432ux43dux435ux43dux438ux44f-ux441-ux437ux43cux456ux43dux43dux438ux43cux438-ux43aux43eux44dux444ux444ux438ux446ux438ux435ux43dux442ux430ux43cux438-ux43eux431ux43bux430ux441ux442ux438-ux43dux435ux43fux440ux430ux432ux438ux43bux44cux43dux43eux439-ux444ux43eux440ux43cux44b-ux438-ux442.-ux434.-ux43cux43eux436ux43dux430-ux440ux435ux448ux438ux442ux44c-ux43cux435ux442ux43eux434ux43eux43c-ux432ux43eux437ux43cux443ux449ux435ux43dux438ux44f-ux431ux43eux43bux435ux435-ux43fux440ux43eux441ux442ux44bux445-ux437ux430ux434ux430ux447.-ux438ux43dux44bux43cux438-ux441ux43bux43eux432ux430ux43cux438-ux43cux44b-ux445ux43eux442ux438ux43c-ux43fux43eux43aux430ux437ux430ux442ux44c-ux43aux430ux43a-ux441ux43bux435ux434ux443ux435ux442-ux438ux437ux43cux435ux43dux438ux442ux44c-ux440ux43eux437ux432ux44fux437ux43eux43a-ux43fux440ux43eux441ux442ux43eux439-ux437ux430ux434ux430ux447ux438-ux447ux442ux43eux431ux44b-ux43eux43dux43e-ux434ux430ux432ux430ux43bux43e-ux43fux440ux438ux431ux43bux438ux436ux435ux43dux43dux43eux435-ux440ux43eux437ux432ux44fux437ux43eux43a-ux442ux440ux443ux434ux43dux43eux439-ux437ux430ux434ux430ux447ux438.}

Наприклад, мы покажем, как нелинейную задачу

(УЧП)

\begin{verbatim}
,  0 < r < 1  ,
\end{verbatim}

\begin{verbatim}
(ГУ)  u(1, \theta) = \cos \theta  ,  0 \le \theta \le 2\pi  .
\end{verbatim}

Очень часто задачу, розв'язок которой не известно, можна путем
непрерывного изменения каких-то параметров свести к близкой, но легко
решаемой задаче. Тогда и розв'язок исходной задачи буде близко решению
модифицирующей задачи.

Наприклад, рівняння Лапласа

в помощью семейства уравнений

\begin{verbatim}
H(x-a) =
\end{verbatim}

\texttt{,\ \textquotesingle{}0\ \textbackslash{}le\ \textbackslash{}varepsilon\ \textbackslash{}le\ 1\textquotesingle{}\ ,}

\begin{verbatim}
\begin{cases} 0, x < a, \\ 1, x \geqslant a, \end{cases}
\end{verbatim}

\texttt{,\ \textquotesingle{}0\ \textless{}\ r\ \textless{}\ 1\textquotesingle{}\ ,\ \textquotesingle{}0\ \textbackslash{}le\ \textbackslash{}theta\ \textless{}\ 2\textbackslash{}pi\textquotesingle{}\ ,\ \textquotesingle{}u(1,\ \textbackslash{}theta)\ =\ g(\textbackslash{}theta)\textquotesingle{}\ ,\ \textquotesingle{}0\ \textbackslash{}le\ \textbackslash{}theta\ \textless{}\ 2\textbackslash{}pi\textquotesingle{}\ ,}

при одних граничных условиях (наприклад,
\texttt{g(\textbackslash{}theta)\ =\ \textbackslash{}cos\ \textbackslash{}theta}
) имеет два различных вещественных розв'язки, а при других (наприклад,
\texttt{g(\textbackslash{}theta)\ =\ A\ \textbackslash{}cos\ \textbackslash{}theta}
с A \textgreater{} 20,65) ця задача не имеет ни одного вещественного
розв'язки. --- Прим. ред.

4) Краевая задача

\texttt{можна\ непрерывно\ модифицировать\ в\ нелинейное\ уравнение\ \textquotesingle{}\textbackslash{}Delta\ u\ =\ u\^{}2\ =\ 0\textquotesingle{}\ 1).}

На рис. 46.1 изображена схема розв'язки нелинейной задачи путем сведения
ее к решению уравнения Лапласа2)

В результате получаем

Рис. 46.1. Схема метода возмущении.

Чтобы найти розв'язок возмущенного уравнения Лапласа

\begin{verbatim}
H(x-a) =
u(0, t) = f(t),
\begin{cases} 0, & x < a, \\ 1, & x \geqslant a, \end{cases}
u_{x}(0, t) = 0.
\begin{cases} 0, & n = 2, 4, \dots, \\ (4/n\pi)\cos(n\pi t), & n = 3, 5, 7, \dots \end{cases}
x + ct = \text{const},
f(x) = \mathcal{F}^{-1}[F] = \frac{1}{\sqrt{2\pi}} \int_{-\infty}^{+\infty} F(\omega) e^{-i\omega x} d\omega \qquad F(\omega) = \mathcal{F}[f] = \frac{1}{\sqrt{2\pi}} \int_{-\infty}^{+\infty} f(x) e^{-i\omega x} dx
\end{verbatim}

\texttt{\textquotesingle{}0\ \textless{}\ r\ \textless{}\ 1,\textquotesingle{}\ \textquotesingle{}0\ \textbackslash{}le\ \textbackslash{}theta\ \textless{}\ 2\textbackslash{}pi,\textquotesingle{}\ \textquotesingle{}u(1,\ 0)\ =\ 0,\textquotesingle{}\ \textquotesingle{}0\ \textbackslash{}le\ \textbackslash{}theta\ \textless{}\ 2\textbackslash{}pi,\textquotesingle{}\ нмеет\ два\ розв\textquotesingle{}язки\ (одно\ —\ нулсьое,\ другое\ —\ нетривиальное),\ в\ то\ время\ как\ задача}

Следовательно, решение нашей задачи можно представить в виде

\begin{verbatim}
u(1, \theta) = 0, \quad 0 \le \theta < 2\pi,
\begin{cases} 1, |x| < a \\ 0, |x| > a \end{cases}
\end{verbatim}

\begin{verbatim}
\begin{array}{ccccccccccccccccccccccccccccccccccc
\end{verbatim}

\section{ЗАМЕЧ
АНИЯ}\label{ux437ux430ux43cux435ux447-ux430ux43dux438ux44f-1}

\begin{enumerate}
\def\labelenumi{\arabic{enumi}.}
\tightlist
\item
  Для решения задачи методом конечного синус- или косинуспреобразования
  краевые условия при x=0 и x=L должны иметь вид
\end{enumerate}

\begin{verbatim}
\begin{array}{ccc} (\text{УЧ}\Pi) & \Delta u = 0, & 0 < r < 1, \\ (\text{ГУ}) & u(1, \theta) = \cos \theta, & 0 \leqslant \theta \leqslant 2\pi, \end{array}
\begin{cases} 1, |x| < 1 \\ 0, |x| > 1 \end{cases}
u_{tt} = c^2 u_{xx}, \quad 0 < x < L, \quad 0 < t < \infty,
u(0, t) = g_1(t),
u_x(0, t) = g_1(t),
\delta(x-a)
u_x(0, t) - \gamma_1 u(0, t) = g_1(t),
\end{verbatim}

\begin{verbatim}
\begin{array}{ll} (\mathrm{Y} \mathrm{U} \Pi) & u_t = u_{xx}, \quad 0 < x < 1, \quad 0 < t < \infty, \\ (\mathrm{\Gamma} \mathrm{Y} & \left\{ \begin{array}{ll} u_x \left( 0, \ t \right) = 0 \\ u_x \left( 1, \ t \right) = 0, \end{array} \end{array}
\end{verbatim}

\begin{verbatim}
u(0, t) = g_1(t),
(1+x^2)^{-1}
u_x(0, t) = g_1(t),
xe^{-a+x^4}, a > 0
u_x(0, t) - \gamma_1 u(0, t) = g_1(t),
\end{verbatim}

\texttt{H(x+a)} - \texttt{H(x-a)}

\begin{verbatim}
\begin{cases} \Delta u_{1} = -r^{2} \cos^{2} \theta = -\frac{r^{2}}{2} \left[1 + \cos(2\theta)\right] = -\frac{r^{2}}{2} - \frac{r^{2}}{2} \cos(2\theta), \\ u_{1}(1, \theta) = 0. \end{cases}
\frac{a}{x^2+a^2}
\Delta u = -\frac{r^2}{2} - \frac{r^3}{2} \cos{(2\theta)}
\end{verbatim}

\begin{verbatim}
\right. & 0 < t < \infty, \\ (\text{НУ}) & u(x, 0) = 0, \quad 0 \leqslant x \leqslant 1. \end{array}
\end{verbatim}

\begin{verbatim}
u_p(r, \theta) = Ar^4 + Br^4 \cos(2\theta).
\frac{2ax}{(x^2+a^2)^2}
u_p(r, \theta) = -\frac{r^4}{32} - \frac{r^4}{24} \cos(2\theta).
\begin{cases} \cos(ax), |x| > \pi/2a \\ 0, |x| > \pi/2a \end{cases}
\sum_{n=0}^{\infty} [a_n \cos(n\theta) + b_n \sin(n\theta)] - \frac{1}{32} - \frac{1}{24} \cos(2\theta) = 0.
\end{verbatim}

\pandocbounded{\includegraphics[keepaspectratio]{images/_page_189_Picture_10.jpeg}}

Отметим, що якщо величина \texttt{u(0,\ t)} положительна, то величина
\texttt{u\_x(0,\ t)} також положительна, якщо \texttt{u(L,\ t)}
положительна, то \texttt{u\_x(L,\ t)}

\pandocbounded{\includegraphics[keepaspectratio]{images/_page_144_Picture_4.jpeg}}

Рис. 19.6. Возникновение неоднородных граничных условий упругого
закрепления: h---коефіцієнт упругости пружины.

отрицательна. Однородные граничні умови (19.2) можна записать в виде

(19.3)

\begin{enumerate}
\def\labelenumi{\arabic{enumi}.}
\setcounter{enumi}{17}
\tightlist
\item
\end{enumerate}

\begin{verbatim}
P_{2}
\begin{cases} 1-|x|, |x| < 1 \\ 0, |x| > 1 \end{cases}
\Delta u = 0, \quad 0 < r < 1 + \frac{1}{4} \sin \theta,
\end{verbatim}

досягається в кінцевій точці сімейства граничних умов виду (див. рисунок
46 3)

\begin{verbatim}
Tu_x(0, t) = -\beta u_t(0, t).
\cos(ax)
Tu_{x}(0, t) = \varphi[u(0, t)],
\frac{2}{\pi} \frac{\sin a\omega}{\omega}
Tu_{x}(0, t) = -hu^{3}(0, t)
\end{verbatim}

Перед нами неоднородное уравнение теплопроводности, так что метод
разделения переменных неприменим. Можпо, конечно, воспользоваться
конечным синус-преобразованием по переменной x или преобразованием
Лапласа по переменной t, но вместо этого мы рассмотрим две подзадачи

\begin{verbatim}
mu_{tt}(L, t) = -ku_x(L, t) + mg.
\frac{2}{\pi} \frac{\sin a\omega}{(\omega^2+a^2)^2}
P_{0}
\frac{\pi}{2} e^{-a+\omega}
\end{verbatim}

Отже, розв'язуючи кожну з задач Дірікла всередині кола для функцій
\texttt{u\_0} , \texttt{u\_1} , \texttt{u\_2} , отримуємо розв'язок

(46.9)

Задачі Діріхле в деформованій області (зверніть увагу, що тут

\begin{verbatim}
Ми вже взяли  \epsilon = 1/4  :)
\end{verbatim}

Як вправу, читача просять знайти пертурбацію першого порядку
\texttt{u\_1} і перевірити, наскільки добре є приблизне розв'язання

якщо розглянути рівняння з параметром

(46.11)

\begin{verbatim}
\frac{\pi}{2} e^{-a+\omega}
\end{verbatim}

і побачити його розв'язок у вигляді

\begin{verbatim}
\frac{\pi}{2} e^{-a+\omega}
\end{verbatim}

Підставляючи цю декомпозицію на (46.11), отримуємо таку послідовність
завдань

\begin{verbatim}
\frac{\pi}{2} e^{-a+\omega}
u_{x}(0, t) = \frac{h}{T} [u(0, t) - \theta_{1}(t)]
\frac{\pi}{2} e^{-a+\omega}
u_{tt} = u_{xx}, \quad 0 < x < 1, \quad 0 < t < \infty,
\end{verbatim}

\begin{verbatim}
Каждая из этих задач легко решается и в то же время ясно, что  \mathit{suma}  решений задач  P_1  и  P_2  является решением исходной вадачи, т. е.
\end{verbatim}

\begin{verbatim}
\Delta u = 0,
\frac{2ax}{(x^2+a^2)^2}
q_{uu} + q_{vv} = 0
\end{verbatim}

Решение задачи \texttt{P\_\{i\}} Розв'язання задачі \texttt{P\_i}
Розв'язання задачі \texttt{P\_i} Можна показати, що в загальному випадку
розв'язок задачі

\begin{verbatim}
w = f(z)
\frac{\cos(ax), |x| > \pi/2a}{0, |x| > \pi/2a}
\end{verbatim}

\begin{verbatim}
\left\{
\end{verbatim}

\begin{verbatim}
\frac{1-|x|, |x| < 1}{0, |x| > 1}
\Rightarrow.
\begin{aligned} (Y & \Pi) & \varphi_{xx} + \varphi_{yy} = 0, & -\infty < x < \infty, & 0 < y < \infty, \\ (47.1) & \varphi(x, 0) = \begin{cases} 0, & |x| > 1, \\ 1, & |x| \le 1. \end{cases} \end{aligned}
w = \ln\left\{\frac{z-1}{z+1}\right\}
\cos(ax)
\end{verbatim}

\begin{verbatim}
\begin{cases} (\text{УЧП}) & \varphi_{uu} + \varphi_{vv} = 0, \quad -\infty < u < \infty, \quad 0 < v < \pi, \\ (\text{47.2}) & \begin{cases} \varphi(u, 0) = 0, \\ \varphi(u, \pi) = 1. \end{cases}
\end{verbatim}

\begin{verbatim}
\frac{\pi}{2} [\delta(\omega+a) + \delta(\omega-a)]
u(x, t) = \sum_{n=1}^{\infty} c_n X_n(x) T_n(t),
u_{tt} = \alpha^2 u_{xx}, \quad 0 < x < L, \quad 0 < t < \infty,
(\Gamma Y)
f(x) = \int_{0}^{\infty} F(\omega) \sin(\omega x) d\omega
\end{verbatim}

Щоб візуально зрозуміти, як працює ця функція, ми радимо читачеві
побудувати її за кількома лініями y=c.

У другому прикладі ми перетворюємо область між двома неконцентричними
колами у кільце.

Задача Діріхле в області між двома некопоцентричними колами

Припустимо, ви хочете знайти потенціал між двома колами

\begin{verbatim}
F(\omega) = \frac{?}{\pi} \int_{0}^{\infty} f(x) \sin(\omega x) dx
\quad 0 < t < \infty,
0 < x < \infty
\end{verbatim}

(HY)

\begin{verbatim}
u = \gamma [(x-s)(x-t)-\gamma^2],
\quad 0 \le x \le L,
0 < \omega < x
u(x, t) = X(x) T(t).
-\omega^2 F(\omega) + \frac{2}{\pi} \omega f(0)
T'' - \alpha^2 \lambda T = 0,
\frac{1}{\alpha} F\left(\frac{\omega}{\alpha}\right)
X'' - \lambda X = 0,
\end{verbatim}

проследим, как оно решается методом конечного синус-преобразования. Вот
что мы должны сделать: разложить еходное воздействие \texttt{f(x,\ t)}
на компоненты, найти отклик \texttt{U\_n} на каждую компоненту, а затем
сложить эти отклики. Математически это может быть и не совсем очевидно,
но давайте внимательно проследим за ходом решения. Разложим входящие в
уравнение с частными производными члены в ряды Фурье по синусам:

\begin{verbatim}
u(x, t) = [C \sin(\beta x) + D \cos(\beta x)][A \sin(\alpha \beta t) + B \cos(\alpha \beta t)]
\frac{2\omega}{\pi (\alpha^2 + \omega^2)}
\beta_n = \frac{n\pi}{L}, \quad n = 0, 1, 2, \ldots
x^{-1/2}
u_n(x, t) = X_n(x) T_n(t) = \sin(n\pi x/L) [a_n \sin(n\pi at/L) + b_n \cos(n\pi at/L)],
[2/\pi\omega]^{1/2}
u_n(x, t) = R_n \sin\left(\frac{n\pi x}{L}\right) \cos\left(\frac{n\pi\alpha (t - \delta_n)}{L}\right).
\end{verbatim}

where

\begin{verbatim}
H(x-a) =
\frac{2}{\pi\omega} [1-\cos(\omega a)]
\end{verbatim}

Функція Хевісайду,

1

\begin{verbatim}
\begin{cases} 0, & x < a, \\ 1, & x \geqslant a, \end{cases}
e^{-1}a + \omega
f(x) = \mathcal{F}^{-1}[F] = \frac{1}{\sqrt{2\pi}} \int_{-\infty}^{+\infty} F(\omega) e^{-i\omega x} d\omega \qquad F(\omega) = \mathcal{F}[f] = \frac{1}{\sqrt{2\pi}} \int_{-\infty}^{+\infty} f(x) e^{-i\omega x} dx
\end{verbatim}

Зверніть увагу, по-перше, що коефіцієнти \texttt{A\_n} , \texttt{B\_n}
\texttt{F\_n} дійсно є функціями часу t, оскільки ми застосували
синусоїдальне перетворення до функції двох змінних x, t, u, а по-друге,
ми розклали вхідний ефект \texttt{f(x,t)} на прості компоненти
\texttt{F\_n}(t) . Тепер знайдіть відповіді \texttt{U\_n}(t) для кожного
компонента \texttt{F\_n}(t) , потім додайте всі відповіді
\texttt{U\_n}(t) і отримайте розв'язок \texttt{u(x,t)}.

Щоб знайти відповіді \texttt{U\_n}(t) , необхідно робити невеликі
розрахунки з коефіцієнтами \texttt{A\_n}(t) і \texttt{B\_n}(t) , щоб у
вирази інтегрування була включена лише функція u, а не її похідні
\texttt{u\_i} і \texttt{u\_\{xx\}} . Інтегруючи на частини, ми отримуємо

\begin{verbatim}
\begin{cases} 1, |x| < a \\ 0, |x| > a \end{cases}
\frac{1}{2}e^{-\alpha}\sin\omega
\delta(x-a)
\frac{1-e^{-a\omega}}{\omega}
a_{n} = \frac{2}{n\pi\alpha} \int_{0}^{L} g(x) \sin(n\pi x/L) dx, \\ b_{n} = \frac{2}{L} \int_{0}^{L} f(x) \sin(n\pi x/L) dx.
u(x, t) = \sum_{n=1}^{\infty} \sin(n\pi x/L) \left[ a_n \sin(n\pi \alpha t/L) + b_n \cos(n\pi \alpha t/L) \right],
u(x, t) = \sum_{n=1}^{\infty} b_n \sin(n\pi x/L) \cos(n\pi\alpha t/L)
\frac{2}{\pi}F''(\omega)
u(x, \theta) = f(x)
f(x) = \int_{0}^{\infty} F(\omega) \cos(\omega x) d\omega \qquad \qquad F(x) = \frac{2}{\pi} \int_{0}^{\infty} f(x) \cos(\omega x) dx
f(x) = \sum_{n=1}^{\infty} b_n \sin(n\pi x/L).
0 < x < \infty \qquad \qquad 0 < \omega < \infty
u_n(x, t) = b_n \sin(n\pi x/L) \cos(n\pi\alpha t/L).
u(x, t) = \sin(\pi x/L)\cos(\pi \alpha t/L) + 0.5\sin(3\pi x/L)\cos(3\pi \alpha t/L) + 0.25\sin(5\pi x/L)\cos(5\pi \alpha t/L).
R_n \sin(n\pi x/L) \cos[n\pi\alpha(t-\delta_n)/L],
\end{verbatim}

j'\,'(x)

\begin{verbatim}
\omega_n = \frac{n\pi\alpha}{L} = \frac{n\pi}{L} \sqrt{\frac{T}{\rho}},
S_n = \frac{2}{\pi} \int_0^{\pi} f(x) \sin(nx) dx.
u(x, 0) = \sin(\pi x/L) + 0.5 \sin(3\pi x/L),
\end{verbatim}

в последнее соотношение, получаем

\begin{verbatim}
u(x, 0) = 0,
y = \frac{\pi}{b - a} (x - a).
u(x, 0) = \sin(3\pi x/L),
\end{verbatim}

и, следовательно, (26.1) принимает вид

\begin{verbatim}
u(x, 0) =
f(x) = \sum_{n=1}^{\infty} S_n \sin(nx).
\begin{cases} 2hx, & 0 \le x \le 0.5, \\ 2h(1-x), & 0.5 \le x \le 1. \end{cases}
\end{verbatim}

\begin{verbatim}
Поскольку это разложение является тождеством по x, его коэффициенты должны быть равны нулю,  \tau  . е.
\end{verbatim}

К а к

\pandocbounded{\includegraphics[keepaspectratio]{images/_page_153_Picture_8.jpeg}}

ие движения буде совершать струна, якщо ее отпустить? 6. Решите задачу о
затухающих колебаниях струны

\begin{verbatim}
f(x) = \sum_{n=1}^{\infty} S_n \sin(nx)
\frac{2ax}{(x^2+a^2)^2}
\end{verbatim}

.

Таким образом, мы получили связь между входом и выходом системы, между
коэффициентами \texttt{F\_n} и \texttt{U\_n} . Прежде чем решать
уравнения для \texttt{U\_n}(t) , обратимся к начальному условию

Представляется ли вам разумным полученное розв'язок? Удовлетворяет ли
оно уравнению з частинними похідними, граничным и начальным условиям?

\begin{enumerate}
\def\labelenumi{\arabic{enumi}.}
\setcounter{enumi}{6}
\tightlist
\item
  Как вы будете решать неоднородное уравнение с заданными начальными и
  граничными условиям?
\end{enumerate}

\begin{verbatim}
S_n = \frac{2}{\pi} \int_{0}^{\pi} f(x) \sin(nx) dx
\cos(ax)
\end{verbatim}

Если разложить \texttt{u(x,\ 0)} в ряд Фурье по синусам и приравнять его
нулю:

\begin{verbatim}
\frac{\pi}{2} [\delta(\omega+a) + \delta(\omega-a)]
\end{verbatim}

\begin{verbatim}
то определяются начальные условия для функции  U_n\left(t\right)  , т. е.
\end{verbatim}

\begin{verbatim}
f(x) = \int_{0}^{\infty} F(\omega) \sin(\omega x) d\omega
n = 1, 2, ...
0 < x < \infty
\end{verbatim}

\begin{verbatim}
,  n = 1,  2, \dots  Мы получили разложение нашей смешанной задачи на ряд простых задач об отклике системы на элементарные входные воздействия:
\end{verbatim}

(OДУ)

\begin{enumerate}
\def\labelenumi{\arabic{enumi}.}
\tightlist
\item
\end{enumerate}

(HУ) \texttt{U\_n}(0) = 0, n = 1, 2, \ldots{} Каждое из этих уравнений
легко решается методом интегрирующего множителя (или с помощью
преобразования Лапласа). В любом случае получаем

\begin{verbatim}
\frac{1}{\alpha} F\left(\frac{\omega}{\alpha}\right)
\frac{\pi}{2} e^{ax}
x^{-1/2}
\end{verbatim}

Таким образом, мы получили отклики \texttt{U\_n}(t) на простейшие
входные воздействия \texttt{F\_n}(t) . Для получения решения исходной
задачи осталось совершить последний шаг---сложить все отклики

\begin{enumerate}
\def\labelenumi{\arabic{enumi}.}
\setcounter{enumi}{4}
\tightlist
\item
  \texttt{H(a-x)}
\end{enumerate}

\begin{verbatim}
C_n = \frac{2}{\pi} \int_0^{\pi} f(x) \cos(nx) dx.
\end{verbatim}

\begin{enumerate}
\def\labelenumi{\arabic{enumi}.}
\setcounter{enumi}{5}
\tightlist
\item
  \texttt{x\^{}\{-1\}}
\end{enumerate}

\begin{verbatim}
Читатель, должно быть, заметил, что каждый отклик  U_n(t)  умножается на весовой множитель  \sin(n\pi x)  . Конечно, в разложении функции f(x, t) на компоненты  F_n(t)  присутствуют те же весовые множители.
\end{verbatim}

\section{ЗАМЕЧАНИЕ}\label{ux437ux430ux43cux435ux447ux430ux43dux438ux435}

Використання скінченного синусоїдального перетворення дозволяє записати
розв'язок у вигляді ряду, але для більшості інших інтегральних
перетворень розв'язки представлені як інтеграли (неперервний розклад).
Читача можна попросити дати інтерпретацію проблеми

\begin{enumerate}
\def\labelenumi{\arabic{enumi}.}
\setcounter{enumi}{6}
\tightlist
\item
  \texttt{\textbackslash{}frac\{x\}\{x\^{}2+a\^{}2\}}
\end{enumerate}

\begin{verbatim}
y = \frac{\pi}{b-a}(x-a).
\end{verbatim}

\begin{enumerate}
\def\labelenumi{\arabic{enumi}.}
\setcounter{enumi}{8}
\tightlist
\item
  \texttt{\textbackslash{}arctan\ \textbackslash{}frac\{a\}\{x\}}
\end{enumerate}

\begin{verbatim}
f(x) = \frac{C_0}{2} + \sum_{n=1}^{\infty} C_n \cos(nx).
\end{verbatim}

Откуда получаем два уравнения

\begin{verbatim}
f(x) = \frac{C_0}{2} + \sum_{n=1}^{\infty} C_n \cos(nx) \qquad C_n = \frac{2}{\pi} \int_{0}^{\pi} f(x) \cos nx \, dx
0 < x < \infty \qquad \qquad 0 < \omega < \infty
0 \le x \le \pi \qquad n = 0, 1, 2, \dots
\end{verbatim}

j'\,'(x)

\begin{verbatim}
S_n = \frac{2}{\pi} \int_0^{\pi} f(x) \sin(nx) dx.
y = \frac{\pi}{b - a} (x - a).
\end{verbatim}

f'\,'(x)

\begin{verbatim}
f(x) = \sum_{n=1}^{\infty} S_n \sin(nx).
-n^{2}C_{n} - \frac{2}{\pi} \{f'(0) - (-1)^{n} f'(n)\}
S_n = \frac{2}{\pi} \int_{0}^{\pi} f(x) \sin(nx) dx
\begin{array}{ll} u\left(x, \ 0\right) = \sin{(\pi x)}, \\ u_{t}\left(x, \ 0\right) = 0, \end{array}
\end{verbatim}

Воспользовавшись ортогональностью семейства функції (sin (nnx)) на
отреже {[}0, 1{]}, находим формулы для коефіцієнтів

(21.6)

\begin{enumerate}
\def\labelenumi{\alph{enumi})}
\tightlist
\item
\end{enumerate}

\begin{verbatim}
\frac{\pi}{2} e^{ax}
\end{verbatim}

\begin{verbatim}
5. \cos(mx), m=1, 2, \dots
\end{verbatim}

\begin{verbatim}
C_n = \frac{2}{\pi} \int_0^{\pi} f(x) \cos(nx) dx.
\end{verbatim}

\begin{verbatim}
\mathbf{\Pi}  для яких рівнянь функція  u_1 + u_2  також буде розв'язком? Які висновки ви можете зробити зі своїх відповідей?
\end{verbatim}

\begin{enumerate}
\def\labelenumi{\arabic{enumi}.}
\setcounter{enumi}{3}
\tightlist
\item
  Найдите четыре смещанные задачи, сумма решений которых дает решение
  следующей задачи:
\end{enumerate}

\begin{verbatim}
y = \frac{\pi}{b-a}(x-a).
f(x) = \frac{C_0}{2} + \sum_{n=1}^{\infty} C_n \cos(nx).
0 \le x \le \pi \qquad n = 0, 1, 2, \dots
\end{verbatim}

\begin{enumerate}
\def\labelenumi{\arabic{enumi}.}
\setcounter{enumi}{4}
\tightlist
\item
  Решите задачу Коши
\end{enumerate}

(OAY)

f'\,'(x)

(HY) \texttt{U\_n}(0) = 0. Можете ли вы проверить полученное решение?

\begin{enumerate}
\def\labelenumi{\arabic{enumi}.}
\setcounter{enumi}{5}
\tightlist
\item
  Предположим, что функции \texttt{u\_1} и \texttt{u\_2} удовлетворяют
  линейным однородным граничным условиям
\end{enumerate}

\texttt{2.\ \textquotesingle{}\textbackslash{}frac\{a\_\{0\}\}\{2\}\ +\ \textbackslash{}sum\_\{n=1\}\^{}\{\textbackslash{}infty\}\ a\_\{n\}\ \textbackslash{}cos(nx)\textquotesingle{}\ \textquotesingle{}a\_\{n\}\textquotesingle{}\ 3.\ \textquotesingle{}f(n-x)\textquotesingle{}\ \textquotesingle{}(-1)\^{}\{n\}\ \textbackslash{}frac\{2\}\{\textbackslash{}pi\}\ C\_\{n\}\textquotesingle{}\ 4.\ 1}

\$\$
\texttt{\textbackslash{}begin\{cases\}\ u\_x(1,\ t)\ +\ h\_2\ u(1,\ t)\ =\ 0.\ \ Чи\ відповідає\ функція\ u\_i\ +\ u\_i\ цим\ умовам?\ ,\ .}

Лекция 27

\section{УРАВНЕНИЯ ПЕРВОГО ПОРЯДКА (МЕТОД
ХАРАКТЕРИСТИК)}\label{ux443ux440ux430ux432ux43dux435ux43dux438ux44f-ux43fux435ux440ux432ux43eux433ux43e-ux43fux43eux440ux44fux434ux43aux430-ux43cux435ux442ux43eux434-ux445ux430ux440ux430ux43aux442ux435ux440ux438ux441ux442ux438ux43a-1}

МЕТА ЛЕКЦІЇ: Ввести поняття диференціального рівняння з похідними в
похідних першого порядку (поки що ми розглядали рівняння другого
порядку) і ознайомитися з важливим методом розв'язання задач із
початковими умовами --- методом характеристик. Проблема, яку ми зараз
розглянемо, виглядає так:

(УЧП)

\begin{verbatim}
5.  \cos(mx), m=1, 2, \dots   \begin{cases}
\end{verbatim}

\end{document}
